\documentclass[11pt]{amsart}
\usepackage[utf8]{inputenc}
\usepackage[T1]{fontenc}
\usepackage{amsmath,amssymb,amsthm,mathtools}
\usepackage[margin=1.1in]{geometry}
\usepackage{enumitem}
\usepackage{booktabs}
\usepackage{array}
\usepackage{longtable}
\usepackage{tabularx}
\usepackage{url}
\usepackage[hypertexnames=false]{hyperref}
\hypersetup{hidelinks}
\usepackage[final,expansion=false]{microtype}
\setlength{\parskip}{3pt plus 1pt}
\emergencystretch=2.5em
\setcounter{tocdepth}{2}

% ---------------------------------------------------------------- theorem environments (one counter per section)
\newtheorem{theorem}{Theorem}[section]
\newtheorem{proposition}[theorem]{Proposition}
\newtheorem{lemma}[theorem]{Lemma}
\newtheorem{corollary}[theorem]{Corollary}
\theoremstyle{definition}
\newtheorem{definition}[theorem]{Definition}
\newtheorem{problem}[theorem]{Problem}
\newtheorem{observation}[theorem]{Observation}
\theoremstyle{remark}
\newtheorem{remark}[theorem]{Remark}
\numberwithin{equation}{section}

% ---------------------------------------------------------------- number sets and operators
\newcommand{\C}{\mathbb{C}}
\newcommand{\R}{\mathbb{R}}
\newcommand{\N}{\mathbb{N}}
\newcommand{\Z}{\mathbb{Z}}
\newcommand{\Q}{\mathbb{Q}}
\DeclareMathOperator{\Rel}{Re}
\DeclareMathOperator{\Iml}{Im}
\DeclareMathOperator{\supp}{supp}
\DeclareMathOperator{\Tr}{Tr}
\DeclareMathOperator{\spn}{span}

% ---------------------------------------------------------------- notation of the paper (see conventions.tex)
\newcommand{\IT}{I_T}                                  % (-T,T)
\newcommand{\Fh}{\widehat F}                           % Fourier transform
\newcommand{\cF}{c_F}                                  % autocorrelation
\newcommand{\PhiF}{\Phi_F}                             % Mellin-type transform
\newcommand{\PT}{\mathcal{P}_T}                        % primitive space
\newcommand{\DT}{\mathcal{D}_T}                        % form domain on I_T
\newcommand{\NT}{\mathcal{N}_T}                        % contact set
\newcommand{\mT}{\mathfrak{m}_T}                       % archimedean form
\newcommand{\pT}{\mathfrak{p}_T}                       % prime form
\newcommand{\aT}{\mathfrak{a}_T}                       % truncated Weil form
\newcommand{\WT}{W_T}                                  % complete Weil form
\newcommand{\sT}{\sigma_T}                             % symbol
\newcommand{\lam}{\lambda(T)}                          % bottom on primitive functions
\newcommand{\lamW}{\lambda^{\mathrm W}(T)}             % bottom of the complete form
\newcommand{\Tst}{T_{*}}                               % first contact window
\newcommand{\Tarch}{T^{*}_{\infty,\mathrm P}}          % archimedean threshold

\begin{document}

\title{Routes to Weil positivity on finite intervals}

\author{David Alejandro Trejo Pizzo}

\date{\today}

\begin{abstract}
For $T>0$, Weil's explicit formula defines a quadratic form on functions supported in $(-T,T)$, the difference of
an archimedean form and a finite sum over the prime powers below $e^{2T}$. By Weil's criterion, the Riemann
hypothesis is equivalent to the nonnegativity of this form, for every $T$, on the functions orthogonal to the two
polar directions. This paper is a survey of the attempts to prove that nonnegativity by working with the form
itself. We define the form and prove the elementary facts about it, describe the windows on which its sign has
been established with explicit bounds, and prove several theorems showing that whole classes of arguments cannot
give the sign: for each class there is a weight sequence close to the von Mangoldt function which satisfies the
premises of the class and for which the form is negative. We then go through the approaches that have been
tried, grouped by idea, and say for each what was obtained and why it stops where it does. The sign has since
been proved by a different route, through a one-sided criterion for Dirichlet $L$-functions; the last section
relates that proof to the obstructions described here and lists the questions about the form that remain open.
\end{abstract}

\maketitle
\tableofcontents

\section{Introduction}\label{w1:sec:intro}

Weil's explicit formula \cite{Weil52} expresses, for a test function $F$ of compact support on the real line, a
Hermitian quantity built from the zeros of the Riemann zeta function as the sum of three terms: a polar term
coming from the poles of $\zeta$ at $s=0,1$, an archimedean term coming from the Gamma factor, and a sum over
prime powers. Weil observed that the Riemann hypothesis is equivalent to the nonnegativity of this quantity for
all test functions; Bombieri's account \cite{Bombieri2000} gives the analytic details and the form of the
criterion that we use. When the support of $F$ is confined to an interval $(-T,T)$ the prime sum becomes finite,
since only the prime powers $n$ with $\log n<2T$ contribute, and one is led to the following problem, which we
call the \emph{window problem} (Problem~\ref{w1:prob:window}): to show that for every $T>0$ the truncated Weil
form
\[
\aT=\mT-\pT ,
\]
the difference of an archimedean form $\mT$ and a prime form $\pT$, is nonnegative on the space $\PT$ of
\emph{primitive} functions, those supported in $(-T,T)$ and orthogonal to the two polar directions. An off-line
zero of $\zeta$ produces primitive functions of compact support on which the form is negative
(Theorem~\ref{w1:thm:criterion}), so the window problem for all $T$ is equivalent to the Riemann hypothesis.

For small windows the problem is elementary. If $T\le\frac12\log2$ no prime power is seen and the form is the
archimedean one, which is positive there; Yoshida \cite{Yoshida1992} proved positivity in this range. Connes's
programme \cite{Connes1999,ConnesConsani2021} places the same positivity in the framework of a trace formula, in
which the archimedean contribution is isolated and studied by itself. For larger windows the nature of the
problem changes. We show in Section~\ref{w2:sec:structure} that the archimedean form has at least
$(2t_0/\pi)T-3$ negative directions on $\PT$, where $t_0\approx6.2898$ is the positive zero of the archimedean
symbol (Theorem~\ref{w2:thm:archcount}). The primes must therefore compensate a number of negative directions
that grows linearly with $T$, and they do so with almost nothing to spare: on the windows where it has been
computed, the least eigenvalue $\lam$ of the form on $\PT$ is positive and extremely small
(Table~\ref{w2:tab:certified} and Section~\ref{w2:sec:margin}).

This paper is a survey of the attempts to prove the window problem directly, that is, by an argument about the
form $\aT$ and the arithmetic function $\Lambda$ that enters it. Many such attempts have been made, in the
literature and in earlier papers of this series, and they are scattered: positivity by sums of squares and by
decomposition over the primes, positivity imported from probability or geometry, dual certificates and linear
programming, Schur complements and determinants, positive kernels, realizations in Kre\u{\i}n and de~Branges
spaces, the coefficients of Li and Keiper, prolate operators and semilocal trace formulas, estimates on the prime
side, Euler products and Beurling's generalized integers, transfers of the Hodge index theorem, and deformations
in $T$ or in the weights. We describe them in one place and in one notation, with proofs of what can be proved
in a few lines, and for each approach we try to say precisely why it does not give the sign.

The reasons turn out to be few. The most useful tool for finding them is to replace $\Lambda$ by another weight
sequence $a$ and to consider the form $\aT[a]$ so obtained. If an argument uses only properties of $\Lambda$ that
$a$ shares, and if $\aT[a]$ is negative somewhere on $\PT$, the argument cannot be correct. The weights that
serve are simple: multiples $c\Lambda$, sequences obtained from $\Lambda$ by changing one value, truncations,
and the weights of Beurling systems or of functions such as the Davenport--Heilbronn function which have a
functional equation but zeros off the line. With them we prove the following statements, among others.

\begin{itemize}
\item No identity that holds for every weight in a set containing all multiples $c\Lambda$ can give the sign on
a window $T>\frac12\log2$ (Theorem~\ref{w3:bar:B1}).
\item No argument that uses only the values $\Lambda(n)$ for $n<N$, together with properties that survive an
arbitrary change of one later coefficient, can decide a window $T>\frac12\log q$ with $q\ge N$
(Theorem~\ref{w3:bar:B11}).
\item Every nonzero finite prime part is indefinite on $\PT$, so that no decomposition prime by prime can work
(Theorem~\ref{w3:bar:B6}).
\item The density of the primes does not suffice: with the prime form replaced by its main term the form is
negative on $\PT$ for large $T$ (Proposition~\ref{w8:prop:kappa}).
\end{itemize}

Each approach in Sections~\ref{w4:sec:squares}--\ref{w8:sec:families} meets one of three ends. Either it does
not see the arithmetic of $\Lambda$ finely enough, and one of the weights above is a counterexample to it; or it
loses, through an inequality, a quantity larger than the margin; or it succeeds in transforming the problem and
arrives at an exact reformulation of it, no easier than the original. Several of these reformulations are of
independent interest and are recorded as such.

The sign of the form on every window is now known. It was obtained in \cite{P46}, not by any of the approaches
surveyed here but through a one-sided criterion of Landau type for Dirichlet $L$-functions and an input from
the theory of automorphic forms due to Alpoge \cite{CharacterFamily}. We have nevertheless kept the survey in
the form in which the problem presented itself before that proof, for two reasons. The obstructions are
theorems about the explicit formula and remain true; they apply equally to $L$-functions for which no such
proof is available. And the proof gives the sign through the zeros, without exhibiting any positive structure
in the form itself, so that the questions raised by the direct approaches remain open. We return to this in
Section~\ref{w8:sec:open}.

This paper is organised as follows. Section~\ref{w1:sec:setting} defines the form on a finite interval, its
archimedean, prime and polar parts, the primitive subspace and the least eigenvalue $\lam$, and proves the
explicit formula in our normalisation. Section~\ref{w2:sec:structure} describes what is known about the window
problem: the structure of the archimedean form, the windows on which the sign is established with explicit
bounds, the reductions to finite problems, and the size of the margin. Section~\ref{w3:sec:falsifiers} introduces
the test weights and proves the theorems that exclude classes of arguments.
Sections~\ref{w4:sec:squares}--\ref{w8:sec:families} go through the approaches: squares, local decompositions
and imported positivity (Section~\ref{w4:sec:squares}); duality, Schur complements and positive kernels
(Section~\ref{w5:sec:dual}); operator-theoretic and function-theoretic criteria (Section~\ref{w6:sec:operators});
arithmetic approaches through prime sums, Euler products and Beurling systems (Section~\ref{w7:sec:primeside});
geometric approaches and deformations (Section~\ref{w7:sec:geometry}). Section~\ref{w8:sec:constraints} draws
from the obstructions a list of properties that any direct proof must have, and Section~\ref{w8:sec:open}
concludes. A table at the end lists the approaches with the place where each is discussed.

Earlier papers of the series contain partial surveys: \cite{P01} is a general survey of approaches to the
Riemann hypothesis, \cite{P06} a discussion of positivity formulations and of the analogy with function fields,
and \cite{P34,P36} accounts of the obstructions met by spectral, adelic and positivity approaches. The present
paper is organised around the single family of window problems and supersedes them on that subject. Beyond the
series it rests on the classical literature on the explicit formula \cite{Weil52,Bombieri2000,Yoshida1992}, on
Connes's programme \cite{Connes1999,ConnesConsani2021}, on Li's criterion \cite{Li1997,BombieriLagarias1999},
on the work on the de~Bruijn--Newman constant and Jensen polynomials
\cite{deBruijn1950,Newman1976,RodgersTao2020,GORZ2019}, and on the theory of Beurling generalized integers \cite{DiamondZhang2016,DMV2006,Zhang2007}.

\section{The Weil form on a finite interval}\label{w1:sec:setting}

We fix the notation of the paper, derive the explicit formula in the
normalization in which it is used, and prove the elementary facts about the
form on which the later sections rely. The notation is collected at the end of
the section.

\subsection{The form and its three parts}
Fix $T>0$ and write $\IT=(-T,T)$; we call $T$, or $\IT$, a \emph{window}.
A function $F\in L^2(\IT)$ is always extended by zero to $\R$. We put
\begin{equation}\label{w1:eq:transforms}
  \Fh(\tau)=\int_\R F(t)e^{i\tau t}\,dt,\qquad
  \PhiF(s)=\int_\R F(t)e^{(s-\frac12)t}\,dt,\qquad
  \cF(u)=\int_\R F(t)\overline{F(t-u)}\,dt .
\end{equation}
Since $F$ has compact support, $\Fh$ and $\PhiF$ are entire, and
$\PhiF(\tfrac12+i\tau)=\Fh(\tau)$; more generally $\PhiF(s)=\Fh(-i(s-\frac12))$
when $\Fh$ is continued to $\C$. The autocorrelation $\cF$ is continuous,
satisfies $\cF(-u)=\overline{\cF(u)}$ and $|\cF(u)|\le\|F\|_2^2$, vanishes for
$|u|\ge2T$, and has Fourier transform
$\widehat{\cF}(\tau)=\int\cF(u)e^{i\tau u}\,du=|\Fh(\tau)|^2$ for real $\tau$.
The two \emph{Tate moments} of $F$ are
\begin{equation}\label{w1:eq:tate}
  A(F)=\PhiF(1)=\int F(t)e^{t/2}\,dt,\qquad
  B(F)=\PhiF(0)=\int F(t)e^{-t/2}\,dt .
\end{equation}
Both are continuous linear functionals on $L^2(\IT)$, since
$e^{\pm t/2}\in L^2(\IT)$.

Let $\psi=\Gamma'/\Gamma$ and define the \emph{archimedean symbol}
\begin{equation}\label{w1:eq:symbol-m}
  m(\tau)=\Rel\psi\bigl(\tfrac14+\tfrac{i\tau}2\bigr)-\log\pi ,\qquad
  \tau\in\R .
\end{equation}
It is even, real-analytic, and $m(\tau)=\log(|\tau|/2\pi)+O(\tau^{-2})$ as
$|\tau|\to\infty$ by Stirling's formula; its value at the origin is
$C_0=m(0)=\psi(\tfrac14)-\log\pi=-\gamma-\tfrac\pi2-3\log2-\log\pi\approx-5.3722<0$.
Let $\Lambda$ be the von Mangoldt function and put
\begin{equation}\label{w1:eq:weights}
  v_n=\frac{2\Lambda(n)}{\sqrt n}\quad(n\ge2),\qquad
  \NT=\{n\ge2:\ \log n<2T\}.
\end{equation}
We call $\NT$ the \emph{contact set} of the window: by the support property of
$\cF$, these are the only $n$ for which $\cF(\log n)$ can be nonzero.
The \emph{form domain} is
\[
  \DT=\Bigl\{F\in L^2(\IT):\ \int_\R\log(2+|\tau|)\,|\Fh(\tau)|^2\,d\tau<\infty\Bigr\},
\]
a Hilbert space for the norm
$\|F\|_{\DT}^2=\int\log(2+|\tau|)|\Fh|^2\,d\tau$ (it is a closed subspace of
a weighted $L^2$ space on the Fourier side, since $L^2(\IT)$ is closed in
$L^2(\R)$). On $\DT$ we define the \emph{archimedean form}, the \emph{prime
form} and the \emph{truncated Weil form}
\begin{align}
  \mT(F)&=\frac1{2\pi}\int_\R m(\tau)\,|\Fh(\tau)|^2\,d\tau,\label{w1:eq:arch}\\
  \pT(F)&=\sum_{n\in\NT}v_n\,\Rel\cF(\log n)
         =\frac1{2\pi}\int_\R\Bigl(\sum_{n\in\NT}v_n\cos(\tau\log n)\Bigr)|\Fh(\tau)|^2\,d\tau,\label{w1:eq:prime}\\
  \aT(F)&=\mT(F)-\pT(F)=\frac1{2\pi}\int_\R\sT(\tau)\,|\Fh(\tau)|^2\,d\tau,
  \qquad \sT(\tau)=m(\tau)-\sum_{n\in\NT}v_n\cos(\tau\log n).\label{w1:eq:trunc}
\end{align}
The second expression for $\pT$ follows from
$\cF(u)=\frac1{2\pi}\int|\Fh(\tau)|^2e^{-i\tau u}\,d\tau$ and
$\Rel\cF(u)=\frac1{2\pi}\int|\Fh|^2\cos(\tau u)\,d\tau$. Informally, $\pT$ is
the pairing of $\cF$ with the ``comb'' $\sum_{n}\frac{\Lambda(n)}{\sqrt n}(\delta_{\log n}+\delta_{-\log n})$.
Finally the \emph{complete Weil form} adds the polar term:
\begin{equation}\label{w1:eq:complete}
  \WT(F)=\aT(F)+2\Rel\bigl(A(F)\overline{B(F)}\bigr).
\end{equation}
All four forms are Hermitian forms on $\DT$; we use the same letter for the
quadratic form and its polarization. The archimedean form is bounded below
on $L^2$, $\mT(F)\ge C_0\|F\|_2^2$, since $m\ge m(0)=C_0$
(Proposition~\ref{w2:prop:levy}); the prime form is bounded on $L^2$,
$|\pT(F)|\le(\sum_{n\in\NT}v_n)\|F\|_2^2$; and $\mT(F)+c\|F\|_2^2$
is equivalent to $\|F\|_{\DT}^2$ for $c$ large, because
$m(\tau)-\frac12\log(2+|\tau|)$ is bounded below and $m(\tau)\le\log(2+|\tau|)$.
Hence $\aT$ and $\WT$ are closed, lower semibounded forms on $\DT$.

\subsection{The explicit formula and the window problem}
The zeros of the completed zeta function
$\xi(s)=\frac12s(s-1)\pi^{-s/2}\Gamma(s/2)\zeta(s)$ are denoted $\rho$, counted
with multiplicity; sums over $\rho$ are taken symmetrically
($|\Iml\rho|<R$, $R\to\infty$), and converge absolutely for the test functions
below.

\begin{theorem}[explicit formula on a window]\label{w1:thm:explicit}
For every $F\in C_c^\infty(\IT)$,
\begin{equation}\label{w1:eq:explicit}
  \WT(F)=\sum_\rho\PhiF(\rho)\,\overline{\PhiF(1-\bar\rho)} .
\end{equation}
\end{theorem}

This is Weil's explicit formula \cite{Weil52} (see \cite{Bombieri2000} for
the analytic details); a proof in this normalization is
\cite[Theorem~3.3]{P45}. Because normalizations of the explicit formula vary
in the literature, we show how \eqref{w1:eq:explicit} follows from the
standard form. In the form of \cite[Theorem~3.3]{P45} it reads: if
$g\in C_c^\infty(\R)$ and $h(z)=\int g(u)e^{izu}\,du$, then
\begin{equation}\label{w1:eq:standardEF}
  \sum_\rho h(\gamma_\rho)
  =h(\tfrac i2)+h(-\tfrac i2)-g(0)\log\pi
   +\frac1{2\pi}\int_\R h(r)\,\Rel\psi\bigl(\tfrac14+\tfrac{ir}2\bigr)\,dr
   -\sum_{n\ge2}\frac{\Lambda(n)}{\sqrt n}\bigl(g(\log n)+g(-\log n)\bigr),
\end{equation}
where $\rho=\frac12+i\gamma_\rho$ ($\gamma_\rho\in\C$).

\begin{proof}[Derivation of \eqref{w1:eq:explicit} from \eqref{w1:eq:standardEF}]
Take $g=\cF$, which lies in $C_c^\infty(-2T,2T)$; then $h(r)=|\Fh(r)|^2$ for
real $r$. We compare the four terms.

For the zero side, with $s-\frac12=i\gamma$ we have
\[
\PhiF(s)\overline{\PhiF(1-\bar s)}
=\iint F(t)\overline{F(t')}e^{(s-\frac12)t}e^{-(s-\frac12)t'}\,dt\,dt'
=\int\cF(u)e^{(s-\frac12)u}\,du=h(\gamma),
\]
using
$\overline{e^{(\frac12-\bar s)t'}}=e^{-(s-\frac12)t'}$ and $u=t-t'$. So the
left-hand side of \eqref{w1:eq:standardEF} is the right-hand side of
\eqref{w1:eq:explicit}. For the polar term,
\[
h(-\tfrac i2)=\int\cF(u)e^{u/2}\,du
=\iint F(t)\overline{F(t')}e^{t/2}e^{-t'/2}\,dt\,dt'=A(F)\overline{B(F)},
\]
and
likewise $h(\frac i2)=B(F)\overline{A(F)}$; their sum is
$2\Rel(A(F)\overline{B(F)})$. For the archimedean term,
$g(0)=\cF(0)=\frac1{2\pi}\int|\Fh|^2$, so
$-g(0)\log\pi+\frac1{2\pi}\int h\,\Rel\psi(\frac14+\frac{ir}2)\,dr=\mT(F)$.
For the prime term, $\cF(-u)=\overline{\cF(u)}$ gives
$g(\log n)+g(-\log n)=2\Rel\cF(\log n)$, which vanishes unless $n\in\NT$; so
the prime sum equals $\sum_{n\in\NT}v_n\Rel\cF(\log n)=\pT(F)$.
Adding up gives $\WT(F)$.
\end{proof}

If every zero lies on the critical line then $1-\bar\rho=\rho$, each term of
\eqref{w1:eq:explicit} equals $|\PhiF(\rho)|^2$, and $\WT(F)\ge0$. Weil's
criterion is the converse; we need it in a form which uses only functions
on which the polar term vanishes.

\begin{definition}\label{w1:def:primitive}
A function $F\in\DT$ is \emph{primitive} if $A(F)=B(F)=0$. The primitive
functions form the closed subspace $\PT$ of $\DT$. We write $\PT^{\mathrm{ev}}$,
$\PT^{\mathrm{odd}}$ for the primitive functions with $F(-t)=F(t)$,
resp.\ $F(-t)=-F(t)$. Since $t\mapsto -t$ exchanges $e^{t/2}$ and $e^{-t/2}$,
for even $F$ one has $A(F)=B(F)$ and for odd $F$ one has $A(F)=-B(F)$; thus
on each parity sector primitivity is a single linear condition.
\end{definition}

On $\PT$ the polar term vanishes, so $\WT=\aT$ there. In terms of the
autocorrelation: $\int\cF(u)e^{su}\,du=\PhiF(\frac12+s)\overline{\PhiF(\frac12-\bar s)}$,
so for primitive $F$ this entire function has a zero of order at least two at
$s=\frac12$ (both factors vanish there). For real $F$, where $\cF$ is even,
the vanishing of the function and of its derivative at $s=\frac12$ are the two
conditions $\int_0^{2T}\cF(u)g_j(u)\,du=0$ ($j=0,1$) with $g_0(u)=\cosh(u/2)$
and $g_1(u)=u\sinh(u/2)$; we call $g_0,g_1$ the \emph{polar functions}.

\begin{theorem}[Weil's criterion for primitive functions]\label{w1:thm:criterion}
If $\xi$ has a zero off the line $\Rel s=\frac12$, there are nonzero real
primitive functions $F_+$ (even) and $F_-$ (odd) in $C_c^\infty(\R)$ with
$\sum_\rho\Phi_{F_\pm}(\rho)\Phi_{F_\pm}(1-\rho)<0$. Conversely, if all zeros
lie on the line, $\WT(F)\ge0$ for every $T$ and every $F\in C_c^\infty(\IT)$.
\end{theorem}

This is \cite[Theorem~6.2]{P45}, a form of Weil's criterion \cite{Weil52}
following Yoshida \cite{Yoshida1992} and Connes--Consani
\cite{ConnesConsani2021}. We use it only to explain why the problems below
are as hard as the Riemann hypothesis.

\begin{definition}\label{w1:def:lambda}
For $T>0$ put
\[
  \lam=\inf\{\aT(F):\ F\in\PT,\ \|F\|_2=1\},\qquad
  \lamW=\inf\{\WT(F):\ F\in\DT,\ \|F\|_2=1\}.
\]
\end{definition}

Both numbers are finite, by the lower bounds of the previous subsection:
$\lam\ge C_0-\sum_{n\in\NT}v_n$, and similarly for $\lamW$ since the polar term
is bounded on $L^2(\IT)$. Since the forms are closed, $\lam$ is the
bottom of the spectrum of the self-adjoint operator on the $L^2$-closure of
$\PT$ associated with $\aT|_{\PT}$, and $\lamW$ is the analogous quantity for
$\WT$. The embedding $\DT\subset L^2(\IT)$ is compact (the weight
$\log(2+|\tau|)$ tends to infinity and the supports lie in $\IT$), so the
associated operators have compact resolvent: $\lam$ and $\lamW$ are the least
eigenvalues, the infima are attained, and $\lam>0$ if and only if $\aT(F)>0$
for every nonzero $F\in\PT$.

\begin{remark}\label{w1:rem:density}
Smooth functions suffice: $C_c^\infty(\IT)$ is dense in $\DT$ for
$\|\cdot\|_{\DT}$ (dilate the support slightly inwards and mollify; both
operations are continuous for the logarithmic weight), and correcting by a
fixed pair of smooth functions dual to $A,B$ shows that $C_c^\infty(\IT)\cap\PT$
is dense in $\PT$. Hence $\lam$ and $\lamW$ can be computed on smooth test
functions, where Theorem~\ref{w1:thm:explicit} applies.
\end{remark}

\begin{problem}[the window problem]\label{w1:prob:window}
Prove that $\lam\ge0$ for every $T>0$. The \emph{complete window problem}
asks for $\lamW\ge0$ for every $T>0$.
\end{problem}

By Theorem~\ref{w1:thm:criterion} and Remark~\ref{w1:rem:density}, each of the
two problems is equivalent to the Riemann hypothesis: an off-line zero gives
a primitive smooth test function with negative form, supported in some
window. The interest of the window formulation is that, for each fixed $T$,
the prime side is a finite sum over $\NT$, so that the question is a
family, indexed by $T$, of spectral problems for explicit operators; all the
difficulty is uniformity in $T$.

\begin{proposition}\label{w1:prop:elementary}
\begin{enumerate}[label=(\roman*)]
\item $\lamW\le\lam$.
\item $T\mapsto\lam$ and $T\mapsto\lamW$ are nonincreasing.
\item If $T\le\frac12\log2$, then $\NT=\emptyset$ and $\aT=\mT$.
\item $\lamW\ge0$, hence $\lam\ge0$, for $T\le\frac12\log2$.
\end{enumerate}
\end{proposition}

\begin{proof}
(i) On $\PT\subset\DT$ the polar term vanishes and $\WT=\aT$, so the infimum
defining $\lamW$ is taken over a larger set.
(ii) Let $T<T'$. A function in $L^2(\IT)$, extended by zero, lies in
$L^2(I_{T'})$, with the same transforms, moments and $\DT$-norm; so
$\PT\subset\mathcal P_{T'}$ and $\DT\subset\mathcal D_{T'}$. For such $F$,
$\cF(\log n)=0$ whenever $\log n\ge2T$, so the extra terms of
$\mathfrak p_{T'}$ vanish and $\mathfrak a_{T'}(F)=\aT(F)$,
$W_{T'}(F)=\WT(F)$. The infima over the larger sets are smaller.
(iii) $\log n<2T\le\log 2$ has no solution $n\ge2$.
(iv) Yoshida \cite[Theorem~1]{Yoshida1992} proved
$\sum_\rho\PhiF(\rho)\overline{\PhiF(1-\bar\rho)}>0$ for every nonzero smooth
$F$ supported in $[-\frac12\log2,\frac12\log2]$, by a computation that is in
part numerical; see the discussion in \cite[Remark~6.3]{P45}. By
Theorem~\ref{w1:thm:explicit} and Remark~\ref{w1:rem:density} this gives
$\lamW\ge0$ for $T\le\frac12\log2$, and (i) gives $\lam\ge0$.
\end{proof}

For $T\le\frac12\log 2$, then, the window problem concerns the archimedean
form alone; the arithmetic enters at $T=\frac12\log2$, when the first contact
$n=2$ appears, and each prime power $n$ enters at $T=\frac12\log n$. The
structure of $\mT$ on $\PT$ and the windows on which the sign is established
are the subject of Section~\ref{w2:sec:structure}.

\subsection{The polar block and weight sequences}\label{w1:subsec:weights}
The two problems are related by a two-dimensional Schur complement. Let
$V=\spn\{e^{t/2}\mathbf 1_{\IT},e^{-t/2}\mathbf 1_{\IT}\}$; it is the
$L^2(\IT)$-orthogonal complement of the $L^2$-closure of $\PT$, and $V\subset\DT$
(the Fourier transform of $e^{\pm t/2}\mathbf 1_{\IT}$ is $O(1/|\tau|)$, and
$\int\log(2+|\tau|)\tau^{-2}d\tau<\infty$). So $\DT=\PT\oplus V$, an
orthogonal sum in $L^2$.

\begin{proposition}[polar block]\label{w1:prop:polar}
Assume $\lam>0$. For $v\in V$ the function $p\mapsto\WT(v+p)$ on $\PT$ has a
unique minimizer $p(v)$, the map $v\mapsto p(v)$ is linear, and
$\mathfrak G_T(v)=\WT(v+p(v))$ is a Hermitian form on the two-dimensional
space $V$ (the \emph{polar capacity}). Moreover
\[
  \lamW>0\iff\bigl(\lam>0\ \text{and}\ \mathfrak G_T>0\ \text{on}\ V\bigr).
\]
\end{proposition}

\begin{proof}
For $p\in\PT$, $A(v+p)=A(v)$ and $B(v+p)=B(v)$, so
$\WT(v+p)=\aT(v+p)+2\Rel(A(v)\overline{B(v)})$. Since $\lam>0$ we have
$\aT(p)\ge\lam\|p\|_2^2$ on $\PT$; combined with the equivalence
$\aT(p)+c\|p\|_2^2\asymp\|p\|_{\DT}^2$ this gives
$\aT(p)\ge\kappa\|p\|_{\DT}^2$ for some $\kappa>0$, that is, $\aT$ is coercive on
the closed subspace $\PT$ of $\DT$. By the Lax--Milgram theorem there is a
unique $p(v)\in\PT$ with $\aT(v+p(v),q)=0$ for all $q\in\PT$, depending
linearly and boundedly on $v$, and completing the square gives
\begin{equation}\label{w1:eq:schur}
  \WT(v+p)=\mathfrak G_T(v)+\aT\bigl(p-p(v)\bigr)\qquad(v\in V,\ p\in\PT).
\end{equation}
If $\mathfrak G_T\ge g\|v\|_2^2$ with $g>0$ (on the finite-dimensional $V$
this is the same as $\mathfrak G_T>0$), then by \eqref{w1:eq:schur}
$\WT(v+p)\ge g\|v\|^2+\lam\|p-p(v)\|^2$, while
$\|v+p\|\le(1+C)\|v\|+\|p-p(v)\|$ with $C$ the norm of $v\mapsto p(v)$; hence
$\lamW>0$. Conversely, if $\lamW>0$, then $\lam\ge\lamW>0$ by
Proposition~\ref{w1:prop:elementary}(i), and
$\mathfrak G_T(v)=\WT(v+p(v))\ge\lamW\|v+p(v)\|^2\ge\lamW\|v\|^2$ because
$v\perp p(v)$ in $L^2$.
\end{proof}

Thus a proof of the window problem with a positive margin reduces the
complete problem, window by window, to the positivity of an explicit
$2\times2$ Hermitian matrix; a version by parity sectors, which also treats
the case $\lam=0$, is \cite[Theorem~8.6]{P46}, and for the von Mangoldt
weights $\mathfrak G_T$ is positive definite for $T\le\frac12\log5$
\cite[Theorem~9.1]{P46}. We do not know whether positivity of $\lam$
implies positivity of $\mathfrak G_T$ for the general weight sequences
introduced below; computations with model weights suggest that it does not,
in which case positivity of $\mathfrak G_T$ is an additional condition that
no argument using only the shape of the weights can derive from $\lam\ge0$.

Many of the obstructions described in this paper compare the von Mangoldt
weights with other weight sequences. For a sequence $a=(a_n)_{n\ge2}$ of real
numbers we write $\aT[a]$ for the form \eqref{w1:eq:trunc} with $v_n$ replaced
by $2a_n/\sqrt n$, $W_T[a]=\aT[a]+2\Rel(A\overline B)$ for the corresponding
complete form, and $\lambda(T;a)$ for the bottom of $\aT[a]$ on unit vectors
of $\PT$; so $\aT=\aT[\Lambda]$ and $\lam=\lambda(T;\Lambda)$. Typical
comparison sequences are the scaled weights $c\Lambda$, the truncations
$\Lambda\mathbf 1_{n\ge P}$, the twists $\Lambda\chi$ by characters, and weights
attached to Beurling systems or to models resembling $L$-functions. A
\emph{falsifier} for a class of arguments is such a sequence (or model) which
satisfies the premises of the class and has $\lambda(T;a)<0$ for some $T$.

\begin{proposition}[convexity of the admissible weights]\label{w1:prop:convex}
Fix $T>0$. The set
$\mathcal K_T=\{(a_n)_{n\in\NT}\in\R^{\NT}:\ \aT[a]\ge0\ \text{on}\ \PT\}$
is a closed convex subset of $\R^{\NT}$; if nonempty, it is a cone only if it
contains $0$ (a closed cone contains its apex), that is, only if $\mT\ge0$ on $\PT$.
\end{proposition}

\begin{proof}
For fixed $F\in\PT$, $a\mapsto\aT[a](F)=\mT(F)-\sum_{n\in\NT}\frac{2a_n}{\sqrt n}\Rel\cF(\log n)$
is an affine function on the finite-dimensional space $\R^{\NT}$, so
$\{a:\aT[a](F)\ge0\}$ is a closed half-space (or all of $\R^{\NT}$).
$\mathcal K_T$ is the intersection of these half-spaces over $F\in\PT$.
The last assertion is the case $a=0$.
\end{proof}

Proposition~\ref{w1:prop:convex} is behind the most common obstruction met
in this paper: a falsifier $a$ with $\lambda(T;a)<0$ lies outside the convex
set $\mathcal K_T$, and an argument that uses only properties of the weights
shared by $\Lambda$ and $a$ cannot prove $\lam\ge0$ at that $T$.

Three other descriptions of $\WT$ are used in later sections. First, writing
$E$ for evaluation at the zeros and $J$ for the involution induced by
$\rho\mapsto1-\bar\rho$, the right side of \eqref{w1:eq:explicit} is
$\langle JEF,EF\rangle$, a Kre\u{\i}n-space inner product with $J=I$ under the
Riemann hypothesis; see \cite[Proposition~6.1]{P08} for this identity on
Paley--Wiener spaces, and Section~\ref{w6:sec:operators} for the approaches
based on it. Second, for real $F$ the self-intersection of the radial class of
$F$ for the canonical trace on semilocal idele class groups equals $-\aT(F)$
\cite[Theorem~4.1(ii)]{P45}, hence $-\WT(F)$ on primitive functions; the
construction uses the correspondence algebra of \cite{P43} and the Frobenius
family of \cite{P44}. Third, in Connes's framework \cite{Connes1999} the
archimedean form $\mT$ is the contribution of the archimedean place, and its
positivity properties on functions of small support are studied in
\cite{ConnesConsani2021}.

\subsection{Notation and conventions}\label{w1:subsec:notation}
For reference we collect the notation fixed above, together with a few symbols
introduced in Section~\ref{w2:sec:structure}.

\smallskip
\noindent
\begin{tabularx}{\linewidth}{@{}lX@{}}
\toprule
$\IT=(-T,T)$ & window; $F\in L^2(\IT)$ extended by zero\\
$\Fh,\ \PhiF,\ \cF$ & Fourier transform, Mellin-type transform, autocorrelation \eqref{w1:eq:transforms}\\
$A(F),\ B(F)$ & Tate moments $\PhiF(1)$, $\PhiF(0)$ \eqref{w1:eq:tate}\\
$m,\ C_0$ & archimedean symbol \eqref{w1:eq:symbol-m}; $C_0=m(0)\approx-5.3722$\\
$v_n,\ \NT$ & weights $2\Lambda(n)/\sqrt n$; contact set $\{n\ge2:\log n<2T\}$\\
$\sT$ & symbol $m-\sum_{n\in\NT}v_n\cos(\tau\log n)$\\
$\mT,\ \pT,\ \aT,\ \WT$ & archimedean, prime, truncated and complete Weil forms\\
$\DT,\ \PT$ & form domain; primitive subspace $\{A=B=0\}$ (Definition~\ref{w1:def:primitive})\\
$\PT^{\mathrm{ev}},\ \PT^{\mathrm{odd}}$ & even and odd primitive functions\\
$g_0,\ g_1$ & polar functions $\cosh(u/2)$, $u\sinh(u/2)$ on $[0,2T]$\\
$\lam,\ \lamW$ & bottoms of $\aT$ on $\PT$ and of $\WT$ on $\DT$ (Definition~\ref{w1:def:lambda})\\
$\mathfrak G_T$ & polar capacity (Proposition~\ref{w1:prop:polar})\\
$\aT[a],\ \lambda(T;a)$ & forms and bottoms with weights $a$ in place of $\Lambda$\\
$n_-(\cdot)$ & number of negative eigenvalues, with multiplicity\\
$\Tst$ & first window with $\lam=0$, if there is one\\
$\Tarch$ & archimedean threshold, where the bottom of $\mT$ on $\PT$ changes sign (Remark~\ref{w2:obs:arch})\\
$e_k(t)$ & Dirichlet sines $T^{-1/2}\sin(\pi k(t+T)/2T)$, $k\ge1$, an orthonormal basis of $L^2(\IT)$\\
\bottomrule
\end{tabularx}

\smallskip
The words \emph{window}, \emph{primitive}, \emph{contact} and
\emph{falsifier} are used in the senses defined above; ``comb'' is used
informally for the prime form.

\section{The window problem: what is known}\label{w2:sec:structure}

Before describing the approaches to the window problem we collect what is known about the form itself: the structure of
the archimedean symbol and the number of negative directions that the prime form has to compensate, the windows on which
a positive lower bound for $\lam$ or $\lamW$ has been established, the exact reductions of the problem at a fixed window
to a finite matrix, and the size of the margin. That $\lam\ge0$ for every $T$ is obtained in
\cite[Corollary~16.3]{P46} through the zeros of $\zeta$ (Section~\ref{w8:sec:open}); the statements of the present
section concern the form directly, and they are the constraints that every approach of the later sections meets.

\subsection{Structure of the problem}

We begin with the archimedean form. Its symbol is the characteristic exponent of a L\'evy process, shifted by the
negative constant $C_0$.

\begin{proposition}[L\'evy form of the archimedean symbol]\label{w2:prop:levy}
Let $C_0=m(0)=\psi(\frac14)-\log\pi=-5.37218\ldots$ and
\[
K(u)=\frac{1}{2\sinh(u/2)}+\frac{1}{2\cosh(u/2)}=\frac{e^{u/2}}{\sinh u}\qquad(u>0).
\]
Then
\begin{equation}\label{w2:eq:levy}
m(\tau)=C_0+\int_0^\infty\bigl(1-\cos\tau u\bigr)K(u)\,du\qquad(\tau\in\R).
\end{equation}
Moreover $m(\tau)-C_0=\varphi(\tau^2)$ with
$\varphi(s)=\sum_{k\ge0}\frac{s}{(k+\frac14)\,(4(k+\frac14)^2+s)}$, a function that is increasing and concave on
$[0,\infty)$. In particular $m$ is even, strictly increasing in $|\tau|$, and has a unique positive zero $t_0$.
\end{proposition}

\begin{proof}
From $\psi(z)=-\gamma+\int_0^\infty\frac{e^{-x}-e^{-zx}}{1-e^{-x}}\,dx$ ($\Rel z>0$) one gets
$\Rel\psi(\frac14+\frac{i\tau}2)-\psi(\frac14)=\int_0^\infty e^{-x/4}\frac{1-\cos(\tau x/2)}{1-e^{-x}}\,dx$. The
substitution $x=2u$ turns the kernel into $2e^{-u/2}/(1-e^{-2u})=e^{u/2}/\sinh u$, and
$\frac{1}{2\sinh(u/2)}+\frac1{2\cosh(u/2)}=\frac{\cosh(u/2)+\sinh(u/2)}{\sinh u}=\frac{e^{u/2}}{\sinh u}$; this gives
\eqref{w2:eq:levy}. For the second statement use the partial fractions
$\Rel\psi(a+iy)-\psi(a)=\sum_{k\ge0}\bigl(\frac1{k+a}-\frac{k+a}{(k+a)^2+y^2}\bigr)
=\sum_{k\ge0}\frac{y^2}{(k+a)((k+a)^2+y^2)}$ with $a=\frac14$, $y=\tau/2$. Each summand is, as a function of
$s=\tau^2$, of the form $\alpha s/(\beta+s)$ with $\alpha,\beta>0$, hence increasing and concave; so is the sum.
Since $\varphi(s)\to\infty$ (by Stirling's formula, $m(\tau)\sim\log|\tau|$), $m$ has exactly one positive zero.
\end{proof}

Numerically $t_0=6.28983\ldots$ (see also \cite[Proposition~3.1]{P17}), below the first ordinate $14.13\ldots$ of
$\zeta$. Thus $\mT$ is negative on functions whose spectrum is concentrated in $|\tau|<t_0$; by
\cite[Proposition~4.1]{P17} the archimedean term is indefinite on primitive test functions in the classical setting,
without windows, with normalized limit $C_0<0$ for spectra concentrating at the origin. The window problem asks whether
the prime form $\pT$ compensates this negativity exactly on $\PT$. The next theorem counts the directions involved.
Recall the Dirichlet sines $e_k(t)=T^{-1/2}\sin(\omega_k(t+T))$, $\omega_k=\pi k/2T$, an orthonormal basis of
$L^2(\IT)$ with $e_k(\pm T)=0$ and $e_k''=-\omega_k^2e_k$; $e_k$ is even for odd $k$ and odd for even $k$.

\begin{theorem}[negative directions of the archimedean form]\label{w2:thm:archcount}
Let $N\ge1$ with $\omega_N<t_0$ and $B=\varphi'(t_0^2)>0$. For every $F\in\spn\{e_1,\dots,e_N\}$,
\[
\mT(F)\le-B\,(t_0^2-\omega_N^2)\,\|F\|^2<0 .
\]
Consequently $n_-(\mT)\ge N(T):=\lceil 2Tt_0/\pi\rceil-1$, of which $\lceil N(T)/2\rceil$ are even and
$\lfloor N(T)/2\rfloor$ odd, and $n_-(\mT|_{\PT})\ge N(T)-2$. Asymptotically $N(T)=(2t_0/\pi)T+O(1)=4.00423\ldots T+O(1)$.
\end{theorem}

\begin{proof}
By concavity of $\varphi$, $\varphi(s)\le\varphi(t_0^2)+B(s-t_0^2)$ for all $s\ge0$, and $\varphi(t_0^2)=-C_0$; hence
$m(\tau)\le B(\tau^2-t_0^2)$. For $F=\sum_{k\le N}x_ke_k$ extended by zero, $F\in H^1(\R)$ and
$\frac1{2\pi}\int\tau^2|\Fh|^2=\|F'\|^2=\sum\omega_k^2|x_k|^2\le\omega_N^2\|F\|^2$, so
$\mT(F)\le B(\omega_N^2-t_0^2)\|F\|^2$. The condition $\omega_N<t_0$ is $N<2Tt_0/\pi$, whose largest solution is
$N(T)$; the parity count is that of $\{e_1,\dots,e_N\}$. Restriction to the subspace $\PT$ of codimension two loses at
most two negative directions, by the min--max principle.
\end{proof}

Thus the prime form must compensate a negative space of $\mT$ whose dimension grows linearly in $T$, half of it even and
half odd, and any proposal that locates the sign in one or two directions is false for large $T$. The count alone forces
a negative direction of $\mT$ on $\PT$ as soon as $N(T)\ge3$, that is, for
\[
T>T_0(0):=\frac{3\pi}{2t_0}=0.74921\ldots ,
\]
but the bottom of $\mT$ on $\PT$ becomes negative earlier.

\begin{remark}[the archimedean threshold]\label{w2:obs:arch}
The bottom of $\mT$ on $\PT$ is strictly decreasing in $T$ \cite[Lemma~7.5]{P46}, so it changes sign at a single window
$\Tarch$, and $\frac12\log2\le\Tarch\le T_0(0)$ \cite[Proposition~7.1]{P46}. Numerically $\Tarch\in[0.58,0.60]$, and
computations indicate that $n_-(\mT|_{\PT})$ equals $1$, with an even negative direction, from $\Tarch$ up to a
threshold in $[\log2,0.81]$ at which an odd negative direction appears, and equals $2$ from there up to $T=0.90$. We
expect $n_-(\mT|_{\PT})=(2t_0/\pi)T+O(\log T)$; the lower bound is Theorem~\ref{w2:thm:archcount}. Past $\Tarch$ the
primes must contribute positively on $\PT$.
\end{remark}

We turn to the two linear conditions that define $\PT$. Both $\mT$ and $\pT$ commute with the reflection
$F(t)\mapsto F(-t)$: the first is a Fourier multiplier with even symbol, the second is
$\sum v_n\Rel\langle F,\tau_{\log n}F\rangle$, where $\tau_u$ denotes translation by $u$, and
$\langle F,\tau_uF\rangle=\langle\tau_{-u}F,F\rangle$. Hence $\aT$ decouples into an even and an odd sector, and so does
$\PT$; the conditions $A=B=0$ become $A\pm B=0$, one condition in each sector.

\begin{proposition}[interlacing]\label{w2:prop:interlace}
Let $\lambda_1\le\lambda_2\le\dots$ be the min--max values of $\aT$ on $\DT$ and $\mu_1\le\mu_2\le\dots$ those on
$\PT$. Then $\lambda_k\le\mu_k\le\lambda_{k+2}$, and in each parity sector $\lambda^{\pm}_k\le\mu^{\pm}_k\le\lambda^\pm_{k+1}$.
\end{proposition}

\begin{proof}
This is Cauchy interlacing for the restriction of a semibounded form to a closed subspace of codimension two (one in each
sector), which follows from the min--max principle.
\end{proof}

So primitivity is worth exactly two directions, one even and one odd: it removes the polar term but cannot by itself
remove the archimedean negative directions, of which there are about $4T$.

\begin{proposition}[sign changes of primitive functions]\label{w2:prop:signchange}
Every real $F\in\PT$, $F\ne0$, changes sign at least twice: there is no $t_1$ such that $F\ge0$ a.e.\ on one side of
$t_1$ and $F\le0$ a.e.\ on the other.
\end{proposition}

\begin{proof}
If $F$ does not change sign, $A(F)=\int Fe^{t/2}\ne0$. If it changes sign once, at $t_1$, put
$g(t)=e^{t/2}-e^{t_1}e^{-t/2}$, which vanishes only at $t_1$ and has the sign of $t-t_1$. Then $Fg$ has constant sign
and is not a.e.\ zero, so $\int Fg\ne0$; but $\int Fg=A(F)-e^{t_1}B(F)=0$.
\end{proof}

\begin{proposition}[the primitive functions as a range]\label{w2:prop:range}
A function $F\in L^2(\IT)$ satisfies $A(F)=B(F)=0$ if and only if $F=(\frac14-\partial^2)G$ for some $G\in H^2_0(\IT)$
(that is, $G$ extended by zero lies in $H^2(\R)$). In that case $\Fh(\tau)=(\frac14+\tau^2)\widehat G(\tau)$, $F\in\PT$
if and only if $\int(\frac14+\tau^2)^2\log(2+|\tau|)|\widehat G|^2<\infty$, and
\[
\mT(F)=\frac1{2\pi}\int_\R m(\tau)\bigl(\tfrac14+\tau^2\bigr)^2|\widehat G(\tau)|^2\,d\tau .
\]
\end{proposition}

\begin{proof}
If $G\in H^2_0$, integration by parts without boundary terms gives
$\int(\frac14-\partial^2)G\,e^{\pm t/2}=\int G(\frac14-\partial^2)e^{\pm t/2}=0$. Conversely, given $F$ with $A=B=0$,
let $G(t)=-\int_{-T}^t2\sinh\frac{t-s}2F(s)\,ds$, the solution of $G''-\frac14G=-F$ with $G(-T)=G'(-T)=0$. Then
$G(T)=-\int(e^{(T-s)/2}-e^{-(T-s)/2})F(s)ds$ and $G'(T)=-\frac12\int(e^{(T-s)/2}+e^{-(T-s)/2})F(s)\,ds$ are combinations of
$A(F)$ and $B(F)$, hence vanish, and $G\in H^2_0$. The Fourier identity follows from $\widehat{G''}=-\tau^2\widehat G$.
\end{proof}

Equivalently, $F$ is primitive if and only if $\Fh(\pm\frac i2)=0$: the primitive condition removes the two poles of the
explicit formula \eqref{w1:eq:explicit}.

It remains to describe the dependence on $T$.

\begin{proposition}[monotonicity]\label{w2:prop:monotone}
If $T<T'$ then $\lambda(T')\le\lambda(T)$ and $\lambda^{\mathrm W}(T')\le\lambda^{\mathrm W}(T)$. Moreover
$\lamW\le\lam$, and $\lim_{T\to\infty}\lam$ exists in $[-\infty,\infty)$.
\end{proposition}

\begin{proof}
The first three assertions are Proposition~\ref{w1:prop:elementary}(i),(ii): extension by zero maps $\PT$ into
$\mathcal P_{T'}$ and $\DT$ into $\mathcal D_{T'}$ isometrically, the additional contacts of the larger window do not
see a function supported in $\IT$, and on $\PT$ the polar term vanishes. A nonincreasing function has a limit.
\end{proof}

Both $\lam$ and $\lamW$ are continuous in $T$ \cite[Lemma~3.10]{P46}, and $\lam>0$ for $T\le\frac12\log5$
(Table~\ref{w2:tab:certified}). Consequently, if $\lam$ vanishes somewhere, there is a first window $\Tst>\frac12\log5$
with $\lambda(\Tst)=0$. The null space of $\aT$ on $\PT$ at $T=\Tst$ is finite-dimensional, by compactness of the
resolvent, and invariant under reflection, so it contains a null vector $F$ of definite parity
\cite[Lemma~10.17]{P46}. The convex hull of the support of $F$ is exactly $[-\Tst,\Tst]$, since it is symmetric and a
smaller one would give $\lambda(T)\le0$ at some $T<\Tst$; and $\aT(F,G)=0$ for all $G\in\PT$, which is an
Euler--Lagrange equation whose right-hand side is a combination of $e^{\pm t/2}$. We expect such an $F$ to lie in
$H^s$ for $s<\frac12$ and not in $H^{1/2}$ in general. The bottoms of the archimedean forms are strictly decreasing in
$T$ \cite[Lemma~7.5]{P46}; it is plausible that $\lamW$ is too, in which case a nonnegative $\lamW$ at every window is
strictly positive at every window. In the even sector, if $\WT\ge0$ on even functions then the form $\aT$ on the even
functions of $\DT$, with neither the primitive condition nor the polar term, has at most one negative eigenvalue,
because the two forms differ there by the rank-one term $2|A(F)|^2$; we expect the converse to hold when both statements
are required for every $T$. Numerically this form has a negative direction, $\cosh(t/2)\mathbf 1_{\IT}$, for
$T\ge0.12$.

To summarize, in each parity sector the window problem is a cancellation problem counted in directions: a negative
archimedean space of dimension at least about $2T$, which by Theorem~\ref{w2:thm:archcount} depends only on the
archimedean symbol, must be lifted by the prime form on a subspace of codimension one. On the windows of
Table~\ref{w2:tab:certified} the resulting bottom is positive but very small (\S\ref{w2:sec:margin}), and this is a
property of the weight $\Lambda$: the falsifiers of
Section~\ref{w3:sec:falsifiers} show that small changes of the weight destroy the cancellation.

\subsection{Windows on which the sign is established}\label{w2:sec:certified}

For a fixed window, positivity of $\lam$ is equivalent to the positive definiteness of a finite matrix
(\S\ref{w2:sec:reductions}), and lower bounds have been obtained on bounded windows by computations with rigorous error
control. Table~\ref{w2:tab:certified} lists them. A bound for $\lamW$ implies the same bound for $\lam$, since
$\lamW\le\lam$, and by Proposition~\ref{w2:prop:monotone} a bound at $T$ holds on $(0,T]$.

\begin{table}[htbp]
\centering\small
\renewcommand{\arraystretch}{1.15}
\begin{tabular}{p{3.0cm}p{5.2cm}p{5.2cm}}
\toprule
window & bound & where it is proved\\
\midrule
$0<T\le\frac12\log2$ & $\lamW\ge0$ & \cite[Theorem~1]{Yoshida1992}; see also \cite{Bombieri2000}\\
$0<T\le\frac12\log3$ & $\lam\ge1.1599\cdot10^{-4}$ & proof of \cite[Theorem~9.18]{P42}\\
$0<T\le\log2$ & $\lam>0$ & \cite[Theorem~9.18, Corollary~9.20]{P42}\\
$\frac12\log2<T\le0.8$ & $\lamW\ge8.9\cdot10^{-18}$ & \cite{Zhu2026}\\
$0<T\le\frac12\log5$ & $\lam\ge2.5\cdot10^{-12}$, $\lamW\ge10^{-18}$ & \cite[Theorem~9.1]{P46}\\
$0<T\le17/16$ & $\lamW\ge2^{-49162}$ & reported by Liu in an unrefereed manuscript \cite{Liu2026}\\
$0<T\le\frac12\log15$ & $\lamW\ge2^{-2904}$ & \cite{Mori15}\\
$0<T\le2.8022$ & $\lamW\ge e^{-120110026807}$ & \cite[\S9.3]{P46}, using \cite{PlattTrudgian2021}\\
\bottomrule
\end{tabular}
\caption{Windows on which a lower bound for $\lam$ (the bottom of $\aT$ on $\PT$) or for $\lamW$ (the bottom of the
complete form) is established.}
\label{w2:tab:certified}
\end{table}

The entries are of different kinds and are not comparable as theorems. The first is classical: for
$T\le\frac12\log2$ there is no contact, and only the archimedean and polar parts are present. All the others are
computer-assisted. The bounds of \cite{P42} are obtained in interval arithmetic for the primitive Weil operator on
$\PT$; the proof of \cite[Theorem~9.18]{P42} gives the margins $6.547\cdot10^{-4}$ in the even sector and
$1.1599\cdot10^{-4}$ in the odd sector at $T=\frac12\log3$. The bounds of \cite[Theorem~9.1]{P46} are obtained in ball
arithmetic. The preprint \cite{Zhu2026} states $8.9\cdot10^{-18}\le\lambda^{\mathrm W}(0.8)\le2.27\cdot10^{-17}$; it and
\cite{Mori15} are unrefereed, and we have not verified them, beyond reproducing the computational part of
\cite{Mori15} in $1280$-bit arithmetic. Further interval computations, which we do not reproduce here, indicate $\lambda(\frac45)\ge2.96437\cdot10^{-17}$ (with a bound $6.34\cdot10^{-15}$ in the odd sector),
and $\lam\ge2\cdot10^{-26}$ for $0<T\le1.035$. Positivity of the complete form in both parities for
$0<T\le\frac12\log11$ is an external computer-assisted result \cite{RStation45}; we have reproduced it, regenerating
the whole chain of the published supplement and repeating in interval arithmetic the elimination of both kernels. The
pivots of that elimination are positive, but they give no bound on $\lamW$. For very short windows there is a bound in closed form: for $0<T\le0.1$, where $\aT=\mT$, one has
$\lam\ge\log(1/T)-1.31$ \cite[Theorem~10.9]{P46}.

The last row of the table differs from the others in that it uses the zeros. A mask in frequency gives an archimedean
margin $\Delta_T>0$ independent of the zeros, the zeros up to the height $3\cdot10^{12}$ to which the Riemann
hypothesis has been verified \cite{PlattTrudgian2021} absorb the remaining spectral contribution, and the comb is
controlled by Schur bounds verified in ball arithmetic. The method works one window at a time. Heuristically, a floor
obtained in this way requires the zeros to be verified up to a height growing doubly exponentially in $T$ (the
logarithm of the height being at least of order $e^{T}$), and computations indicate that the comparison of Schur
weights used at present stops near $T\approx3.3$--$3.4$. This is a limitation of the method and says nothing about the
sign.

Three kinds of results should be distinguished from the bounds of the table. Positivity of the complete form has been
found, by computations that we do not reproduce here, on several explicit infinite-dimensional but special subspaces: functions whose
Mellin transforms are rational with prescribed poles (up to $5\cdot10^8$ poles in the largest family treated),
combinations of up to $10^4$ powers with exponents larger than $1$, and families with a bounded number of sign changes
or crossings. These spaces are not windows; we return to them in Section~\ref{w5:sec:kernels}. Second, numerical
verifications of one-sided arithmetic criteria equivalent to the Riemann hypothesis, for $x$ in bounded ranges (for
instance up to $10^{20}$, or in $[10^{12},10^{38}]$), concern other formulations and no window; see
Section~\ref{w7:sec:primeside}. Third, a lower bound on a window is a statement about that window, and establishing
such bounds one window at a time is not a method for the window problem: no entry of the table extends to all $T$, and
the precision required grows with $T$ at the rate at which the margin decays, which we discuss in
\S\ref{w2:sec:margin}.

\subsection{Finite reductions}\label{w2:sec:reductions}

Several approaches begin by rewriting the window problem as the positivity of an explicit matrix, or of a sequence of
finite matrices. Such a rewriting is an exact restatement, which neither gains nor loses information. We describe the
reductions and then the attempts to decide the resulting finite problems. We start with two elementary facts.

\begin{proposition}[approximate positivity suffices]\label{w2:prop:soft}
Suppose that $\lam\ge-\varepsilon(T)$ for every $T>0$, with $\varepsilon(T)\to0$ as $T\to\infty$. Then $\lam\ge0$ for
every $T>0$.
\end{proposition}

\begin{proof}
If $\lambda(T_0)<0$ for some $T_0$, then $\lam\le\lambda(T_0)<0$ for every $T\ge T_0$ by
Proposition~\ref{w2:prop:monotone}, contradicting $\lam\ge-\varepsilon(T)\to0$.
\end{proof}

The same argument gives the version along a sequence: $\lambda(T_k)\ge-\varepsilon_k$ with $T_k\to\infty$ and
$\varepsilon_k\to0$ suffices. In particular it is not true that a proof must decide, at each window, a sign of the size
of $\lam$: a proof of approximate positivity with an error tending to zero would be enough.

\begin{proposition}[the prime form on Dirichlet sines]\label{w2:prop:combsine}
Put $D_T(\omega)=\sum_{n\in\NT}v_n\sin(\omega\log n)$ and
$C_T(\omega)=\sum_{n\in\NT}v_n\bigl(1-\frac{\log n}{2T}\bigr)\cos(\omega\log n)$. For real $F=\sum_kx_ke_k$,
$\pT(F)=\sum_{k,l}x_kx_l\Pi_{kl}$, where $\Pi_{kl}=0$ if $k\not\equiv l\pmod 2$, and otherwise
\[
\Pi_{kk}=C_T(\omega_k)+\frac{D_T(\omega_k)}{2T\omega_k},\qquad
\Pi_{kl}=\frac1{2T}\Bigl[-\frac{D_T(\omega_k)-D_T(\omega_l)}{\omega_k-\omega_l}+\frac{D_T(\omega_k)+D_T(\omega_l)}{\omega_k+\omega_l}\Bigr]
\quad(k\ne l).
\]
\end{proposition}

\begin{proof}
For $0\le u<2T$, put $s=t+T$ and $L=2T$, so that $\omega_kL=\pi k$. Then
$\langle e_k,\tau_ue_l\rangle=\int e_k(t)e_l(t-u)\,dt=\frac1T\int_u^L\sin(as)\sin(b(s-u))\,ds$ with $a=\omega_k$,
$b=\omega_l$. For $a\ne b$, writing $\sin x\sin y=\frac12[\cos(x-y)-\cos(x+y)]$, integrating, and using
$\sin(\pi m\pm x)=(-1)^m(\pm\sin x)$ with $m=k\mp l$, one finds, when $k\equiv l\pmod2$,
\[
\langle e_k,\tau_ue_l\rangle=\frac1{2T}\Bigl[-\frac{\sin au-\sin bu}{a-b}+\frac{\sin au+\sin bu}{a+b}\Bigr],
\]
which is symmetric in $(a,b)$, while for $k\not\equiv l$ the same computation gives an expression antisymmetric in
$(a,b)$, which drops out of the quadratic form of a real $F$. For $a=b$ the integral equals
$(1-\frac u{2T})\cos au+\frac{\sin au}{2Ta}$. Multiplying by $v_n$ at $u=\log n$ and summing over $\NT$ gives the
stated matrix, since $\pT(F)=\sum_nv_n\cF(\log n)$ for real $F$.
\end{proof}

The prime form is thus, in each parity sector, a divided-difference (Loewner) matrix of the sine sum $D_T$ plus a
matrix of Cauchy type, together with a diagonal; the phases $\omega_k\log n$ enter exactly. Combined with
Theorem~\ref{w2:thm:archcount}, positivity of $\aT$ on the primitive functions in the span of $e_1,\dots,e_N$ is the
statement that the matrix $\bigl(\mT(e_k,e_l)\bigr)_{k,l\le N}-\Pi$, compressed to the subspace where $A=B=0$, is
positive semidefinite. This is a Galerkin truncation, hence a necessary condition for the window problem and not an
equivalent one. Numerical computations without error control, with up to $80$ sines at windows $T\le\log2$, confirm
the condition, and show that for $T>\Tarch$ it fails once the prime $2$ is removed from the comb: the bottom is then
$\approx-0.15$ at $T=0.6$ and $\approx-0.44$ at $T=\log2$. This matrix is the starting point of the approaches of
Sections~\ref{w5:sec:dual} and~\ref{w5:sec:schur}.

An exact reduction at a fixed window is obtained in \cite[\S10]{P46}. Divide $\IT$ into $k$ consecutive intervals
$B_1,\dots,B_k$ of length $2T/k$, let $\mathcal V_k\subset\DT$ consist of the functions both of whose Tate moments
vanish on every block, and let $E_k\subset\PT$ be the space of primitive functions which on each block are
combinations of $e^{t/2}$ and $e^{-t/2}$; then $\dim E_k=2k-2$ and $\PT=\mathcal V_k\oplus E_k$, an orthogonal sum in
$L^2$. On $\mathcal V_k$ the archimedean form is coercive with a constant that grows with $k$:
for $k\ge T$ one has $\mT(F)\ge(\log\frac kT-C_B(\frac Tk))\|F\|^2$, with an explicit function $C_B$ which tends to
$1.6399$ as $T/k\to0$ \cite[Theorem~10.11]{P46}. Since
the prime form is bounded, the bottom $c=c(T,k)$ of $\aT$ on the unit sphere of $\mathcal V_k$ is positive for
$k\ge k_0(T)$ \cite[Corollary~10.13]{P46}. For such $k$ there is a linear map $\mathsf K\colon E_k\to\mathcal V_k$ with
$\aT(\mathsf Ke,v)=\aT(e,v)$ for all $v\in\mathcal V_k$, and the Schur complement
$\mathfrak s_{T,k}(e)=\aT(e)-\aT(\mathsf Ke)$ is a quadratic form on $E_k$; let $S_{T,k}$ be its matrix in an
orthonormal basis, a real symmetric matrix of size $2k-2$, and $\lambda_S$ its least eigenvalue. Then
\cite[Theorem~10.14]{P46}
\begin{equation}\label{w2:eq:schurred}
\lam>0\iff S_{T,k}\succ0,\qquad \lam\ge0\iff S_{T,k}\succeq0,
\end{equation}
and quantitatively
\[
\frac{\min\{c,\lambda_S\}}{(1+\|\mathsf K\|)^2}\le\lam\le\min\{c,\lambda_S\}\quad(\lambda_S\ge0),\qquad
\lambda_S\le\lam\le\frac{\lambda_S}{1+\|\mathsf K\|^2}\quad(\lambda_S<0).
\]
So at each fixed window the sign of $\lam$ is that of one finite matrix, and all the difficulty of the window problem
is uniformity in $T$; the lower bounds of Table~\ref{w2:tab:certified} are instances of this. The entries of $S_{T,k}$
involve the inverse of the form on $\mathcal V_k$ and are not given in closed form: a Galerkin approximation bounds
$S_{T,k}$ from above, and a lower bound needs an error estimate which uses $c>0$. The value of $k_0(T)$ given by the
proof grows doubly exponentially with $T$, while $k=25$ suffices at $T=\frac12\log5$. The statement that
$S_{T,k}\succeq0$ for all $T>\frac12\log5$ and suitable $k=k(T)$ is the window problem beyond the range of the table,
and we know no proof of it in these terms.

Two further reformulations of the same kind are expected to be equivalent to the window problem; we state them without
proof. In the first, one fixes a profile $F_0$ supported in a window $T_0$ with $2T_0<\log2$ and considers its
translates sampled at Dirichlet frequencies. Positivity for all $T$ should then be equivalent to the positivity, for
every $M$, of a matrix
\begin{equation}\label{w2:eq:nearform}
\mathrm{Arch}_M-\|F_0\|^2\mathrm{Div}_M-\mathrm{Near}_M,
\end{equation}
where $\mathrm{Div}_M$ collects the exact divisibility relations $m=kn$ among the sampled indices (it does not depend
on $\Lambda$ and has nonnegative entries) and $\mathrm{Near}_M$ collects the near coincidences $m\approx kn$, $m\ne kn$.
A divisible enlargement of the sampling set, by $q=\operatorname{lcm}(1,\dots,N)$ or by the product of the primes up
to $N/2$, rewrites the condition as $B_{\mathrm{arch},M}\ge C_M+C_M^*$ for all $M$, with explicit $B$ and $C$. The
second reformulation uses the unitary change of variables $f(x)=x^{-1/2}F(\log x)$ from $L^2(\R,dt)$ to
$L^2((0,\infty),dx)$. Translation by $\log n$ becomes the dilation $D_nf(x)=\sqrt nf(nx)$, one has
$\langle f,D_nf\rangle=\cF(\log n)$ for real $F$, and the conditions $A(F)=B(F)=0$ become $\int f\,dx=\int f\,dx/x=0$.
Since $\aT$ is invariant under translation, which preserves primitivity, one may consider step functions
$f=\sum_{1\le j\le N}a_j\mathbf 1_{[j,j+1)}$ on unit cells. On them the constraints read
\begin{equation}\label{w2:eq:cell}
\sum_j a_j=0,\qquad \sum_j a_j\log\frac{j+1}{j}=0,
\end{equation}
and the form is a matrix $M_N=A_N-C_N$, with an explicit archimedean part $A_N$ (built from second differences of
$x\log x$, a diagonal, and a Stieltjes moment matrix entering with a minus sign) and a prime part $C_N$ expressed
through the integer-division operators $(B_na)_j=a_{\lfloor j/n\rfloor}$, with $a_0=0$. The expected statement is that
$\aT\ge0$ on $\PT$ for all $T$ if and only if $M_N\succeq0$ on the subspace \eqref{w2:eq:cell} for all $N$, and that it
suffices to test integer vectors satisfying \eqref{w2:eq:cell}, of all sizes.

Many attempts were made to prove the positivity of $S_{T,k}$, of $M_N$ or of the matrices \eqref{w2:eq:nearform} by a
general principle: norm bounds, block Schur complements of fixed size, stability in $k$, localization, or transfer
from smaller windows. None succeeded, for one of two reasons. Since the reductions are exact, the finite problem is as
hard as the original one, and a general principle applied to it either loses a constant factor, which is fatal against
a margin that decays like $e^{-cT}$ or faster (\S\ref{w2:sec:margin}), or applies to every weight sequence in a
neighbourhood of $\Lambda$, which the falsifiers of Section~\ref{w3:sec:falsifiers} exclude. We go through the
attempts in turn.

For $S_{T,k}$, the coupling between adjacent blocks has bilinear norm exactly $\pi$, by Hilbert's inequality and its
sharpness, so that subtracting the norms of the coupling and of the prime form separately costs a constant and already
at $T=\frac12\log5$ requires $k\ge70$ in place of $25$ \cite[Remark~10.16]{P46}; an estimate by operator norms
loses a constant, and the margin leaves none. Numerically, the lower constant of uniform $4\times4$ Schur blocks drops to $0.013$ already at $T=0.69$,
and a localization of IMS type gives nothing for $T>0.822$. Refining the partition does not help: when $S_{T,k}$ is
nonsingular, $n_-(S_{T,k})$ is the negative index of $\aT$ on $\PT$ and so does not depend on $k$
\cite[Remark~10.15]{P46}. Computations indicate that the slack of the archimedean block with the primes removed is
negative for $T>\frac12\log5$, and that the hypotheses of a scalar transfer inequality between windows are satisfied by
$2\times2$ models in which the conclusion fails. Determinants of $S_{T,k}$ do not help, since every determinant that
decides the sign has $\lam$ as a factor, and linear certificates in a fixed basis are excluded; see
Section~\ref{w3:sec:barriers}, \ref{w3:bar:B3} and~\ref{w3:bar:B2}.

For the cell matrices $M_N$, partial positivity is found numerically on subfamilies: the Rayleigh minima on the
subspace \eqref{w2:eq:cell} remain positive up to $N=40$; an explicit family $F_k$ ($k\ge3$) supported on three cells
has $\aT(F_k)\ge10^{-3}\|F_k\|^2$, and every plane spanned by $F_3$ and $F_k$ is positive; and on the kernel of the
prime valuations at $N=8$ the form is a positive combination of six squares $|a_i-a_j|^2$. None of this extends. The
same planes are positive for the weight $1.1\Lambda$, while a combination of the $F_k$ at $N=256$ has Rayleigh quotient
below $-0.0931$ for $1.1\Lambda$ and above $0.0367$ for $\Lambda$: local positivity does not control the collective
coupling. Sparse positive representations of $M_N$, of Laplacian type or by squares in two coordinates, appear to be
excluded at all sizes $N\ge6$, resp.\ $N\ge8$, by explicit rational dual certificates; see
Section~\ref{w4:sec:squares}. For the matrices \eqref{w2:eq:nearform}, $\mathrm{Near}_M$ has no sign of its own:
numerically, a packet breaking positivity has entries of size $O(M)$ against a margin of size $\asymp M^{-3}$. The
coefficients of the divisor defect have a definite sign, but the resulting matrix is not self-adjoint, and its
Hermitian part appears to be indefinite on primitive profiles.

The finite problems themselves remain without a direct proof; we return to them in Section~\ref{w8:sec:open}. What
would be needed is a structure that does not inherit the margin: for $S_{T,k}$, a relative bound of the off-diagonal
block by a fixed fraction of the archimedean block; for the cell matrices, an explicit square which uses the
integrality of the dilations (for $z>0$, the identity $\lfloor\lfloor u\rfloor/z\rfloor=\lfloor u/z\rfloor$ holds for
all $u\ge0$ exactly when $z$ is a positive integer) and which explains why it fails for $1.1\Lambda$; for
\eqref{w2:eq:nearform}, joint control of $\mathrm{Near}_M$ and $\mathrm{Arch}_M$ on vectors that see the coincidences
$m\approx kn$. Related finite structures appear in \cite{P26}, on the operator framework of window forms and their
symbols, and in \cite{P39}, on exact transport identities between finite quotients of the semilocal construction.

\subsection{The size of the margin}\label{w2:sec:margin}

The bottom $\lam$ is very small and decreases very fast, and this explains the failure of most approaches that rely on
estimates with slack. We describe what is known about its size, the numerical structure of the low spectrum, and the
sensitivity statements that follow from the form being affine in the weight.

The only general upper bound that is proved is exponential: from an explicit family of test functions one gets
\cite[Corollary~9.11]{P46}
\[
\lamW\le\lam\le4\cdot10^{9}\,T\,e^{-15.449\,T}\qquad(T\ge1).
\]
At particular windows, explicit test functions give numerically
\[
\lambda(\log2)\le9.8326\cdot10^{-8},\quad \lambda(\tfrac12\log3)\le9.4484\cdot10^{-4},\quad
\lambda(\tfrac12\log5)\le3.1029\cdot10^{-12},
\]
and at $T=\frac12\log15$ R.~Mori reports a test function with complete-form value $\approx4.98\cdot10^{-70}$, and the
range $[4.4,8.5]\cdot10^{-70}$ for such values; these are upper bounds for $\lamW$, not values of it. The true decay
is expected to be doubly exponential. Zhu conjectures $-\log\lamW\asymp e^{2T}$ \cite[Conjecture~12.1]{Zhu2026}, and
Mori \cite{MoriUpper} reports a ceiling of the form $C\mu^{9/2}\log\mu\,e^{-4\pi\mu}$, $\mu=e^{2T}$, for $\lamW$, obtained with data from
the spectral side. We expect at least a bound $\lam\le Ce^{cT}e^{-2\pi e^{2T}}$ for large $T$, for the following
heuristic reason: $e^{-2\pi e^{2T}}=(e^{-\pi e^{2T}})^2$ is the square of the Gaussian tail of the term $n=1$ of the
theta function $\sum_n e^{-\pi n^2x^2}$ at the edge $x=e^{T}$ of the window in the multiplicative variable, and an
almost-null test vector built from the kernel of P\'olya's study of the Fourier integral of $\xi$, cut off at the edge
of the window, has a form value of that size. If the decay is of this kind, a lower bound on a window must be
doubly exponentially fine in $T$, and heuristically a certificate with a constant floor needs a basis of dimension at
least of the order of the archimedean count $N(T)\approx4T$ of Theorem~\ref{w2:thm:archcount}.

Galerkin computations of $\lam$, without error control, give
\[
\lambda(0.8)=5.2940(3)\cdot10^{-12},\qquad \lambda(1.0)=1.320(1)\cdot10^{-23},
\]
and, in the even sector, $1.953\cdot10^{-3}$ at $T=0.534$, $3.547\cdot10^{-7}$ at $0.678$, $7.402\cdot10^{-10}$ at $0.75$,
$1.954\cdot10^{-14}$ at $0.85$, $6.57\cdot10^{-19}$ at $0.93$ and $1.80\cdot10^{-27}$ at $1.05$. Beyond $T\approx1.1$
only upper bounds are reliable (for instance $2.79\cdot10^{-37}$ at $T=1.25$). A least-squares fit over $71$
points with $X=e^{2T}\in[2.94,8.17]$,
\[
\log\lam\approx-13.25\,X+13.78\log X+17.7\qquad(\text{rms }0.085),
\]
is far from a leading coefficient $-2\pi$ and of the order of $-4\pi$, without being consistent with it within its
error (another fit gives $-13.07\pm0.23$). Part of the data are not converged, the third digit of the fitted
coefficients is not stable across fits, and neither the exponent nor the prefactor is determined: the size
$e^{-4\pi e^{2T}}$ is suggested by the data, not established.

The bottom is not an isolated small eigenvalue. Numerically, at $T=1$ there are eight directions with Rayleigh
quotient below $4\cdot10^{-6}$, the successive low eigenvalues in each sector are separated by factors $10^2$--$10^3$,
and the number of eigenvalues below any fixed threshold grows with $T$. The bottom vector is even, with two
sign changes, as Proposition~\ref{w2:prop:signchange} requires, and an edge behaviour scaling like $\sqrt{\lam}$; its
interior singularities sit at the points $\pm(T-\log n)$. A counting heuristic gives a cluster of almost-null
directions of size $\asymp e^{2T}$, much larger than the archimedean count $\approx4T$. The bottom also reacts to the
position of individual atoms: moving the atom at $3$ keeps positivity at $T=0.65$ only within $[3.0000,3.001]$, and at
$T=0.85$ only within an interval of length below $10^{-10}$ around $3$. Integers are not distinguished from nearby
non-integers by this test.

The sensitivity to the size of the weights is elementary. Recall from \S\ref{w1:subsec:weights} the form $\aT[a]$ with
weights $2a_nn^{-1/2}$ and its bottom $\lambda(T;a)$ on $\PT$.

\begin{proposition}[concavity and slack]\label{w2:prop:concave}
\emph{(i)} For each $T$, $a\mapsto\lambda(T;a)$ is concave; in particular $\{c\in\R:\lambda(T;c\Lambda)\ge0\}$ is an
interval.
\emph{(ii)} If $F\in\PT$, $\|F\|=1$, has $\pT(F)>0$, then $\lambda(T;(1+\delta)\Lambda)<0$ for every
$\delta>\aT(F)/\pT(F)$.
\end{proposition}

\begin{proof}
For fixed $F$, $\aT[a](F)=\mT(F)-\sum_n2a_nn^{-1/2}\Rel\cF(\log n)$ is affine in $a$, and an infimum of affine functions
is concave; a superlevel set of a concave function of one variable is an interval. For (ii),
$\lambda(T;(1+\delta)\Lambda)\le\aT[(1+\delta)\Lambda](F)=\aT(F)-\delta\,\pT(F)<0$.
\end{proof}

Numerically, on near-minimizers $\pT(F)$ is positive and of size $10^{-1}$ at the windows computed ($\approx0.15$
at $\frac12\log3$ and $\approx0.09$ at $\log2$, Galerkin values without error control), while $\aT(F)$ can be taken
close to $\lam$. Then (ii) says that the weight $(1+\delta)\Lambda$ is negative in the window $T$ for $\delta$ of the
order of the margin. Computations indicate that the interval of (i) contains a neighbourhood of $c=1$ of exactly that
size, and that at $T=\log2$ a relative change of $10^{-5}$ in the weights changes the sign.

\begin{remark}[the precision required of a proof]\label{w2:obs:precision}
An argument that proves $\lambda(T;a)\ge0$ for every weight $a$ in a neighbourhood of $\Lambda$ whose size does not
shrink with the margin cannot succeed: by Proposition~\ref{w2:prop:concave}(ii) such a neighbourhood contains a weight
that is negative in the window. In this sense estimates with slack are refuted by the rescalings $(1+\delta)\Lambda$;
the precise statements are in Section~\ref{w3:sec:barriers}. This does not mean that a proof must be exact to the size
of $\lam$ at each window: by Proposition~\ref{w2:prop:soft} an error $\varepsilon(T)\to0$ is allowed, provided it is not
uniform over a neighbourhood of the weight. The constraints which this imposes on a proof are collected in
Section~\ref{w8:sec:constraints}.
\end{remark}

\section{Test weights and barrier theorems}\label{w3:sec:falsifiers}

The form $\aT$ is affine in the von Mangoldt function, and nothing prevents us from replacing $\Lambda$ by another
weight, or $\zeta$ by another Dirichlet series with a functional equation. Many of the approaches described in
Sections~\ref{w4:sec:squares}--\ref{w8:sec:families} stop for the same reason: every premise they use holds for such
a modified object, and for that object the conclusion is false. This section introduces the modified weights and
series that serve as tests, says what each of them shares with $\zeta$, and proves the general statements that follow
from them. It uses the size of the margin discussed in Section~\ref{w2:sec:margin}.

\subsection{Test weights and what a counterexample weight shows}
\label{w3:subsec:logic}\label{w3:subsec:models}\label{w3:subsec:table}\label{w3:subsec:use}

For a real sequence $a=(a(n))_{n\ge2}$ put $v_n(a)=2a(n)n^{-1/2}$ and define, on the form domain $\DT$,
\begin{equation}\label{w3:eq:aTa}
  \aT[a](F)=\mT(F)-\sum_{n\in\NT}v_n(a)\,\Rel\cF(\log n),
  \qquad
  \lambda(T;a)=\inf\{\aT[a](F):F\in\PT,\ \|F\|_2=1\}.
\end{equation}
Thus $\aT=\aT[\Lambda]$ and $\lam=\lambda(T;\Lambda)$. Only the finitely many values $a(n)$, $n<e^{2T}$, enter
$\aT[a]$. When the test object is a Dirichlet series $L(s)$ rather than a weight, its weight is the sequence of
coefficients of $-L'/L=\sum a_L(n)n^{-s}$, and its archimedean symbol is the one attached to its own gamma
factor and conductor; we then speak of the window form of $L$. In this generality the weight need not be
nonnegative or supported on prime powers.

\begin{definition}\label{w3:def:falsifier}
Let $\mathcal C$ be a class of arguments, described by the set $\mathcal A(\mathcal C)$ of weights (or
Dirichlet series) satisfying every premise used by the arguments of $\mathcal C$. A \emph{falsifier} of
$\mathcal C$ is an element $a'\in\mathcal A(\mathcal C)$ with $\lambda(T;a')<0$ for some $T>0$ (respectively,
for a series, whose window form takes a negative value on primitive functions in some window).
\end{definition}

\begin{proposition}\label{w3:prop:principle}
If an argument of class $\mathcal C$ proves $\lambda(T;a)\ge0$ for every $a\in\mathcal A(\mathcal C)$ and a
fixed $T$, then $\mathcal C$ has no falsifier at that $T$. Consequently, if $\Lambda\in\mathcal A(\mathcal C)$
and $\mathcal C$ has a falsifier at $T$, no argument of $\mathcal C$ proves $\lam\ge0$.
\end{proposition}

\begin{proof}
The conclusion of the argument applies to every element of $\mathcal A(\mathcal C)$, the falsifier included.
\end{proof}

The proposition is a tautology. To apply it one needs two things: a description of $\mathcal A(\mathcal C)$, that is,
of the premises an argument really uses, and a proof that some element of $\mathcal A(\mathcal C)$ has a negative
window. A falsifier says nothing about an argument that uses one more property, which the falsifier lacks. For some
of the test objects below negativity is proved in Section~\ref{w3:sec:barriers} or in the literature; for the others
it rests on computations in interval arithmetic or in floating point which we do not reproduce, and which are never
used as premises. The values are collected in Table~\ref{w3:tab:falsifiers}.

\emph{Rescalings.} The weights $c\Lambda$, $c>0$, keep everything that is homogeneous in the weight: support on
prime powers, nonnegativity, the formal Euler product of $\zeta^c$, monotonicity and every linear identity. They lose
the normalization, that is, the fact that the coefficient of $-\zeta'/\zeta$ at $p$ is exactly $\log p$, and with it
the functional equation. For every window $T>\frac12\log2$ we have $\lambda(T;c\Lambda)<0$ once $c$ is large enough
(Theorem~\ref{w3:bar:B1}\ref{w3:B1i}). For $c$ close to $1$ the evidence is computational: an interval computation
gives $\lambda(1;1.1\Lambda)\le-0.078$; in floating point the bottom crosses zero near $T\approx0.493$ for
$1.1\Lambda$ and near $T\approx0.645$ for $0.9\Lambda$; and at the window $\frac12\log15$ the reported upper bound
for $\lam$ of Section~\ref{w2:sec:certified}, if correct, makes $(1+\varepsilon)\Lambda$ negative as soon as
$\varepsilon>8\cdot10^{-69}$. The last figure shows the scale: the true weight lies at a distance of the order of the
margin from the boundary of the set of weights with nonnegative bottom (Section~\ref{w2:sec:margin}). An argument
that is homogeneous in the weight, or that has slack, or that deforms in the parameter $c$ without using the exact
value $c=1$, applies to these weights and therefore cannot give the sign.

\emph{Finite changes.} Deleting one prime power, or all $n<N$ (the weight $\Lambda\mathbf 1_{n\ge N}$), or adding
$c\mathbf 1_{n=n_0}$ with $c\ge0$, keeps nonnegativity, the support and all asymptotic information, and changes
finitely many values. Changing one coefficient by a large enough amount makes the bottom negative in every window
containing that contact (Theorem~\ref{w3:bar:B11}). For changes of moderate size we rely on interval computations of
Rayleigh quotients on explicit primitive test functions; they give negative values for the deletions of $n=3$, $7$
and $8$, for the empty comb, for $2\Lambda$ and $1.15\Lambda$, for the truncation $a_3$ defined below and for
the shift of the atom $3\mapsto3.2$, in the windows listed in Table~\ref{w3:tab:falsifiers}. Numerically, perturbations of relative size $10^{-5}$ already change the sign in
the windows where the margin is that small. Classical estimates are insensitive to finitely many coefficients, and
arguments monotone in the weight, or localized at one prime or one window, are insensitive to these changes; none of
them can give the sign.

\emph{Truncations of the Euler product.} For $P\ge2$ let
\[
  a_P(n)=\Lambda(n)\mathbf 1_{p\ge P}\qquad(n=p^k),
\]
the weight of $-\frac{d}{ds}\log\prod_{p\ge P}(1-p^{-s})^{-1}$. It is nonnegative, supported on prime powers, and
comes from an Euler product over a cofinite set of primes; the functional equation is lost. For
$T\le\frac12\log P$ every contact $n<e^{2T}\le P$ has all its prime factors below $P$, so $\aT[a_P]=\mT$ on such
windows, and $\lambda(T;a_P)$ is the bottom of the archimedean form on $\PT$. That bottom is negative for all
$T$ beyond a finite threshold $\Tarch$ (Proposition~\ref{w3:prop:archneg}). Hence $a_P$ is a falsifier in the windows
\[
  \Tarch<T\le\tfrac12\log P ,
\]
a range that is nonempty exactly when $P>e^{2\Tarch}$, in particular for all sufficiently large $P$. Numerically
$\Tarch\in[0.58,0.60]$, and \cite[Proposition~7.1]{P46} gives the bound $\Tarch\le0.7493$; since
$\frac12\log5\approx0.8047$, the range is nonempty for every $P\ge5$. For $P\le3$ it is empty
($\frac12\log3\approx0.549<0.58$; note $a_2=\Lambda$), and negativity of $a_3$ on primitive functions has to be
obtained otherwise: an interval computation gives $\lambda\le-0.505$ at $\frac12\log5+10^{-4}$. The truncations
satisfy the premises of every argument that uses the Euler product without the functional equation: cone conditions
on the logarithm of an Euler product, premises stable under removing finitely many primes, local functional
equations at finitely many places.

\emph{The Davenport--Heilbronn function.} Let $\chi$ be the character modulo $5$ with $\chi(2)=i$ and
\[
  L_{\mathrm{DH}}(s)=\tfrac{1-i\kappa}{2}L(s,\chi)+\tfrac{1+i\kappa}{2}L(s,\bar\chi),\qquad
  \kappa=\frac{\sqrt{10-2\sqrt5}-2}{\sqrt5-1}=0.2840790\ldots
\]
It has real coefficients, satisfies a functional equation $s\mapsto1-s$ with the gamma factor and conductor of
$L(s,\chi)$, and has zeros off the critical line, including zeros with real part greater than $1$
\cite{DavenportHeilbronn1936}, \cite[\S10.25]{Titchmarsh1986}. It has no Euler product. A computation in ball
arithmetic indicates that its window form is negative on primitive functions for all $T\ge9/5$. The premises
satisfied by $L_{\mathrm{DH}}$ are the functional equation, the shape of the gamma factor and the properties of
partial sums on the line; they underlie approximate functional equations, mean values, mollifiers, repulsion of
zeros, flows and virial identities, and realizations in Kre\u{\i}n spaces built from the functional equation alone.
Its partial sums on the critical line converge and do not reveal the off-line zeros, and the phase statistics
computed in \cite{P04} do not separate it numerically from $L$-functions with an Euler product. We use
$L_{\mathrm{DH}}$ only qualitatively, as a series with a functional equation and a negative window; comparisons of
the sizes of its window quantities with those of $\zeta$ turned out to be unreliable.

\emph{Epstein zeta functions.} The Epstein function of $x^2+5y^2$ (discriminant $-20$, class number $2$)
has nonnegative coefficients, a pole at $s=1$ and a functional equation with the gamma factor of a degree-two
$L$-function, and has infinitely many zeros off the critical line \cite{DavenportHeilbronn1936}. Up to
normalization it is the sum of two Euler products with the same gamma factor, namely $\zeta_K$ for
$K=\Q(\sqrt{-5})$ and the product $L(s,\chi_{-4})L(s,\chi_5)$ of the two genus $L$-functions, but it is not
itself an Euler product. An interval computation indicates that its complete form is negative for $T\ge1.59$; for
the form on primitive functions we have only a heuristic argument. The Epstein function satisfies the premises of
arguments from the functional equation and a lattice or theta-function representation which never use
multiplicativity, and of every argument that survives averaging two Euler products with the same gamma factor
(Observation~\ref{w3:bar:B9}).

\emph{The shifted product.} For $0<a<\frac12$ let $\Xi_a(s)=\xi(s+a)\xi(s-a)$. It satisfies
$\Xi_a(1-s)=\Xi_a(s)$, is the transform of a positive even kernel (the convolution of the twists
$\Phi(u)e^{\pm au}$ of the positive Riemann kernel $\Phi$ of $\xi$, \cite[\S10.1]{Titchmarsh1986}), has the Euler
product of $\zeta(s+a)\zeta(s-a)$ and the nonnegative weight $2\Lambda(n)\cosh(a\log n)$, and has infinitely many
zeros off the line $\Rel s=\frac12$, at real parts $\frac12\pm a$, by Hardy's theorem (all of its zeros lie off the
line if and only if no zero of $\zeta$ has real part $\frac12\pm a$). Every positive structure built only from the
symmetry $s\mapsto1-s$, positive kernels and positive weights is therefore present for $\Xi_a$ as well, and the Weil
form of $\Xi_a$ is negative on some compactly supported test function whose transform vanishes at the polar points
of $\Xi_a$ \cite[Proposition~9.3]{P47}; the proof given there adapts \cite[Theorem~6.2]{P45}. In particular no
positivity functional of Hodge-index type exists for it \cite[Theorem~9.4]{P47}. This excludes arguments by
compactness and passage to a limit, arguments from the symmetry alone, transfers of the Hodge index without an
arithmetic model, and separations by coefficients that do not read the exact normalization of the gamma factor.

\emph{Dedekind zeta functions and Artin $L$-functions.} The functions $\zeta_K$ for number fields $K$, and the Artin
$L$-functions of non-abelian ($S_3$) extensions, are expected to have nonnegative Weil forms. They are not
falsifiers; they test uniformity. An argument that would prove the sign for $\zeta$ by a property not shared by
$\zeta_K$ must name that property, and an argument uniform in the field must also apply to $\zeta_K$ and to
non-abelian Artin $L$-functions (Observation~\ref{w3:bar:B8}).

\emph{Beurling systems.} A Beurling system of generalized primes $1<q_1\le q_2\le\dots$ has a weight
$\Lambda_{\mathcal B}$ supported on the generalized prime powers. The system of Diamond,
Montgomery and Vorhauer \cite{DMV2006} has integer counting function $N(x)=cx+O(x^\theta)$ with $\frac12<\theta<1$,
a zeta function whose zeros approach the line $\Rel s=1$, and a prime-counting function with a remainder no better
than that of de la Vall\'ee-Poussin; Zhang \cite{Zhang2007} constructs systems satisfying the analogue of the Riemann
hypothesis and systems with large oscillation (see \cite{DiamondZhang2016,BDR2023}). For small windows we use a
discrete system: the primes $2,\frac{16}{5},5,7,11,\dots$, in which the prime $3$ is replaced by $16/5$. Here
$N_{\mathcal B}(x)=\frac{32}{33}x+O(\log^2x)$, and because $\frac{16}{5}\cdot5=2^4$ the generalized integers are
formal products counted with multiplicity. An interval computation gives a window form
$\le-0.2026\|F\|^2$ on an explicit primitive test function for all $T\ge\frac12\log5+10^{-4}$; in floating point the
form is negative already from $T\approx0.566$. Beurling systems satisfy the premises of every use of the convolution
identity $\Lambda\ast1=\log$ that does not involve the additive structure of $\Z$, of sieve and large-sieve arguments
valid for generalized integers, and of every square built prime by prime. In each construction we examined, the only
step at which $\N$ and the Beurling system separate is the passage from the formal multiplicative semigroup to the
additive lattice $\Z$; this is an observation about those constructions, not a theorem.

\emph{Finitely many zeros off the line.} For $\xi$ itself, under the hypothesis of exactly $m$ off-line quartets,
the explicit formula for $\psi(x)$ acquires an explicit defect term \cite[Theorem~3.3]{P19}; formally, the
hypothesis is realized by multiplying $\xi$ by a rational function that moves finitely many zeros off the line. On
the zero side each off-line quartet contributes one indefinite block, so under the same hypothesis the negative
index on even test functions is at most $m$ \cite[Corollary~6.5]{P18}, and on the full space it is exactly $2m$
\cite[Proposition~4.7]{P25}. These models, and quartets planted with small displacement, satisfy the premises of
arguments that read only finitely many zeros, finite inertia counts or patterns of finitely many window
observables.

\emph{Deformations, twisted weights and sign twists.} Deformations $f_t=\zeta+t\,g$ that keep the functional
equation and the coefficients $\Lambda(n)$ for $n$ below a given bound are a heuristic test for arguments that use
finitely many exact values together with integrality; we have not constructed them rigorously. The nonnegative
twisted weights $b_\vartheta(n)=\Lambda(n)\cdot\frac12\big(1+\cos(\gamma\log n+\vartheta)\big)\le\Lambda(n)$
satisfy every premise of positivity, support and upper-sieve type and reproduce oscillations of size $X^{1/2}$; we
expect that no bound for twisted prime sums below that size follows from those premises alone, and we know no
negative window for them. The comb with inverted sign, $a=-\Lambda$, and more generally the sign twists
$\epsilon(n)\Lambda(n)$ with $\epsilon(n)\in\{\pm1\}$ and $\epsilon(2)=-1$, has a negative complete form from
$T\approx0.38$ according to an interval computation, and is negative on primitive functions numerically from
$T\approx0.53$; it is a test for arguments by character sums and cocycles.

\emph{One modified Euler factor.} Computations indicate that a modification of a single Euler factor
of $\zeta$, keeping the gamma factor and conductor one, always produces a window with a negative value; we do not
have a proof. The function
$\zeta_{\mathrm{odd}}(s)=\sum_{n\ \mathrm{odd}}n^{-s}=(1-2^{-s})\zeta(s)$ has coefficients $0$ and $1$, a
rational Euler product and weight $0\le\Lambda_{\mathrm{odd}}\le\Lambda$. Its complete form, which contains the
zeros $2\pi ik/\log2$ of the removed factor on the line $\Rel s=0$, is negative for all $T\ge0.4608$ according to an
interval computation. This is a statement about the complete form. On primitive functions the window form of
$\zeta_{\mathrm{odd}}$ is $\aT[a_3]$: it equals $\mT$ for $T\le\frac12\log3$, hence is nonnegative there by the
numerical value of $\Tarch$; it is negative at $\frac12\log5+10^{-4}$ by the interval computation quoted above, and
in floating point already from $T\approx0.574$, below $\Tarch$.

\emph{Controls.} Three objects that should not be negative serve as calibration: $L(s,\chi_{-4})$, for which no
negative eigenvalue appears in any computed window; $\zeta$ itself; and curves over finite fields, where the Riemann
hypothesis holds and yet the frame of Hodge-index type constructed by analogy has a defect
\cite[Theorem~10.8]{P47}. The last control shows that a structural obstruction appearing for $\zeta$ need not
indicate a failure of positivity.

\begin{table}[htbp]
\centering\small
\caption{Windows in which a negative value is known for the test objects. The window given is the smallest one for
which the negative value has been obtained by the indicated means; it is not claimed to be optimal. ``Complete''
refers to the complete form, ``primitive'' to the form on $\PT$.}\label{w3:tab:falsifiers}
\begin{tabular}{llll}
\toprule
Test object & Negative window & Form & Source\\
\midrule
$c\Lambda$, $c>c_0(T)$ & every $T>\frac12\log2$ & primitive & Theorem~\ref{w3:bar:B1}\ref{w3:B1i}\\
$1.1\Lambda$ & $T=1$, $\lambda\le-0.078$ & primitive & interval arithmetic\\
$1.1\Lambda$; $0.9\Lambda$ & from $T\approx0.493$; $T\approx0.645$ & primitive & floating point\\
$1.15\Lambda$, $2\Lambda$ & $T=1$ & primitive & interval arithmetic\\
$\Lambda+c\mathbf 1_{n_0}$, $c$ large & every $T>\frac12\log n_0$ & primitive & Theorem~\ref{w3:bar:B11}\\
deletion of $n=3$ & $T=0.75$; $\frac12\log5+0.01$ & primitive & interval arithmetic\\
deletion of $n=7$; of $n=8$ & $T=1$; $T=\log3$ & primitive & interval arithmetic\\
empty comb & $T=\frac12\log5$; $T=1$ & primitive & interval arithmetic\\
$a_P$, $P\ge5$ & $\Tarch<T\le\frac12\log P$ & primitive & Theorem~\ref{w3:bar:B1}\ref{w3:B1iii}\\
$a_3$ (also $\zeta_{\mathrm{odd}}$) & $\frac12\log5+10^{-4}$, $\lambda\le-0.505$ & primitive & interval arithmetic\\
atom $3$ shifted to $3.2$ & some window & primitive & interval arithmetic\\
$a_3$ & from $T\approx0.574$ & primitive & floating point\\
$L_{\mathrm{DH}}$ & $T\ge9/5$ & primitive & ball arithmetic\\
Epstein, $x^2+5y^2$ & $T\ge1.59$ & complete & interval arithmetic\\
$\Xi_a$ & some window & primitive for $\Xi_a$ & \cite[Proposition~9.3]{P47}\\
Beurling $2,\frac{16}5,5,7,\dots$ & $T\ge\frac12\log5+10^{-4}$, $\le-0.2026$ & primitive & interval arithmetic\\
Beurling $2,\frac{16}5,5,7,\dots$ & from $T\approx0.566$ & primitive & floating point\\
$-\Lambda$ & from $T\approx0.38$ & complete & interval arithmetic\\
$-\Lambda$ & from $T\approx0.53$ & primitive & floating point\\
$\zeta_{\mathrm{odd}}$ & $T\ge0.4608$ & complete & interval arithmetic\\
\bottomrule
\end{tabular}
\end{table}

The falsifiers share a great deal of structure with $\zeta$: the Euler product ($a_P$, Beurling systems,
$\Xi_a$), the functional equation ($L_{\mathrm{DH}}$, the Epstein function, $\Xi_a$), finitely many exact
coefficients (deformations, finite changes). A list of premises that excludes all of them at once must contain the
Euler product over all primes, the functional equation with the gamma factor of $\zeta$ at conductor one and the
exact integer structure of the multiplicative semigroup, and must use them jointly; the resulting constraints on a
proof are collected in Section~\ref{w8:sec:constraints}. In practice a proposed construction can be tested against a
short list: a finite erasure, a single amplified coefficient, $c\Lambda$ on both sides of $c=1$, a truncation $a_P$
with $P\ge5$, $L_{\mathrm{DH}}$, the Epstein function of $x^2+5y^2$, a planted off-line quartet, a deformation with
finitely many exact coefficients, a discrete Beurling system and a twisted weight $b_\vartheta$. For each of them
one should be able to name the step of the construction that fails.

\subsection{The barrier theorems proved here}
\label{w3:sec:barriers}\label{w3:subsec:lemmas}\label{w3:subsec:proved}

By a barrier we mean a statement of the form: no argument of the class $\mathcal C$ proves $\lam\ge0$. By
Proposition~\ref{w3:prop:principle} it suffices to exhibit an element of $\mathcal A(\mathcal C)$ with a negative
window; the difficulty lies in describing the class. In this subsection we prove four such statements, each with
its class stated explicitly, and an auxiliary fact about the archimedean form. Throughout, $T>0$ is a window,
$\NT$ its contact set, and for a finitely supported real sequence $w=(w_n)_{n\in\NT}$ we write
\[
  \mathfrak p[w](F)=\sum_{n\in\NT}w_n\,\Rel\cF(\log n)
\]
for the comb form with weights $w$; thus $\pT=\mathfrak p[v]$ with $v_n=2\Lambda(n)n^{-1/2}$. The proofs rest on
the following lemma, which produces a primitive function that sees a single contact.

\begin{lemma}[isolating one contact]\label{w3:lem:isolate}
Let $n_0\in\NT$ and $s\in\{1,-1\}$. There is a real $F\in C_c^\infty(\IT)\cap\PT$, $F\ne0$, such that for every
finitely supported $w$ on $\NT$,
\[
  \mathfrak p[w](F)=s\,w_{n_0}\,\tfrac12\|F\|_2^2 .
\]
\end{lemma}

\begin{proof}
Put $u=\log n_0<2T$ and choose $\varepsilon>0$ with $\varepsilon<T-\frac u2$, $2\varepsilon<\log2$ and
$2\varepsilon<|\log m-\log n_0|$ for every $m\in\NT\setminus\{n_0\}$ (a finite set). The two linear
functionals $\varphi\mapsto\int\varphi(x)e^{\pm x/2}dx$ on the infinite-dimensional space of real functions in
$C_c^\infty(-\varepsilon,\varepsilon)$ have a common nonzero element $\varphi$ of their kernels. Put
$F(t)=\varphi(t+\frac u2)+s\,\varphi(t-\frac u2)$. Its support lies in $(-\frac u2-\varepsilon,\frac u2+\varepsilon)
\subset\IT$, and $A(F)=(e^{-u/4}+se^{u/4})\int\varphi(x)e^{x/2}dx=0$, likewise $B(F)=0$; so $F\in\PT$. The two
translates have disjoint supports (since $u\ge\log2>2\varepsilon$), hence $\|F\|_2^2=2\|\varphi\|_2^2$, and
\[
  \cF(v)=2c_\varphi(v)+s\,c_\varphi(v-u)+s\,c_\varphi(v+u),
\]
where $c_\varphi$ vanishes outside $(-2\varepsilon,2\varepsilon)$ and $c_\varphi(0)=\|\varphi\|_2^2$. For
$m\in\NT$ we have $\log m\ge\log2>2\varepsilon$ and $\log m+u>2\varepsilon$, while
$|\log m-u|<2\varepsilon$ only for $m=n_0$. Hence $\cF(\log m)=0$ for $m\ne n_0$ and
$\cF(\log n_0)=s\|\varphi\|_2^2=\frac s2\|F\|_2^2$.
\end{proof}

\begin{proposition}[the archimedean form is indefinite in large windows]\label{w3:prop:archneg}
Let $\mu(T)=\inf\{\mT(F):F\in\PT,\ \|F\|_2=1\}$. Then $\mu$ is nonincreasing in $T$, and there is $L_0<\infty$
such that $\mu(T)\le\frac12C_0<0$ for all $T\ge L_0$, where $C_0=m(0)=\psi(\frac14)-\log\pi$. In particular
$\Tarch:=\inf\{T:\mu(T)<0\}$ is finite and $\mu(T)<0$ for every $T>\Tarch$.
\end{proposition}

\begin{proof}
Extension by zero maps $\mathcal P_{T}$ into $\mathcal P_{T'}$ for $T<T'$ and does not change $\mT$, which
depends only on $\Fh$; so $\mu$ is nonincreasing. Fix three linearly independent real
$g_1,g_2,g_3\in C_c^\infty(-1,1)$ and, for $1\le L\le T$, put $G_j(t)=L^{-1/2}g_j(t/L)$, supported in $\IT$.
Then $\langle G_i,G_j\rangle=\langle g_i,g_j\rangle$ and $\widehat G_j(\tau)=L^{1/2}\widehat g_j(L\tau)$, so
\[
  \mT(G_i,G_j)=\frac1{2\pi}\int_\R m(\sigma/L)\,\widehat g_i(\sigma)\overline{\widehat g_j(\sigma)}\,d\sigma .
\]
Since $m$ is continuous and $|m(\tau)|\le C\log(2+|\tau|)$, we have $|m(\sigma/L)|\le C\log(2+|\sigma|)$ for
$L\ge1$, and the $\widehat g_j$ are Schwartz functions; by dominated convergence and Plancherel the $3\times3$
matrix $(\mT(G_i,G_j))$ tends to $C_0(\langle g_i,g_j\rangle)$ as $L\to\infty$. As the Gram matrix is positive
definite, for $L\ge L_0$ we get $\mT(G)\le\frac12C_0\|G\|_2^2$ for every $G$ in the span of $G_1,G_2,G_3$. The
two conditions $A(G)=B(G)=0$ cut this three-dimensional span in a nonzero subspace, whose elements belong to
$\PT$. Hence $\mu(T)\le\frac12C_0$ for $T\ge L_0$.
\end{proof}

A machine computation, which we do not reproduce, gives $\Tarch\in[0.58,0.60]$. The threshold concerns the
archimedean part alone and is not a first-contact value for $\lam$. For $T\le\frac12\log2$ the contact set is empty,
$\aT=\mT$, and positivity in those windows is classical \cite{Yoshida1992}. The proposition shows that the
archimedean form alone cannot carry positivity beyond a finite window, so the primes must enter every proof for
$T>\Tarch$.

\begin{theorem}[every finite comb is indefinite on primitive functions]\label{w3:bar:B6}
For every $T>\frac12\log2$ and every nonzero finitely supported real $w$ on $\NT$, the comb form
$\mathfrak p[w]$ takes both signs on $\PT$; its Rayleigh quotient on $\PT$ reaches $\pm\frac12|w_{n_0}|$ for each
$n_0$ in the support.
\end{theorem}

\begin{proof}
Apply Lemma~\ref{w3:lem:isolate} at any $n_0$ with $w_{n_0}\ne0$, with both signs $s$.
\end{proof}

In particular positivity cannot be localized at a single prime power, or at any finite part of the comb taken
alone: each piece is indefinite on $\PT$, and the sign can only come from the interaction of the whole comb with
the archimedean form.

\begin{theorem}[rescalings, identities and truncations]\label{w3:bar:B1}
Let $T>\frac12\log2$.
\begin{enumerate}[label=(\roman*)]
\item\label{w3:B1i} There is $c_0(T)<\infty$ with $\lambda(T;c\Lambda)<0$ for all $c>c_0(T)$.
Hence no class of arguments whose admissible weights contain $\{c\Lambda:c>0\}$, in particular no class closed
under multiplication of the weight by positive constants, proves $\lam\ge0$.
\item\label{w3:B1ii} There is no identity
$\aT[a](F)=Q_a(F)+\rho_a(F)$ on $\PT$, with $Q_a\ge0$ and $\rho_a\ge0$, valid for every weight $a$ in a set
containing $\{c\Lambda:c>0\}$; in particular no identity valid for every nonnegative weight carries a sign.
\item\label{w3:B1iii} For every $P$ with $\frac12\log P>\Tarch$ and every
$T\in(\Tarch,\frac12\log P]$, $\lambda(T;a_P)<0$; such $P$ exist. Hence no class of arguments whose admissible
weights contain the truncations $a_P$ for all large $P$ (for instance, arguments using nonnegativity, support
on prime powers and an Euler product over a cofinite set of primes) proves $\lam\ge0$ in every window.
\end{enumerate}
\end{theorem}

\begin{proof}
\ref{w3:B1i}: $2\in\NT$. Take $F$ from Lemma~\ref{w3:lem:isolate} with $n_0=2$, $s=1$; then
$\aT[c\Lambda](F)=\mT(F)-\frac c2v_2\|F\|_2^2$, $v_2=\sqrt2\log2$, which is negative for
$c>c_0=2\mT(F)/(v_2\|F\|_2^2)$. The second sentence follows from Proposition~\ref{w3:prop:principle}.
\ref{w3:B1ii}: such an identity would give $\lambda(T;c\Lambda)\ge0$ for all $c$, contradicting \ref{w3:B1i}.
\ref{w3:B1iii}: for $T\le\frac12\log P$ every $n<e^{2T}$ has all prime factors $<P$, so $\aT[a_P]=\mT$ and
$\lambda(T;a_P)=\mu(T)<0$ for $T>\Tarch$ by Proposition~\ref{w3:prop:archneg}.
\end{proof}

We expect the same conclusion for a much larger family
of modifications of the comb: rescalings on either side of $c=1$, erasure of initial segments, deletion or
amplification of a single term, shift of an atom, replacement by a Beurling system, sign twists; about fifteen
elementary operations in all, two of which seem to require a non-degeneracy hypothesis on the minimizer. We have
not proved this; for rescalings by $1\pm\delta$ with small $\delta$ and for a single deletion at a specific window
the evidence is the computed negative values of Table~\ref{w3:tab:falsifiers}. In either form the statement means
that a proof must use a property that fails for each modified weight.

\begin{theorem}[one coefficient beyond any finite list]\label{w3:bar:B11}
Let $N\ge2$, let $q\ge N$ be an integer and let $T>\frac12\log q$. For each sign $s\in\{1,-1\}$ there is a weight
$a$ with $a(n)=\Lambda(n)$ for all $n\ne q$, $s\,a(q)$ arbitrarily large, and $\lambda(T;a)<0$; for $s=1$ the
weight is nonnegative. Consequently no argument that uses only the values $\Lambda(n)$, $n<N$, together with
properties of the weight preserved by changing one coefficient $a(q)$ with $q\ge N$ by a large amount, proves
$\lam\ge0$ in any window $T>\frac12\log q$.
\end{theorem}

\begin{proof}
Since $q<e^{2T}$, $q\in\NT$. Take $F$ from Lemma~\ref{w3:lem:isolate} with $n_0=q$ and the sign $s$; then
$\aT[a](F)=\mT(F)-\frac s2v_q(a)\|F\|_2^2$, independent of $a(n)$ for $n\ne q$, and negative once
$s\,a(q)>\sqrt q\,\mT(F)/\|F\|_2^2$.
\end{proof}

The theorem uses a single amplified coefficient, and does not cover arguments which also use upper bounds such as
$0\le a(n)\le\log n$. We expect sharper statements to hold. First, no finite list of exact linear data of $\Lambda$
should suffice for all windows. Second, under asymptotics of Landau--Widom type for the time--frequency
concentration operators \cite{LandauWidom1980}, an averaging over windows indicates that positivity for all $T$
forces the contribution of the primes entering at a window $T$ to cancel the symbol $\sigma_{T_0}$ of an earlier
window, with exact sign and magnitude, at frequency resolution $\pi/T$; the converse fails. Third, a lower bound
$\lam\ge-Ce^{\kappa T}$ should force the abscissa of the zeros to deviate from $\frac12$ by at most $O(\kappa)$,
and this should be sharp. None of the three is proved here or used.

\begin{theorem}[criteria on the low part of the symbol]\label{w3:bar:B4}
Fix $T>0$ and put $\eta_0(T)=\sup\{-q(F)/\|F\|_2^2:\,F\in\PT\setminus\{0\}\}$, where
$q(F)=\frac1{2\pi}\int\min(\sT,0)|\Fh|^2$. Then $\eta_0(T)>0$, and for every $0<\eta<\eta_0(T)$ the quadratic form
with symbol $s_\eta=\min(\sT,\eta)$,
\[
  F\mapsto\frac1{2\pi}\int_\R\min(\sT(\tau),\eta)\,|\Fh(\tau)|^2\,d\tau ,
\]
takes negative values on $\PT$, although $s_\eta$ has the same negative set as $\sT$ and coincides with $\sT$ on
the sublevel set $\{\sT\le\eta\}$. Hence no criterion which, applied to an arbitrary continuous symbol, certifies
nonnegativity of its form on $\PT$ from the restriction of the symbol to a sublevel set $\{\sT\le\eta\}$ with
$\eta<\eta_0(T)$ can certify $\lam\ge0$.
\end{theorem}

\begin{proof}
$\sT$ is continuous and $\sT(0)=C_0-\sum_{n\in\NT}v_n<0$, so $\{\sT<0\}$ contains an interval. For $F\ne0$ in
$L^2(\IT)$ the transform $\Fh$ is entire and not identically zero, so it vanishes only on a discrete set and
$q(F)<0$; $\PT\ne\{0\}$, hence $\eta_0(T)>0$. Given $\eta<\eta_0(T)$ choose $F\in\PT$ with
$-q(F)>\eta\|F\|_2^2$. Since $\min(\sT,\eta)\le\min(\sT,0)+\eta$, Plancherel gives
$\frac1{2\pi}\int\min(\sT,\eta)|\Fh|^2\le q(F)+\eta\|F\|_2^2<0$.
\end{proof}

\subsection{Further barriers and what they have in common}
\label{w3:subsec:announced}\label{w3:subsec:mechanism}\label{w3:tab:barriers}

The statements of this subsection have the same form as the theorems above, but, with one exception
(Observation~\ref{w3:bar:B5}), we cannot prove them in the generality in which we state them. Each is given with the
evidence we have for it, and none is used as a premise.

\begin{observation}[linear certificates in a fixed basis]\label{w3:bar:B2}
Certificates of positivity that are linear in a fixed finite basis, such as linear programming bounds,
congruence to a diagonal form, Schur complements and inertia computations in a basis chosen independently of
$T$, certify finitely many windows but do not decide a new window: the dimension needed grows with
$T$. For the two-sector Schur reduction our computations indicate growth at least like a negative power of $\lam$.
We have observed this for the specific bases used; as a general statement it is heuristic.
\end{observation}

\begin{observation}[determinants]\label{w3:bar:B3}
Every determinant whose sign we have considered as a criterion (of a finite Galerkin block, of a Schur complement,
of a Fredholm or Feshbach reduction) factors as $\lam\cdot R$ with $R$ of fixed sign in the regime considered, so it
contains no information beyond $\lam$. Bounds for the determinant with a multiplicative error $O(1)$ do not decide
the sign, because $\lam$ is exponentially small (Section~\ref{w2:sec:margin}).
\end{observation}

\begin{observation}[moments, patterns, bounded-range families]\label{w3:bar:B4b}
We expect that no criterion depending on finitely many traces or moments of the window operator, no criterion
depending on a pattern of finitely many window observables, and no functional of $\{\lambda(T'):|T'-T|\le R\}$ for
bounded $R$, decides the sign: in each case a model sharing the data with $\Lambda$ (finitely many zeros off the
line, planted quartets, deformations) should have a negative window. Two-sided asymptotic laws, which constrain
$\lambda(T)$ as $T\to\infty$ from both sides, are not concerned by these statements.
\end{observation}

\begin{observation}[finite zero sets]\label{w3:bar:B5}
In the dual formulation, positivity of the complete form in a window $T\ge\frac12\log2$ is certified by an even
weight $w$ equal to $1$ on $[-2T,2T]$ for which the function $f_w(\tau)=\sum_\rho\widehat w(\tau-\gamma_\rho)$, the
smoothing of the zeros by $\widehat w$, is nonnegative \cite[Theorems~6.2 and~8.2]{P46}. No finite set of zeros
carries such a certificate: if $Z$ is a nonempty finite symmetric multiset of real ordinates and $w$ is even,
continuous, real and integrable with $w=1$ on $[-2T,2T]$, then $\sum_{\gamma\in Z}\widehat w(t-\gamma)$ takes
negative values \cite[Theorem~6.5]{P46}. In particular no certificate acts term by term on the zeros. This is the
reason why the zero side of the explicit formula~\eqref{w1:eq:explicit} cannot be used as a source of the sign.
\end{observation}

\begin{observation}[sensitivity]\label{w3:bar:B7}
The sign of $\lam$ is not a bounded (Lipschitz) functional of the weights $(v_n)$ or of the symbol $\sT$:
numerically, relative changes of size $10^{-5}$ in a single coefficient change it in windows near $T=1$. More
precisely, the true weight lies at a distance of order $\lam$ from the boundary of the convex set of weights with
nonnegative bottom. The set is convex because $a\mapsto\lambda(T;a)$ is an infimum of affine functions, and
$\lambda(T;(1+\varepsilon)\Lambda)\le\lam-\varepsilon\,\pT(F_*)$ at a minimizer $F_*$, with
$\pT(F_*)=\mT(F_*)-\lam$; see Section~\ref{w2:sec:margin}.
\end{observation}

\begin{observation}[insufficient axioms; uniformity in the field]\label{w3:bar:B8}
(i) A list of axioms consisting of the functional equation with a gamma factor, a pole at $s=1$, nonnegativity
of coefficients and an integral (theta-function) representation is satisfied by the Epstein function of
$x^2+5y^2$. Granting the computation which gives a negative complete form for $T\ge1.59$, these axioms do not
suffice, and a proof must use the Euler product, the nonnegativity of $\Lambda$ or the absence of zeros on
$\Rel s\ge1$ in an essential step. (ii) A proof that uses only properties shared by Dedekind zeta functions
$\zeta_K$ and Artin $L$-functions of $S_3$-extensions must be uniform in the field; otherwise it must name a
property separating $\zeta$ from them. The stronger assertion that a proof must distinguish $\zeta$ from $\zeta_K$
is not justified: a field-uniform proof is admissible. Part (ii) is a constraint of method, not a theorem.
\end{observation}

\begin{observation}[averages of Euler products]\label{w3:bar:B9}
No inequality that is linear or convex in the Dirichlet series separates two Euler products with the same gamma
factor from their average. Indeed, if a premise is a convex condition on $L$ (on its coefficients, its values on
a line, or its functional equation data) and holds for $\zeta_K$ and for $L(s,\chi_{-4})L(s,\chi_5)$, it holds for
their average, which is, up to normalization, the Epstein function of $x^2+5y^2$ and has zeros off the line
\cite{DavenportHeilbronn1936}. This part is immediate. The conclusion that such premises cannot prove
positivity requires in addition the negativity of the window form of the Epstein function, for which the evidence
is the computation quoted in Section~\ref{w3:subsec:models}.
\end{observation}

\begin{observation}[imported positivity]\label{w3:bar:B10}
No positivity imported from a construction that is uniform in the comb (positive kernels, completely monotone
functions, representations of Bochner type, Gram matrices of fixed families, correlation inequalities, semigroups
preserving positivity) equals $\aT$ on $\PT$. (i) A positive object independent of the weights cannot coincide with
a form that is negative for $c\Lambda$ with $c$ large. (ii) A positive object depending linearly on the weights is
excluded by Theorem~\ref{w3:bar:B1}\ref{w3:B1ii}. (iii) One depending on the weights through a monotone operation
fails for the deletions, granting the computed values of Table~\ref{w3:tab:falsifiers}. (iv) In the remaining,
non-uniform cases we have only heuristic arguments that the construction reduces to the problem itself. Items (i)
and (ii) follow from Theorem~\ref{w3:bar:B1} for any precise reading of ``uniform in the comb'' that contains the
rescaled weights. An auxiliary positive object is admissible when its arithmetic identity with the form is proved.
\end{observation}

\begin{observation}[orientation of the prime-side inequalities]\label{w3:bar:B12}
The unconditional tools on the prime side (prime number theorem with remainder, sieves, large sieve, mean-value
theorems, zero-density estimates) give lower bounds for the comb $\pT$, or two-sided bounds with an error, whereas
positivity of $\aT=\mT-\pT$ needs an upper bound for $\pT$ exact to the size of the margin. In the eight methods of
this kind that we examined, the inequality produced has the wrong orientation. This is a heuristic statement.
\end{observation}

\begin{observation}[both primitivity conditions are needed]\label{w3:bar:B6b}
Theorem~\ref{w3:bar:B6} shows that the comb alone is indefinite on $\PT$. In the other direction, the two
conditions $A=B=0$ cannot be dropped: computations indicate that no representation of the form ``sum of squares
plus a nonnegative potential'' holds on all of $L^2(\IT)$ for $T\in[0.2,0.75]$. Primitivity removes only two
directions, while the negative space of the archimedean part on $\PT$ grows linearly in $T$
(Section~\ref{w5:sec:schur}).
\end{observation}

Three facts lie behind all the statements of this section. The first is the size of the margin: the numerical
values of Section~\ref{w2:sec:margin} and the upper bound $\lam\le4\cdot10^9\,Te^{-15.449T}$ for $T\ge1$
\cite[Corollary~9.11]{P46} leave no room for estimates with slack; this is what stops linear certificates in a fixed
basis, determinant criteria and the prime-side estimates (Observations~\ref{w3:bar:B2}, \ref{w3:bar:B3},
\ref{w3:bar:B7} and~\ref{w3:bar:B12}). The second is that the form sees each value $\Lambda(n)$, $n<e^{2T}$,
exactly: changing one coefficient, rescaling or truncating destroys the sign (Theorems~\ref{w3:bar:B1},
\ref{w3:bar:B6} and~\ref{w3:bar:B11}). The third is that the falsifiers share finite data with $\zeta$: finitely
many zeros, finitely many coefficients, the functional equation, or the Euler product over a cofinite set
(Theorem~\ref{w3:bar:B4}, Observations~\ref{w3:bar:B4b}, \ref{w3:bar:B5}, \ref{w3:bar:B8} and~\ref{w3:bar:B9}).
What is excluded are classes of arguments. An argument that uses the Euler product over all primes, the functional
equation of $\zeta$ and the integer structure of $\N$ jointly and non-linearly is not excluded, and neither is a
two-sided asymptotic law in $T$; we return to these in Section~\ref{w8:sec:open}.

\section{Squares, local decompositions and imported positivity}\label{w4:sec:squares}

We now turn to the approaches themselves, beginning with the three most direct ones: to write $\aT$ on $\PT$ as a sum of
squares with a manifestly nonnegative remainder, to split it into local pieces each of which has a sign, and to transport
to it a positivity theorem from another field. The three meet the same difficulty. By an elementary identity the Weil
form is a Dirichlet form minus a constant, and every positivity that such a structure supplies holds for all nonnegative
weights, among them the test weights of Table~\ref{w3:tab:falsifiers} for which the form is negative. Numerical values
quoted in this section without a reference come from computations that were not certified; none of them is a premise of a
stated result.

Throughout, $\aT[a]$ denotes the truncated form with $\Lambda$ replaced by a weight sequence $a=(a_n)_{n\ge2}$ in
$v_n$, and $V_uF(t)=F(t-u)$ is translation on $\R$ (functions on $\IT$ are extended by zero).

\subsection{Squares and combinations}\label{w4:fam:R1}

The most direct way to prove $\lam\ge0$ is to exhibit, for every $F\in\PT$, an identity
\begin{equation}\label{w4:eq:sosshape}
\aT(F)=\sum_k\|S_kF\|^2+R(F),\qquad R(F)\ge0\ \text{manifestly},
\end{equation}
with operators $S_k$ and a remainder $R$ written in closed form. Four elementary facts frame every attempt of this kind.

\begin{proposition}[L\'evy--Khintchine form of the Weil form]\label{w4:prop:levy}
Put $K(u)=e^{u/2}/\sinh u$ for $u>0$. For every real weight sequence $a$, every $T>0$ and every $F\in\DT$,
\begin{equation}\label{w4:eq:levy}
\aT[a](F)=\frac12\int_0^\infty K(u)\,\|F-V_uF\|^2\,du+\sum_{n\in\NT}\frac{v_n}{2}\,\|F-V_{\log n}F\|^2-\kappa_T[a]\,\|F\|^2,
\end{equation}
where $v_n=2a_nn^{-1/2}$ and $\kappa_T[a]=\sum_{n\in\NT}v_n-C_0$. Write $\mathcal E_T[a](F)$ for the sum of the first two terms on the
right. In particular, when $a\ge0$ the form $\mathcal E_T[a]=\aT[a]+\kappa_T[a]\|\cdot\|^2$ is
the Dirichlet form of the symmetric L\'evy process on $\R$ with L\'evy measure $\frac12\big(K(|u|)\,du+\sum_{n\in\NT}v_n(\delta_{\log n}+\delta_{-\log n})\big)$,
killed on leaving $\IT$ (that is, restricted to functions supported in $\IT$), and the window problem for the weight $a$ is the Poincar\'e-type inequality
$\mathcal E_T[a](F)\ge\kappa_T[a]\|F\|^2$ on $\PT$.
\end{proposition}

\begin{proof}
Gauss's integral $\psi(z)=\int_0^\infty\big(e^{-x}/x-e^{-zx}/(1-e^{-x})\big)\,dx$ ($\Rel z>0$) gives
$\psi(z)-\psi(\tfrac14)=\int_0^\infty(e^{-x/4}-e^{-zx})(1-e^{-x})^{-1}dx$. Taking $z=\frac14+\frac{i\tau}2$, real parts, and substituting $x=2u$,
\[
m(\tau)-C_0=\int_0^\infty\frac{2e^{-u/2}}{1-e^{-2u}}\,(1-\cos\tau u)\,du=\int_0^\infty K(u)(1-\cos\tau u)\,du ,
\]
the integral converging since $K(u)\sim1/u$ at $0$ and $K(u)\sim2e^{-u/2}$ at $\infty$. By Plancherel,
$\frac1{2\pi}\int(1-\cos\tau u)|\Fh(\tau)|^2d\tau=\|F\|^2-\Rel\cF(u)=\frac12\|F-V_uF\|^2$. Integrating against $K$ (Tonelli; the integrand is
nonnegative) gives $\mT(F)=C_0\|F\|^2+\frac12\int_0^\infty K(u)\|F-V_uF\|^2du$ for $F\in\DT$. For the prime part,
$\Rel\cF(\log n)=\|F\|^2-\frac12\|F-V_{\log n}F\|^2$. Subtracting gives \eqref{w4:eq:levy}.
\end{proof}

For the weight $\Lambda$ this is also \cite[Theorem~3.3]{P46}. Two consequences are used repeatedly below. First, every
positivity property of Markov or Dirichlet-form type (positivity preservation, Beurling--Deny contraction, curvature
bounds for the jump generator) holds for \emph{every} nonnegative weight, hence for the scaled weights $c\Lambda$ with
$c>1$ and for the other test weights of Table~\ref{w3:tab:falsifiers}; it cannot by itself carry the sign. Second,
increasing the weight increases both the Dirichlet form and the constant $\kappa_T[a]$; whether positivity survives is
decided by the balance between the two, which is what a proof must control. We note also that $\aT[a](F)$ depends on $F$
only through $\cF$ and is therefore invariant under translation, while translation multiplies $A(F)$ and $B(F)$ by
$e^{\pm u/2}$; so $\PT$ may be replaced by the primitive functions supported in any interval of length $2T$.

\begin{proposition}[the tautological square]\label{w4:prop:taut}
There is a densely defined operator $U$ on $\PT$, depending on $\Lambda$, with $\aT(F)=\|UF\|^2$ for all $F$ in the form domain, if and only
if $\lam\ge0$.
\end{proposition}

\begin{proof}
Since $m(\tau)\ge\frac12\log(2+|\tau|)-c$ for a constant $c$ and the prime part is bounded on $L^2$, $\aT$ is a closed semibounded
form on $\DT$, and $\DT\cap\PT$ is closed for the form norm because $A$ and $B$ are continuous for it. If $\lam\ge0$, this form is represented by a nonnegative self-adjoint operator $H$ (Friedrichs), and
$U=H^{1/2}$ works. The converse is trivial.
\end{proof}

Thus a square whose ingredients may contain $\Lambda$ without an independent formula is a restatement of the problem, and
an approach by squares has to name its ingredients explicitly.

\begin{proposition}[polarization]\label{w4:prop:polar}
Let $L$ be a real-linear map from a space of weights into a real Hilbert space. For all weights $w$ and every reference weight $w_0$,
\[
\|Lw\|^2=2\langle Lw,Lw_0\rangle-\|Lw_0\|^2+\|L(w-w_0)\|^2 .
\]
\end{proposition}

\begin{proof}
Expand $\|L(w-w_0)\|^2$.
\end{proof}

A quantity that is quadratic in the weight can therefore always be written as a term linear in the weight plus a
nonnegative remainder. The remainder is nonnegative for every $w$, so whatever sign it supplies is supplied equally to
every test weight. Since $\aT[a]$ is affine in $a$, one expects that a square which is quadratic in the weight can always
be rewritten through such a decomposition; this is a heuristic, not a theorem, but it holds in each construction we
examined. The decision then lies entirely in the linear part, to which the dual-certificate analysis of
Section~\ref{w5:sec:dual} applies.

\begin{proposition}[certificates valid for a test weight]\label{w4:prop:controls}
Let $\mathcal C$ be a set of weights containing $\Lambda$ and let $T>0$. If for every $a\in\mathcal C$ there are operators $S_k[a]$ and a form
$R[a]\ge0$ on $\PT$ with $\aT[a](F)=\sum_k\|S_k[a]F\|^2+R[a](F)$ for all $F\in\PT$ in the form domain, then $\lambda(T;a)\ge0$ for
every $a\in\mathcal C$. Hence, if $\lambda(T;a')<0$ for some $a'\in\mathcal C$, no such family of identities exists on $\mathcal C$.
\end{proposition}

The proof is immediate. The proposition is applied as follows. An argument whose steps use only properties shared by all
members of $\mathcal C$ produces such a family of identities, and so cannot exist for $\Lambda$ when $\mathcal C$ contains
a weight with $\lambda(T;a')<0$. This is a principle about proofs and not a theorem; its content lies in the choice of
$\mathcal C$. We use three test weights: the scaled weight $c\Lambda$ (typically $c=1.1$), the deletion of one prime
power, and a Beurling-type system with generators $2,\frac{16}5,5,7,11,\dots$; Table~\ref{w3:tab:falsifiers} says in
which windows, and at which level of verification, each is known to be negative. A construction that is linear in
$\Lambda$ with signless coefficients is valid for $c\Lambda$; one that uses only the formal multiplicative semigroup is
valid for the Beurling system; one that uses only positivity of coefficients is valid after deletion. A useful
certificate must therefore break, at a named step, under each of the three. Two cautions apply. An identity $\aT=Q$ with
$Q\ge0$ derived from an explicit arithmetic source need not lose positivity at a test weight if its \emph{proof} is not
valid there, so the principle excludes only certificates whose proof transfers. And the Schur complement of a Gram block
is a positive projection residue, which does not exclude a Gram identity exactly identified with $\aT$ by a proved
arithmetic formula.

\phantomsection\label{w4:tab:sos}%
We come to the squares themselves. The natural ingredients are those of the arithmetic of $\Lambda$: the convolution
identity $\Lambda*1=\log$, gcd and Smith matrices \cite{Smith1876}, plane-wave and theta families, interpolation
formulas, and translation operators producing the atoms of the prime part at $\log n$. About forty constructions of
this type have been examined. One identity among them is exact and general.

\begin{proposition}[convolution identity]\label{w4:prop:conv}
Let $a_\infty(F)=\mT(F)-\sum_{n\ge2}v_n\Rel\cF(\log n)$ for compactly supported $F$ in the form domain (the sum is finite and
$a_\infty=\aT$ on functions supported in $\IT$). For real $F,G$ with compact support, $G$ integrable, and $\Delta_sF=F(\cdot+\frac s2)-F(\cdot-\frac s2)$ and $c_G(s)=\int_\R G(t)G(t+s)\,dt$ the autocorrelation of $G$,
\[
a_\infty(F*G)=\widehat G(0)^2\,a_\infty(F)-\frac12\int_{\R}c_G(s)\,a_\infty(\Delta_sF)\,ds .
\]
If $F$ is primitive then so are $F*G$ and every $\Delta_sF$; if $\supp F\subset\overline{I_T}$ and $\supp G\subset[-S,S]$, then $F*G$ is
supported in $\overline{I_{T+S}}$ and $\Delta_sF$ in $\overline{I_{T+|s|/2}}$.
\end{proposition}

\begin{proof}
$a_\infty(F)=\frac1{2\pi}\int\sigma|\Fh|^2$ with $\sigma(\tau)=m(\tau)-\sum_nv_n\cos(\tau\log n)$, understood on compactly supported $F$.
Since $|\widehat G(\tau)|^2=\int c_G(s)e^{i\tau s}ds=\widehat G(0)^2-\int c_G(s)(1-\cos\tau s)\,ds$ ($c_G$ is even and $\int c_G=\widehat G(0)^2$)
and $|\widehat{\Delta_sF}|^2=4\sin^2(\tau s/2)|\Fh|^2=2(1-\cos\tau s)|\Fh|^2$, the identity follows from $|\widehat{F*G}|^2=|\Fh|^2|\widehat G|^2$
and Fubini. Primitivity: $\Phi_{F*G}=\PhiF\Phi_G$ and $\Phi_{\Delta_sF}(z)=\PhiF(z)\big(e^{-(z-\frac12)s/2}-e^{(z-\frac12)s/2}\big)$ vanish at $z=0,1$.
\end{proof}

The identity transports a sign from a window to a larger one only through a remainder of the wrong kind: when $c_G\ge0$
the remainder $-\frac12\int c_G\,a_\infty(\Delta_sF)$ has the sign opposite to what is being proved, and it is evaluated
in windows up to $T+S$. In the same spirit, the cone of window forms is not monotone: a form realized on $\PT$ by a
square need not extend to $\mathcal P_{T'}$, $T'>T$.

The other constructions fall into a small number of classes, and in each the obstruction is one of four. None of the
statements of this and the next two paragraphs is proved in this paper or, to our knowledge, elsewhere; they describe what
happens in the constructions examined and what we expect in general. Consider first squares of operators
$S=\sum_gM_{a_g}V_g$, sums of multiplications composed with translations. In $S^*S$ the atoms on the lines
$t-t'=\log n$ arise only from products of two translation pieces, and when they are matched exactly the remainder is the
archimedean form minus a nonnegative multiplication operator. This remainder would have to be nonnegative on $\PT$,
which is impossible for $T>\Tarch$, where $\mT$ itself has a negative direction on $\PT$; numerically
$\Tarch\in[0.58,0.60]$, and $\Tarch\le T_0(0)=0.74921$ by \cite[Proposition~7.1]{P46}. In addition the orthogonality
between coprime translations on which such squares rely fails for $2T>\log3$. Mixed squares $D+C$ with $C$
Hilbert--Schmidt meet the same obstruction, since the line atoms come only from $D^*D$; local square integrability of
the kernel of $C$ is not enough to conclude, as a Fourier extension operator has $C^*C=I$. The alternative source of
atoms is an infinite family of orthonormal type, and such a family reproduces the canonical factorization of
Proposition~\ref{w4:prop:taut}. Theta squares $\int|\ell_y(F)|^2d\mu(y)$ fail for a different reason: their kernels are
real-analytic and cannot reproduce the term $C_0\delta$, the hypersingular part of the archimedean form, or the atoms on
the lines.

The second obstruction is linearity in the weight. Plane-wave families and interpolation formulas at the points
$\sqrt n$ or $n$ are linear in the weight, hence valid for $c\Lambda$ by the remark following
Proposition~\ref{w4:prop:controls}; with an integer index measure their Gram matrix vanishes for $T>0.4812$, and an explicit
unimodular Euler-product family that reproduces $-\Lambda(N)$ exactly is a reformulation of the problem. Purely
arithmetic sums of squares of Dirichlet polynomials (built from unimodular transfers and resolvent rules) do not seem
to exist for $T\ge\frac12\log3$; we do not know whether they exist for $\frac12\log2<T<\frac12\log3$. The squares built
from $\Lambda*1=\log$ (a separator combined with Poisson summation, the exterior tail of the hyperbola, the Smith
factorization $[\log\gcd]=E\,\mathrm{diag}(\Lambda)E^{\mathsf T}$, a two-prime channel, the divisor defect) are, in
every case examined, either linear in $\Lambda$, and then equivalent to a shift of the window or valid for $c\Lambda$,
or quadratic with a remainder that is independent of the weight in the sense of Proposition~\ref{w4:prop:polar}.
Translation-invariant squares whose symbol is a function are equivalent to dual certificates and are treated in
Section~\ref{w5:sec:dual}; squares adapted to cells and lattices are taken up in Section~\ref{w4:sec:local}. Gaussian
non-homogeneous profiles with polynomially bounded atomic coefficients do not appear to give a realization either.

The third obstruction is validity in the formal semigroup. The collision identity
$\|Ea\|^2=\|a\|^2+(\text{collisions})$ for the multiplicative semigroup separates $\N$ from the Beurling system of
Table~\ref{w3:tab:falsifiers} exactly, but with the wrong sign; moreover the monoid generated by $2,\frac{16}5,5,\dots$
has no multiplicative collision below $16$, while that system already fails earlier. The separation between $\N$ and
the Beurling system must therefore come from the additive position of the atoms and not from the semigroup. In the
opposite direction, positivity of the Hessians of given classes of potentials with respect to positive measures does not
exclude an arithmetic identification, so no general prohibition of squares of this kind holds. For comparison,
positive-kernel realizations of the zero side are semidefinite and cannot see the critical line
\cite[Theorems~5.1--5.3]{P18}, and transfers of positivity of auxiliary objects in the Li--Keiper setting fail relative to
their stated inputs \cite{P37}; Bochner's theorem \cite{Bochner1932} is the common background.

These obstructions leave several constructions untouched, none of which has been carried out: a positive block that is
not a Gram matrix, has its own source of positivity, is not exponential in the weight and couples distinct primes; odd
and non-linear pairings in the near-coincidence matrix; sieve weights defined by an extremal problem on the exact values
of $\Lambda$, which would give a square non-linear in $\Lambda$; a square mixing discrete translations with a singular
continuous part, whose cross term carries the additive position of the atoms; and an explicit positive matrix at low
frequency that uses $\Fh(\pm\frac i2)=0$ arithmetically.

\phantomsection\label{w4:fam:R2}%
A related idea is to combine two objects: one with a sign that does not see $\Lambda$ (an archimedean square, a Gram
matrix of translates) and one that sees $\Lambda$ but has no sign (the prime part, a M\"obius-twisted kernel), in such a
way that the weight appears as the equality case of an inequality. Fourteen operations of combination were considered:
sum, oriented pairing, Gram matrix, Schur complement, conjugation, commutator, flow, Legendre transform, min--max,
tensor product, angle, Kre\u{\i}n-space pairing, Birkhoff ordering and convex functionals; also the weight as the
equality case of Cauchy--Schwarz, an identity quadratic in the weight in the quotient channel $m/n$, and boundary
concomitants of the product and quotient channels. Eleven of the fourteen reduce to constructions discussed elsewhere in
this paper and fail with them, and for the other three we found no object to which the operation applies.
Proposition~\ref{w4:prop:controls} applies to every combination whose sign comes only from squares independent of the
weight and from relations valid for a test weight; a combination in which the weight enters only linearly is valid for
$c\Lambda$. Transport from products $mn$ to quotients $m/n$ should be impossible for finite families because of unique
factorization. For transitions between windows other than restriction there is the following elementary constraint: if
$\Phi:\mathcal P_{T'}\to\PT$ is linear with $\mathfrak a_{T'}(F)=c\,\aT(\Phi F)+R(F)$, $c>0$, $R\ge0$, and
$\|\Phi F\|\ge s\|F\|$ on $\mathcal P_{T'}$, then $\lambda(T')\ge c\,s^2\lambda(T)$ whenever $\lambda(T)\ge0$. So
$c\,s^2\le\lambda(T')/\lambda(T)$ when $\lambda(T)>0$, and since the margin $\lambda(T')$ is small
(Section~\ref{w2:sec:margin}), $\Phi$ must be almost non-injective or the identity nearly lossless.

Two strengthenings were examined more closely. The first completes a square and asks that the remainder be nonnegative
on its own. In a divisor and logarithmic-sampling basis the window matrix splits as $B_{\mathrm{arch}}-(C+C^*)$ and one
writes, for a scalar $\alpha>0$,
$c^*(B_{\mathrm{arch}}-C-C^*)c=\|\sqrt\alpha\,u-\alpha^{-1/2}v\|^2+\rho_\alpha(c)$ with $u,v$ linear in $c$, demanding
$\rho_\alpha\ge0$. Computations indicate that already for the first basis vector $\rho_\alpha\to-\infty$ uniformly in
$\alpha$ as the sampling level $Y$ grows, the norm of $u$ growing like $\prod_{p\le Y}(1+1/p)$, although the quadratic
form itself is unaffected; and that in the matrix version, with any positive metric $H$ in place of $\alpha$, an
archimedean negative direction forces the remainder to be negative. Averaging over admissible enlargements of the
sampling modulus gives nothing, because the contact matrix is the same for each. What fails here is the separation of the archimedean part from the contacts, not the positivity of the form.

The second strengthening asks for a matrix $H$ that agrees with the cell matrix of Section~\ref{w4:sec:local} on the
primitive subspace and has nonpositive off-diagonal entries (a weighted graph Laplacian plus a nonnegative diagonal),
or is a sum of two-coordinate squares. Computations in interval arithmetic, which we do not reproduce, indicate that
scalar null corrections are incompatible at size $5$, general rank-two null corrections at size $6$, Laplacian
representations at every size $\ge6$ (an explicit dual witness extends by zeros), and two-coordinate squares with
complex coefficients at every size $\ge8$. On the other hand, on the three-dimensional subspace where all
prime-valuation moments for $p\le7$ vanish, the form at size $8$ appears to be exactly a positive combination of six
edge terms $w_{ij}|\alpha_i-\alpha_j|^2$, and a frozen extension of this local representation fails at the next
composite size. Finite certificates of this kind do not accumulate into a uniform rule; a sparse rule valid at all
sizes would have to use congruences or more coordinates.

\subsection{Local decompositions}\label{w4:sec:local}

The approaches of this subsection try to make positivity local: one splits $\aT$ into pieces indexed by primes, places,
cells of a lattice or blocks of a truncation, proves a sign for each piece, and adds. The common outcome is that
positivity, where it holds, is collective over the whole initial segment of primes. Local pieces lose their sign near
the archimedean threshold $\Tarch\in[0.58,0.60]$ (a numerical value) or in the range $0.70$--$0.75$, around the value
$T_0(0)=0.74921$ of \cite[Proposition~7.1]{P46}, beyond which the count of negative directions of the archimedean form
already forces a contribution of the primes.

\phantomsection\label{w4:fam:R3}%
\emph{Primes and places.} The first idea is to distribute the archimedean form among the prime contacts,
$\mT=\sum_{p}\mathfrak m_p+\mathfrak m_{\mathrm{rest}}$ with $\mathfrak m_p\ge0$, so that each
$q_p(F)=\mathfrak m_p(F)-\sum_{k:\,p^k\in\NT}v_{p^k}\Rel\cF(k\log p)$ is nonnegative on $\PT$. Variants use potentials
bounding each contact, coupled distributions of two primes, distributions with atoms split by position, a local form
$G_p$ for each Euler factor tested for reflection positivity, a Hasse--Minkowski factorization
$Q=Q_\infty\times\prod_pQ_p$, and the semilocal forms of \cite{P43,P45} in which only the places of a finite set $S$ are
kept.

Positivity does not localize in this way. Numerically, each local kernel computed ($p=2,3,5,7$) is indefinite,
positive at frequency $0$ and negative at $\pi/\log p$. A two-position distribution of the archimedean form among the
primes $2$ and $3$ attains the optimal value $\lam$ of its class up to $T\approx0.697$ and fails after it, the optimal
diagonal pair ceasing to be optimal at a window in $(0.696,0.699)$; a distribution with atoms split by position reaches
$0.75$. Potentials bounding each contact separately appear to fail before $\Tarch$. On windows in $(0.59,\log2]$, where
the contacts are $n=2,3$, the sign is obtained across the two primes and not prime by prime: computations indicate that
the minimal rank of a transfer between the two pieces needed to restore positivity equals $n_-(q_2)+n_-(q_3)$, and we
know no closed form for the transfer vector. The proposed law that the number of missing squares equals the number of
contacts in the window is false in this class, numerically. The reason behind these failures is that the negative
directions of the archimedean form on $\PT$, of which there is at least one beyond $\Tarch$, must be compensated by the
joint action of several contacts; and the prime part is a distribution, so an envelope bounding it pointwise does not
exist.

For the semilocal forms, with any finite set $S$ of places, one expects coercivity on very short windows (a lower bound
by $\frac3{10}$ of a weighted norm for $2T\le\frac1{150}$; compare the bound of \cite[Theorem~10.9]{P46} for the
archimedean form on short windows) and negativity at low frequency on long windows, for every finite $S$ including
$S=\emptyset$. Model forms with two or $n$ primes and an archimedean term, in which the cross terms vanish exactly, are
positive; but they place the divisor on the critical line by hypothesis, which is the easy direction of Weil's
criterion, and they do not bear on the full case.  The Osterwalder--Schrader reflection \cite{OsterwalderSchrader1973} of a single Euler factor
gives an oscillatory two-point function, and reflection positivity between places appears to fail as soon as a single
atom is present. The Hasse--Minkowski factorization fails because the form on the zero side does not factor, and there
is in any case no Hasse--Minkowski principle in infinite dimension. Independently of the sign, an Euler product forces
the per-prime anatomy of a localized explicit-formula form to vanish off prime powers \cite[Theorem~3.1]{P10}; the
converse holds under a non-vanishing hypothesis on the autocorrelation of the ground state
\cite[Theorems~5.1--5.2]{P10}, and model forms of this shape have per-prime profiles that change by at most an explicit
Gaussian tail when the coefficients beyond a cutoff $P$ are changed, under a spectral-gap condition
\cite[Theorem~4.1]{P09}.

Several questions remain here. An explicit formula for the distribution between the two primes for windows between
$\Tarch$ and $0.697$ would be a proof by mechanism above the archimedean threshold, though a local one. The function
$n_{\mathrm{miss}}(T)=\min\big(n_-(q_2)+n_-(q_3)\big)$ has not been determined. A joint constraint between primes would
require a new identity, and neither a version of $G_p$ across places nor a reflection between places has been built.
Finally, an entangled Osterwalder--Schrader model of the full form, not a product over places, is not excluded: the
archimedean jump density of Proposition~\ref{w4:prop:levy} is completely monotone,
$K(u)=2\sum_{k\ge0}e^{-(2k+\frac12)u}$, hence reflection positive, since
$\iint F(t)F(s)K(t+s)\,dt\,ds=2\sum_k\big|\int F(t)e^{-(2k+\frac12)t}dt\big|^2\ge0$ for $F$ supported in $t>0$.

\phantomsection\label{w4:fam:R4}%
\emph{Cells and jump coordinates.} A second idea is to pass to the multiplicative variable and use integrality. For
$F$ on $\R$ put $f(x)=x^{-1/2}F(\log x)$, $x>0$, so that $\|f\|_{L^2(dx)}=\|F\|_2$, and for $n\ge1$ let
$D_nf(x)=\sqrt n\,f(nx)$. Then $\langle f,D_nf\rangle=\int F(t)F(t+\log n)\,dt=\cF(\log n)$ for real $F$,
$A(F)=\int f\,dx$ and $B(F)=\int f\,dx/x$. After a translation of the window, primitive functions become functions on
$[1,N+1)$ with $\int f=\int f/x=0$. Integer dilations respect the structure of unit cells $[j,j+1)$, and non-integer
ones do not:

\begin{proposition}[integer division on unit cells]\label{w4:prop:cells}
Let $\mathcal S$ be the space of step functions constant on the unit cells $[j,j+1)$, $j\ge0$, and $\mathcal H_0$ the space of functions
with zero mean on every unit cell.
\begin{enumerate}
\item For $z>0$, the map $f\mapsto f(\cdot/z)$ preserves $\mathcal S$ for every $f$ if and only if $z\in\N$; for $z=n\in\N$ it acts on
coefficients by $(B_n\alpha)_j=\alpha_{\lfloor j/n\rfloor}$, and $B_mB_n=B_{mn}$.
\item For $n\in\N$, $D_n$ maps $\mathcal H_0$ into itself; $D_{16/5}$ does not.
\end{enumerate}
\end{proposition}

\begin{proof}
(1) For $x\in[j,j+1)$ and $n\in\N$, $x/n\in[\lfloor j/n\rfloor,\lfloor j/n\rfloor+1)$; and $\lfloor\lfloor j/m\rfloor/n\rfloor=\lfloor j/(mn)\rfloor$.
If $z\notin\N$, take $f=1_{[0,1)}$: then $f(x/z)=1_{[0,z)}$ jumps at the non-integer point $z$, so it is not in $\mathcal S$.
(2) If $x\in[j,j+1)$ then $nx$ runs over $[nj,nj+n)$, a union of $n$ unit cells, so the mean of $f(n\,\cdot)$ on $[j,j+1)$ is a sum of cell means of $f$.
For $z=\frac{16}5$ and $f=1_{[3,7/2)}-1_{[7/2,4)}\in\mathcal H_0$, the mean of $f(zx)$ over $[1,2)$ is $\frac5{16}\int_{16/5}^{32/5}f=\frac5{16}\big(\tfrac3{10}-\tfrac12\big)=-\tfrac1{16}\ne0$.
\end{proof}

Write a primitive step function as $f=\sum_j\beta_j1_{[0,j)}$ in jump coordinates. Then $\|f\|^2=\beta^*G\beta$ with $G_{jk}=\min(j,k)$, and
$f$ is supported in $[1,\infty)$ and primitive if and only if
\begin{equation}\label{w4:eq:jumpprim}
\textstyle\sum_j\beta_j=0,\qquad\sum_jj\beta_j=0,\qquad\sum_j\beta_j\log j=0,
\end{equation}
because $\int f=\sum j\beta_j$ and, once $\sum\beta_j=0$, $\int f/x=\sum\beta_j\log j$.

\begin{proposition}[the prime part in jump coordinates]\label{w4:prop:jump}
For a real weight sequence $a$ put $\mathcal T_a(x)=\sum_{n\ge2}\frac{a_n}n(x-n)_+$, a convex piecewise linear function when $a\ge0$, with
$\mathcal T_a''=\sum_n\frac{a_n}n\delta_n$. If $\beta$ satisfies \eqref{w4:eq:jumpprim}, then
\[
-\sum_{n\ge2}2a_nn^{-1/2}\cF(\log n)=\beta^*(R+R^*)\beta,\qquad R_{kj}=j\,\mathcal T_a(k/j),
\]
with no dependence on a cutoff.
\end{proposition}

\begin{proof}
$\langle1_{[0,j)},D_n1_{[0,k)}\rangle=\sqrt n\min(j,k/n)$, so $2a_nn^{-1/2}\cF(\log n)=2a_n\sum_{j,k}\beta_j\beta_k\min(j,k/n)$. Write
$\min(j,k/n)=k/n-(k/n-j)_+$; the first term contributes $0$ because $\sum_j\beta_j=0$. Hence the left side equals
$2\sum_{j,k}\beta_j\beta_k\sum_n\frac{a_n}n(k-nj)_+=2\sum_{j,k}\beta_j\beta_kj\,\mathcal T_a(k/j)$, a finite sum.
\end{proof}

The archimedean part should have an equally explicit kernel in jump coordinates, in terms of $x\log x$, $\gamma$ and
$\log2\pi$, and in cell coordinates the shape $D-\frac12H-\frac12T$ with a Hankel matrix $H$ that is a Stieltjes moment
matrix, hence positive semidefinite, entering with a minus sign; we have not written out a proof. One also expects that
positivity of the cell matrices $M_N$ (archimedean minus prime part on $\alpha\in\R^N$) on the primitive subspace for
every $N$ is equivalent to the window problem; no reduction to a fixed size is known. Numerically, the Rayleigh minima on
$\sum\alpha=0$ decrease from $0.2636$ at $N=3$ to $0.0052$ at $N=40$, remain positive for $1.1\Lambda$, and become
negative for the Beurling system from $N=10$: the cell block detects a lower bound on the weight but is slow to see an
excess.

Integrality enters through rational tests. If $\beta\in\Q^L$ satisfies $\sum\beta_j\log j=0$ then, since the numbers
$\log p$ are linearly independent over $\Q$, $\sum_j\beta_jv_p(j)=0$ for every prime $p$; for example
$\beta=(-1,3,-1,-2,0,1)$ satisfies \eqref{w4:eq:jumpprim}. We expect that nonnegativity on all integer unit-cell primitive
tests of all sizes is equivalent to the window problem, with new obligations arising only at composite sizes. The
difficulty is that positivity on small carriers does not control the coupling between them. Computations in interval
arithmetic indicate that a three-cell primitive family $f_k$ has $\aT(F_k)\ge10^{-3}\|F_k\|^2$ for all $k\ge3$ and that
all planes $\spn(F_3,F_k)$ are positive, but that the same holds for $1.1\Lambda$, and that an explicit combination of
the $F_k$ is negative for $1.1\Lambda$ while positive for $\Lambda$. A gcd law for cell projections, which would
organize the coupling, is false: for the conditional expectations $E_m$ onto functions constant on cells of length
$1/m$ one has $E_2E_3\ne E_1$ (this is elementary), and the M\"obius combination $E_6-E_3-E_2+E_1$ appears to be
indefinite on primitive tests. In a sampled matrix model of the window form, the difference between the contact matrix
and the M\"obius-transformed logarithmic kernel has entries built from
$d_m(n)=\Lambda(n)-1_{(m,n)=1}\log n-\Lambda((m,n))\le0$, yet its Hermitian part appears to be indefinite on primitive
profiles; and $d\le0$ for all $m,n$ holds exactly when the weight is at most $\Lambda$, which is an upper bound only.
Dual certificates supported on a lattice with $M$ nodes do not seem to exist for $M\le8$ at $T=\frac12\log3$ nor for
$M\le24$ at windows in $(0.68,\log2]$, according to computer-assisted searches, and lattice sampling costs a loss of
order $1/M$; see Section~\ref{w5:sec:dual}.

Translations and bands behave in the same way. Spans of translates of a profile with $\aT(F)>0$ within a radius
depending on the profile are positive, and heads together with bands of moderate length are jointly positive. But two
distant tubes have normalized energy tending to $-\infty$, because the prime part collects whole bands of primes and
$\sum1/p=\infty$; majorants by absolute values between bands cost $\gg\sqrt R$ against a budget $O(\log R)$; and the
bounds on the negative index obtained in this way hold equally for $c\Lambda$. These are expectations supported by
computation, not theorems. The energies of plateaus and the normalized defect $\Gamma=\lim\beta_N/N$ of $N$-cell tests
lead back to the original inequality: one expects that the window problem holds if and only if $\Gamma=0$, which is a
reformulation. A different cell scheme bounds the archimedean form from below on each cell of a mesh in $t$ and bounds
the prime part separately by its norm. It gives positivity of $\aT$ on the subspace of functions satisfying the two
primitivity conditions on each of $k$ blocks as soon as $\log(k/T)$ exceeds a constant plus the norm of the prime form
\cite[Corollary~10.13]{P46}; since the available bound for that norm is of order $e^T$, the number of blocks grows
doubly exponentially in $T$. The cost lies in the separation of the two parts and not in the sign.

The mechanism common to these failures is the following. Integer dilations are triangular along divisibility chains for
\emph{every} weight on $\N$, so the cross terms are linear in the high part and have no sign; the convexity of
$\mathcal T_a$ and the nonnegativity of the entries of $R$ hold for every nonnegative weight, in particular for
$1.1\Lambda$; and the archimedean block is indefinite on the constraint space. Arguments that proceed term by term, layer by layer or plane by plane therefore do not succeed. The global sign of the cell carrier is the window problem itself. What would be new is an explicit square of
the cell block, for instance a closed formula for its $LDL^{\mathsf T}$ pivots that uses the integrality of
Proposition~\ref{w4:prop:cells} and explains why it does not separate $1.1\Lambda$; an explicit rule for updating the
cross terms at every composite size; or a signed control of the crossings between bands.

\phantomsection\label{w4:fam:R5}%
\emph{Truncations and regularized self-energies.} Two further approaches replace the prime part by an arithmetic
truncation that is easier to sign. In the first, a stationary identity writes $\aT(F)$ for a primitive smooth test as
the sum of an indefinite logarithmic energy
$Q(H)=\frac12\int_{|\kappa|\ge1}\log|\kappa|\,|H|^2-\frac12\int_{|\kappa|<1}|\log|\kappa||\,|H|^2$ of a Poisson-type
transform $H$ of the test and a moment term $\langle tF,F\rangle$; one truncates $H$ by a Ramanujan-sum expansion over
denominators $d\le D$ coprime to a fixed modulus and seeks lower bounds
$Q(H_D)+\langle tF,F\rangle\ge-\varepsilon_D(F)$ with $\varepsilon_D\to0$. What one expects to have, without a written
proof, is the following: decay bounds for $H_D$ and for $H_D-H$, membership of $H_D$ in the energy space, the one-sided
Fatou limit $Q(H)\ge\limsup_DQ(H_D)$, strong $L^2$ convergence $H_D\to H$ under the desired lower bounds, and a
regularization by compensated cells giving absolutely convergent expansions. The lower bounds themselves are not
known, and they cannot be obtained sample by sample: by an elementary change of variables the energy of the $\ell$-th
sample decays only like $(\log\ell)/\ell$, so the diagonal of the expansion diverges and the expansion is not absolutely
summable even at fixed $D$. Nor is the Poisson reference a limit in the energy space, since such convergence would
force $\aT(F)=0$. The indispensable cancellation lies in the cross terms between samples, and negative energy can
concentrate near $\kappa=0$ under the logarithmic weight with vanishing $L^2$ mass.

In the second approach a step test is encoded by its cyclotomic product
$\prod_j(1-z^j)^{\beta_j}=\prod_d\Phi_d(z)^{c_d}$, $c_d=\sum_{d\mid j}\beta_j$, and one studies the regularized
self-energy
\[
J_{\mathrm{sc}}(\beta)=-\sum_{i,j}\gcd(i,j)\,\beta_i\beta_j\log\frac{i+j}{\gcd(i,j)},
\]
which, like $\aT$ in the variable $x$, is multiplied by $k$ under the dilation $\beta_j\mapsto\beta_{j/k}$ (with $\beta_{j/k}=0$ for $k\nmid j$).

\begin{proposition}[self-energy]\label{w4:prop:selfenergy}
$A(\beta):=\sum_d\varphi(d)c_d^2=\sum_{i,j}\gcd(i,j)\beta_i\beta_j$, and $A(\beta)>0$ for $\beta\ne0$.
\end{proposition}

\begin{proof}
$\sum_{d\mid(i,j)}\varphi(d)=\gcd(i,j)$ gives the identity; $A(\beta)=0$ forces all $c_d=0$, hence $\beta=0$ by M\"obius inversion.
\end{proof}

For $\beta=(-1,3,-1,-2,0,1)$ a direct evaluation gives
$J_{\mathrm{sc}}(\beta)=-72\log2+36\log3+10\log5-2\log7\in(1.846,1.847)$ and $A(\beta)=14$. Nevertheless $J_{\mathrm{sc}}$
does not appear to be positive semidefinite: sums of prime dilations of this $\beta$ over $\sqrt X<p\le X$, $X=2^k$,
$k\ge4096$, preserve both moments and all valuation conditions, and computations indicate that they have
$J_{\mathrm{sc}}<0$, the negative cross terms between prime dilations exceeding the diagonal. One then expects
$\inf J_{\mathrm{sc}}/A=-\infty$ on admissible integer tests, so that neither a repair $J_{\mathrm{sc}}+cA$ with a fixed
constant nor a constant rescaling of the cutoff can work. Neither observation refutes $\aT\ge0$, nor the minorant
inequality $J_{\mathrm{sc}}\le\aT$, which we do not know how to decide. What remains to be tried is a collective
compensation between blocks by an explicit operation (grouping before integrating), and a joint compensation for
$J_{\mathrm{sc}}$ other than $cA$.

\subsection{Imported positivity}\label{w4:sec:imported}

\phantomsection\label{w4:fam:R6}%
The third idea is to import a positivity theorem from another field and transport it to $\aT$. The candidates range from positive-definite kernels and log-concavity to curvature conditions for Markov
generators, reflection positivity, correlation inequalities, KMS states and hyperbolic polynomials. About forty such
imports have been examined. A general statement for positivity imported by a construction that is uniform in the prime
part is Observation~\ref{w3:bar:B10}; here we describe the individual mechanisms. The standard references are
\cite{Schoenberg1951,Prekopa1973,BrascampLieb1976,BakryEmery1985,BeurlingDeny1959,OsterwalderSchrader1973,BostConnes1995,Garding1959}.

The starting point is Proposition~\ref{w4:prop:levy}: for every nonnegative weight $a$, $\aT[a]$ is a L\'evy (Dirichlet)
form minus the constant $\kappa_T[a]$.

\begin{observation}[Markov structure is shared by all nonnegative weights]\label{w4:obs:markov}
Every property of $\mathcal E_T[a]$ that follows from the positivity of its jump measure (Beurling--Deny contraction, positivity
preservation of the semigroup, the Markov property of the killed process) holds for every weight $a\ge0$, in particular for $c\Lambda$
with $c>1$, for which the window problem fails in suitable windows (Table~\ref{w3:tab:falsifiers}).
\end{observation}

Moreover, every nonzero real $F\in\PT$ changes sign at least twice in $\IT$ (Proposition~\ref{w2:prop:signchange}).
Perron--Frobenius arguments, which give a positive ground state for the killed L\'evy process, therefore do not apply
on $\PT$: the constraint $A=B=0$ removes the positive cone. The ground-state transform is nevertheless available as an
identity.

\begin{proposition}[ground-state identity]\label{w4:prop:groundstate}
Let $M(t,s)\ge0$ be a bounded symmetric kernel on $\IT\times\IT$, $V$ a bounded real function, and $q(F)=\int VF^2-\iint M(t,s)F(t)F(s)\,dt\,ds$
for real $F\in L^2(\IT)$. If $\varphi>0$ is bounded and bounded below on $\IT$ and $L\varphi=V\varphi-\int M(\cdot,s)\varphi(s)ds$, then
\[
q(F)=\int\frac{L\varphi}{\varphi}\,F^2+\frac12\iint M(t,s)\varphi(t)\varphi(s)\Big(\frac{F(t)}{\varphi(t)}-\frac{F(s)}{\varphi(s)}\Big)^2dt\,ds .
\]
\end{proposition}

\begin{proof}
Put $F=\varphi g$ and use $g(t)g(s)=\frac12\big(g(t)^2+g(s)^2\big)-\frac12\big(g(t)-g(s)\big)^2$ together with the symmetry of $M$.
\end{proof}

For the Weil form, Proposition~\ref{w4:prop:levy} supplies the kernel (the L\'evy density together with the atoms at
$\log n$, all nonnegative), but the archimedean density is singular at the diagonal and the identity is formal unless
regularized. More seriously, a positive supersolution $\varphi$ with $L\varphi\ge0$ would yield positivity on all of
$L^2(\IT)$, and this fails when $T$ is not small: $\aT(F)\le\mT(F)$ for $F\ge0$, because then $\cF(\log n)\ge0$, and the bottom of
$\mT$ on $L^2(\IT)$ tends to $m(0)=C_0<0$ as $T\to\infty$. So one must work with the restriction to $\PT$, where the
identity is a reformulation. Numerically, on the windows where positivity is known, the inequality obtained in this way
is an equality to between $10^{-8}$ and $10^{-15}$, which suggests that every method that proves a gap with an intrinsic
loss must fail.

\phantomsection\label{w4:tab:import}%
The individual imports fail in one of three ways. The first is that the imported positivity holds for every nonnegative
jump measure, as in Observation~\ref{w4:obs:markov}. A multiplicative locally positive-definite weight times an
archimedean kernel independent of the prime gives no sign, so the theorems of Bochner and Schoenberg do not apply as
they stand. For log-concavity, $-D^2\log Z$ of a marginal equals an expectation of a Schur complement plus a
nonnegative remainder \cite{BrascampLieb1976}, and in the application to $\aT$ the remainder does not depend on the
weight. The curvature condition $CD(0,\infty)$ holds for translation-invariant pure-jump generators with finitely many
prime-power jumps, using only $\Lambda\ge0$ \cite[Theorems~5.2--5.3]{P35}; integrated curvature conditions are
equivalent to a spectral gap without sign, and the ceiling for the gap appears to be negative already at
$T=\frac12\log3$, where positivity holds. Indeed, integrating the curvature-dimension inequality reproduces the
Poincar\'e inequality of Proposition~\ref{w4:prop:levy}, so curvature methods can at best restate the problem. Cheeger's
bound gives $\frac12$ where $1-O(1/M)$ is needed, and the discrete curvatures of Ollivier and of Erbar and Maas reduce
to the Bakry--\'Emery condition. Reflection positivity fails for a single Euler factor and between places, as described
in Section~\ref{w4:sec:local}, and chessboard estimates force a constant at two primes. For the correlation inequalities
(GKS, FKG, Ginibre), the equality class of FKG is the class of Euler-product vectors, and no $\pm1$ gauge makes the
near-coincidence matrix ferromagnetic, the cycle $m\to\dots\to2m$ being closed by the multiplicative edge $(m,2m)$.

The second way is that the positivity is exponential in the weight of each prime. Closures built by Schur products,
Pick matrices, total positivity or products of Euler factors from factors exponential in the weight of each prime
remain positive for $2\Lambda$, after deletion of a prime and for Beurling weights. The third is that the import
reproduces $\zeta$ itself: the partition function of the Bost--Connes system is $\zeta(\beta)$ \cite{BostConnes1995},
and a KMS argument repeats the definition. In the same way, hyperbolic polynomials, Lorentzian signatures and
Hodge--Riemann relations lead, in the forms tried, to a restatement of the problem, a $2\times2$ polar reformulation in
the Lorentzian case; see \cite{P14} for the Hodge-index and hyperbolicity reformulations.

Some imports fail more simply. Positive-kernel reproducing spaces give semidefinite forms, and smooth families do not
detect the critical line \cite[Theorems~5.1--5.3]{P18}. For Doob and Markov representations we found no representation
by a positive family of measures and no positive Laplace measure on a ray for the relevant transform, and an additive
positive comparator of size $o(1/u)$ must vanish; a list of obstructions in the Doob coordinate is given in \cite{P41},
and in the Li--Keiper setting the polar factor of a completed cocycle is not the Laplace transform of a positive measure
\cite[Proposition~3.6]{P37}. Real-positive comparators in a ladder of powers are excluded relative to three classical
theorems. Sieves, gcd and Smith matrices and the Nyman--Beurling--B\'aez-Duarte criterion are either free of $\Lambda$
or repeat the sieve arguments of Section~\ref{w7:sec:primeside}.

Two facts specific to the Weil form belong here. Adding primes raises the background: in the language of
Proposition~\ref{w4:prop:levy}, adding a contact increases $\kappa_T$ by $v_n$ and the Dirichlet form by at most
$2v_n\|F\|^2$. This is the only fact specific to $\zeta$ found among these imports, and it acts as an obstruction. And
the scaled weight $(1+c)\Lambda$ still defines a Markov form for every $T$, while its bottom is expected to tend to
$-\infty$. None of the imports uses the interplay of addition and multiplication of integers at near-coincidences
$m\approx kn$, which is where the window form differs from its test weights.

The following have not been tried and are not covered by the reasons above: a convex functional whose convexity comes
from a source other than a positive measure; a potential that is convex only for the true weight; an entangled
reflection-positive model of the full form, as at the end of the discussion of primes and places in
Section~\ref{w4:sec:local}; a Poincar\'e or Wirtinger inequality on $\PT$ with an explicit constant, to be compared first
with the weighted Wirtinger inequalities already excluded; elements built from the additive shift in an algebra with an
exact mixed sum--product relation; and a pointwise density $\tilde\Gamma\ge0$, not arising from nonnegative rates, whose
integral is $\aT$ plus a term independent of the weight.

\section{Duality, Schur complements and kernel positivity}\label{w5:sec:dual}

The approaches of this section work directly with the form $\aT$ on $\PT$ and use classical tools: a lower bound by
duality on the frequency side, the reduction to a finite Schur complement with a count of negative directions, and
the replacement of the family of windows by one kernel or one scale of explicit test spaces. Each gives exact identities,
and the first gives proofs of positivity on single windows; none gives the sign for all $T$.

\subsection{Dual certificates and linear programming}\label{w5:fam:R7}

The oldest way to bound a quadratic form from below on a subspace is by duality: one exhibits a function on the
frequency side that agrees with the symbol $\sT$ on everything the form can see and is bounded below pointwise. The Weil
form on a window sees only the autocorrelation $\cF$ on $(-2T,2T)$ and the two polar conditions.

\begin{definition}[certificate]\label{w5:def:cert}
Fix $T>0$. A \emph{certificate} for the window $\IT$ is a triple $(\nu,\beta_0,\beta_1)$ where $\nu$ is a finite real even
Borel measure on $\R$ with $\supp\nu\subset\{|u|\ge2T\}$ (an \emph{exterior measure}) and $\beta_0,\beta_1\in\R$. With
$G_0(u)=\cosh(u/2)\mathbf 1_{(-2T,2T)}(u)$ and $G_1(u)=u\sinh(u/2)\mathbf 1_{(-2T,2T)}(u)$ (the polar functions $g_0,g_1$ introduced after
Definition~\ref{w1:def:primitive}, extended evenly and cut off to $(-2T,2T)$), its \emph{value function} is
\begin{equation}\label{w5:eq:valuefn}
  f(\tau)=\sT(\tau)+\widehat\nu(\tau)+\beta_0\widehat{G_0}(\tau)+\beta_1\widehat{G_1}(\tau),\qquad
  \widehat\nu(\tau)=\int e^{i\tau u}\,d\nu(u),
\end{equation}
and its \emph{value} is $\inf_{\tau\in\R}f(\tau)\in[-\infty,\infty)$.
\end{definition}

Since $\nu$, $G_0$, $G_1$ are even and real, $\widehat\nu$, $\widehat{G_0}$, $\widehat{G_1}$ are real, even and bounded;
$f$ is a real even function with $f(\tau)=\log|\tau|+O(1)$ at infinity.

\begin{proposition}[weak duality]\label{w5:prop:weakdual}
If $(\nu,\beta_0,\beta_1)$ is a certificate with $f\ge c$ on $\R$, then $\aT(F)\ge c\|F\|_2^2$ for every $F\in\PT$; that is,
$\lam\ge c$.
\end{proposition}

\begin{proof}
Let $F\in\PT$. By Plancherel, $\frac1{2\pi}\int|\Fh|^2=\|F\|_2^2$, and $\aT(F)=\frac1{2\pi}\int\sT|\Fh|^2$ by
\eqref{w1:eq:trunc}. It suffices to show that the three added terms of \eqref{w5:eq:valuefn} integrate to zero against
$|\Fh|^2=\widehat{\cF}$; then $\aT(F)=\frac1{2\pi}\int f|\Fh|^2\ge c\|F\|_2^2$.

\emph{Exterior term.} The function $\widehat{\cF}=|\Fh|^2$ is integrable and $\cF$ is continuous, so
$\cF(u)=\frac1{2\pi}\int\widehat{\cF}(\tau)e^{-i\tau u}d\tau$ for every $u$. By Fubini (the measure is finite and
$\widehat{\cF}\in L^1$),
\[
  \int\widehat\nu(\tau)\,|\Fh(\tau)|^2\,d\tau=\int\Bigl(\int\widehat{\cF}(\tau)e^{i\tau u}d\tau\Bigr)d\nu(u)
  =2\pi\int\cF(-u)\,d\nu(u)=0,
\]
because $\cF$ vanishes on $|u|\ge2T$ and $\nu$ is carried by that set.

\emph{Polar terms.} For a bounded function $G$ supported in $[-2T,2T]$ the same computation gives
$\frac1{2\pi}\int\widehat G\,|\Fh|^2=\int G(u)\cF(-u)\,du$. Now
$\int\cF(u)e^{su}\,du=\PhiF(\tfrac12+s)\overline{\PhiF(\tfrac12-\bar s)}$ (Section~\ref{w1:sec:setting}); for primitive $F$
both factors vanish at $s=\pm\frac12$, so this entire function of $s$ vanishes to second order at $s=\pm\frac12$. Hence
$\int\cF(u)e^{\pm u/2}du=0$ and $\int\cF(u)\,u\,e^{\pm u/2}du=0$, and therefore
$\int\cF(-u)G_0(u)du=\int\cF(-u)G_1(u)du=0$ (both $G_j$ are combinations of $e^{\pm u/2}$, $ue^{\pm u/2}$ on
$(-2T,2T)$, where $\cF$ lives).

Finally $F\in\DT$ and $|f-\sT|$ is bounded, so all integrals converge absolutely.
\end{proof}

Three remarks on the proof. It uses nothing about $\sT$: the same inequality holds for $\aT[a]$ and any weight sequence
$a$, so a certificate carries no arithmetic information of its own, and the lower bound has to be supplied by the choice
of $(\nu,\beta_0,\beta_1)$. Next, primitivity enters only through the two identities $\int\cF(-u)G_j(u)\,du=0$, $j=0,1$.
Writing $\lambda^{\mathrm P}(T)$ for the infimum of $\aT(F)$ over the unit vectors $F\in\DT$ that satisfy these two
identities, the proof therefore gives $\lambda^{\mathrm P}(T)\ge c$; and
\[
  \lamW\le\lambda^{\mathrm P}(T)\le\lam ,
\]
the second inequality because primitive functions satisfy the identities, the first because
$\int\cF(-u)G_0(u)\,du=\Rel\bigl(A(F)\overline{B(F)}\bigr)$ for every $F\in\DT$, so that $\WT=\aT$ on the larger class.
Finally the same identity gives $\WT(F)=\frac1{2\pi}\int(\sT+2\widehat{G_0})|\Fh|^2$ on $\DT$: for the complete form the
polar coefficients are not free, and a certificate with $(\beta_0,\beta_1)=(2,0)$ and $f\ge c$ gives $\lamW\ge c$.

One is thus led to the dual problem
\begin{equation}\label{w5:eq:dual}
  \lambda^{\ast}(T)=\sup\Bigl\{\inf_\tau f(\tau):\ (\nu,\beta_0,\beta_1)\ \text{a certificate}\Bigr\}\ \le\ \lam ,
\end{equation}
which one can solve for each window by linear programming (with $\nu$ discretized on a lattice of exterior
frequencies, or in a family of profiles), certifying the resulting value function by interval arithmetic; one then
looks for a certificate whose shape is uniform in $T$. The structure of the dual problem is known.

\begin{theorem}[strong duality and non-attainment {\cite[Theorems~4.4 and~4.6]{P46}}]\label{w5:obs:strong}
For every $T>0$:
\begin{enumerate}[label=(\roman*)]
\item $\lambda^\ast(T)=\lambda^{\mathrm P}(T)$, and the supremum in \eqref{w5:eq:dual} is approached by certificates
whose exterior measure has finitely many atoms;
\item the supremum is not attained: every certificate has value strictly less than $\lambda^{\mathrm P}(T)$.
\end{enumerate}
\end{theorem}

The proof of (ii) shows why no construction by complementary slackness is available. The primal infimum is attained
at some $F_\star$, and $\widehat{F_\star}$ is a nonzero entire function, which vanishes on no interval; a certificate
of value $\lambda^{\mathrm P}(T)$ would have $f$ constant on the support of $|\widehat{F_\star}|^2$, hence on $\R$, which
is impossible. In the linear programming bounds of Cohn--Elkies and Viazovska \cite{CohnElkies2003,Viazovska2017} the
extremal function vanishes exactly on the support of the dual object; here there is no such ``magic function''.

Thus the dual formulation is an exact restatement: since $\lamW\le\lambda^{\mathrm P}(T)\le\lam$ and the two window
problems are equivalent (Section~\ref{w1:sec:setting}), $\lam\ge0$ holds for all $T$ if and only if for every $T$ and
$\varepsilon>0$ there is a certificate of value at least $-\varepsilon$.

A special class of certificates comes from \emph{flat-top weights}: even functions $w\in C_c^\infty(\R)$ with $w\equiv1$
on $[-2T,2T]$. Multiplying by $w$ the distribution on the $u$-side whose Fourier transform is the zero counting measure
leaves the truncated form on $|u|<2T$ and produces on $|u|\ge2T$ the finite exterior measure
\[
  d\nu_w=\Bigl[\tfrac12\bigl(1-w(u)\bigr)K(|u|)+2w(u)\cosh\tfrac u2\Bigr]\mathbf 1_{|u|\ge2T}\,du
  -\sum_{\log n\ge2T}\tfrac{v_n}2\,w(\log n)\,\bigl(\delta_{\log n}+\delta_{-\log n}\bigr),
\]
with $K$ the L\'evy density of Proposition~\ref{w2:prop:levy}. We put $f_w=\sT+\widehat{\nu_w}+2\widehat{G_0}$, the value
function of the certificate $(\nu_w,2,0)$ for the complete form.

\begin{theorem}[flat-top certificates and the explicit formula {\cite[Theorems~6.2, 6.4 and~8.2]{P46}}]\label{w5:obs:flattop}
Let $w$ be a flat-top weight for the window $\IT$ and $\widehat w(t)=\int w(u)e^{itu}\,du$. Then
\[
  f_w(\tau)=\sum_\rho\widehat w(\tau-\gamma_\rho),\qquad\rho=\tfrac12+i\gamma_\rho ,
\]
the sum converging absolutely. If $T\ge\frac12\log2$, then $\sup_w\inf_\tau f_w(\tau)=\lamW$, the supremum being taken over
the flat-top weights of the window, and there is a flat-top weight with $f_w\ge0$ on $\R$ if and only if $\lamW>0$.
\end{theorem}

The first assertion is the explicit formula \eqref{w1:eq:explicit} read as an identity of distributions:
multiplication by $w$ on the $u$-side is convolution by $\widehat w$ on the zero side. Thus the existence, for every $T$
and $\varepsilon>0$, of a flat-top weight with $f_w\ge-\varepsilon$ restates the complete window problem as a one-sided
smoothing property of the zero counting measure. Such a property can never be checked zero by zero.

\begin{proposition}[no positive-definite flat top]\label{w5:prop:noPD}
Let $w$ be continuous on $\R$ with $w(u)\to0$ as $|u|\to\infty$, $w(0)=1$, and suppose $w(u_0)=1$ for some $u_0\ne0$
(in particular, if $w$ is flat on a neighbourhood of $[-2T,2T]$, $T>0$). Then $w$ is not positive definite; in particular its
Fourier transform (a tempered distribution) is not a nonnegative measure, and so is not a nonnegative function.
\end{proposition}

\begin{proof}
Suppose $w$ is positive definite. By Bochner's theorem $w(u)=\int e^{iu\xi}d\mu(\xi)$ with $\mu\ge0$ of mass $w(0)=1$.
For all $u$, by Cauchy--Schwarz,
\[
  |w(u+u_0)-w(u)|^2=\Bigl|\int e^{iu\xi}(e^{iu_0\xi}-1)\,d\mu\Bigr|^2\le\int|e^{iu_0\xi}-1|^2d\mu=2\bigl(1-\Rel w(u_0)\bigr)=0 .
\]
So $w$ is periodic with period $u_0$; being periodic and tending to $0$ at infinity, it vanishes identically, contradicting $w(0)=1$. If
$\widehat w$ were a nonnegative measure, $w$ would be positive definite by the Bochner--Schwartz theorem.
\end{proof}

Consequently the value function $f_w=\sum_\rho\widehat w(\cdot-\gamma_\rho)$ of a flat-top certificate always contains
terms of both signs, even if every $\gamma_\rho$ is real: nonnegativity of $f_w$ is never termwise. Nor can it be
obtained from finitely many zeros. If $Z$ is a nonempty finite symmetric multiset of real numbers and $w$ is even, real,
continuous and integrable with $w=1$ on $[-2T,2T]$, then $\sum_{\gamma\in Z}\widehat w(\cdot-\gamma)$ takes negative
values \cite[Theorem~6.5]{P46}; the certificate is intrinsically a statement about infinitely many zeros
(compare Section~\ref{w3:sec:barriers}).

\begin{remark}[mass and value]\label{w5:obs:massvalue}
Computations with certificates of a fixed shape indicate that coming within $\delta$ of the best attainable value
requires exterior total variation of order $\log(1/\delta)$, that this logarithmic exchange is sharp, and that adding
positive exterior measure alone never raises the value. We have no proof of these statements.
\end{remark}

Restricting $\nu$ to atoms on a lattice $\{\pm(2T+jd)\}_{0\le j<M}$ turns \eqref{w5:eq:dual} into a finite linear
program. Its optimal value is an upper bound for what certificates of that class can prove, and it can be far below
$\lam$.

\begin{remark}[lattice classes]\label{w5:obs:ceil}
The following is numerical, obtained with interval enclosures in the cases where a number is given, and is not proved
in any paper we can cite. (i) The optimal value over lattice certificates with $M\le8$ atoms at $T=\frac12\log3$ is at
most $-5.6974\cdot10^{-3}$, and the optimal values with $M\le24$ atoms at $T=0.68$ and $T=\log 2$ are negative; so these
small lattice classes contain no certificate of nonnegative value, although $\lam>0$ there. (ii) We expect the optimal
values to converge to the dual value as $M\to\infty$, at rate $O(1/\log M)$ in general and $O(1/M)$ under an additional
hypothesis. (iii) Sampling the exterior measure on a lattice of mesh $d$ appears to cost a loss of exact order
$d\cosh T$ (the constant observed is $4/\pi$), caused by the jump of the polar functions at $|u|=2T$. In particular the
lattice loss is not negligible.
\end{remark}

The certificates obtained from sums of squares of Dirichlet polynomials (positivity of $\sT$ up to an exterior
correction, built from $|\sum_na_nn^{-i\tau}|^2$) are discussed with the other square representations in
Section~\ref{w4:sec:squares}. In the dual language, a comparison of the bottom of the archimedean form on $\PT$ (at most
$0.33$ at $\frac12\log3$) with the total prime mass seen by the window (at least $0.980$) indicates that such
certificates do not exist for $T\ge\frac12\log3$; we do not know whether they exist for
$T\in(\frac12\log2,\frac12\log3)$.

On single windows the method works. Certificates found by linear programming and then certified support several of the
positivity statements of Table~\ref{w2:tab:certified}; at $T=\frac12\log3$, for instance, a certificate with $437$ atoms
has value $8.8577\cdot10^{-5}$ on the primitive form (computer-assisted, and not reproduced here; proved bounds at this window
are \cite[Theorem~9.18]{P42} and \cite[Theorem~9.1]{P46}). The search for a shape uniform in $T$ gave only negative
results. The following values come from floating-point computations, uncertified, for which no program accompanies this
paper. A flat-top program using the first $1000$ zeros produces a function that is at least $1.37\cdot10^{-3}$ on
$[0,16.75]$ and takes the value $-7.321$ at $\tau=17.436$; it is not a certificate, and shows the finite-zero
obstruction numerically. Bounded-profile weights $0\le w\le1$, weights vanishing on the comb, and two-frequency ladders
interpolating from $T=0.58$ to $T=\log2$ (best value $-3.18\cdot10^{-3}$ with $24$ atoms) all give negative values. A
zero-side program with one shape for all windows is negative in each of $103$ windows tested. A certificate with polar
coefficient fixed in advance exists only with exterior measures of strongly mixed sign (ratio of total variation to
total mass about $52$). Finally, heuristically, certificates of the form ``constant floor plus tail'' need a dimension
at least of the order of $(2/\pi)T\tau_+(T)$, where $\tau_+(T)$ is the frequency beyond which the symbol is positive,
and we expect their cost to grow superpolynomially in $T$.

As a method for all $T$ the dual approach does not give the sign, for three reasons. First, by
Theorems~\ref{w5:obs:strong} and~\ref{w5:obs:flattop} the dual problem is the window problem written on the zero side: a
certificate for all windows is the explicit formula smoothed by a flat-top weight, and its nonnegativity is the
conclusion (see Section~\ref{w8:sec:open}). Second, this nonnegativity is not termwise and does not come from finitely
many zeros, so no finite verification of zeros produces it. Third, there is no induction in $T$: the map $T\mapsto\lam$
is nonincreasing (Proposition~\ref{w1:prop:elementary}), so a certificate at $T$ says nothing at $T'>T$, and
$\lam\le4\cdot10^9\,T\,e^{-15.449\,T}$ for $T\ge1$ \cite[Corollary~9.11]{P46} (Section~\ref{w2:sec:margin}), so that
the certificates become exponentially fine and the lattice and profile classes must grow without bound
(Remark~\ref{w5:obs:ceil}). Moreover weak duality is blind to the weights: the arithmetic enters only through the
inequality $f\ge c$, which is what has to be proved.

Two questions remain. Only upper bounds for what a class of certificates can prove are known; a lower bound for the
best value of lattice certificates of mesh $d$ and frequency range $U$, tending to $\lam$ in a controlled way, would
turn the linear program into a proof scheme for a range of windows. The second question is whether an exterior measure
or a flat-top weight of nonnegative value can be built from the primes, in particular from the prime powers outside the
window; the unknown is then an arithmetic object and not a smoothing of the zeros. The two constructions of this kind
known to us fail, and we do not know a general obstruction.

For the comparison of reformulations of band-limited Weil forms see \cite{P12}, and for the Toeplitz structure of window
forms \cite{P26}. Extremal functions of Beurling--Selberg type, the classical source of one-sided band-limited
approximation, are surveyed in \cite{Vaaler1985}; Hilbert-type inequalities for the off-diagonal terms are in
\cite{MontgomeryVaughan1974}. The linear programming method for extremal problems with a dual measure is that of
\cite{CohnElkies2003}, and its exact solution in dimension $8$ is \cite{Viazovska2017}.

\subsection{Schur complements, inertia, determinants and Wiener--Hopf operators}\label{w5:sec:schur}\label{w5:fam:R8}

A second group of approaches treats $\aT$ on $\PT$ as a bounded perturbation of the archimedean form by the prime form, a
finite sum of translation pairings, and tries to decide the sign by finite linear algebra: reduce to a finite Schur
complement, count negative directions, compute a determinant, or use the asymptotic theory of truncated Toeplitz and
Wiener--Hopf operators. We use two classical tools. For a Hermitian form $q$ on a Hilbert space write $n_\pm(q)$,
$n_0(q)$ for its inertia.

\begin{proposition}[inertia additivity]\label{w5:prop:haynsworth}
Let $q$ be a lower semibounded closed Hermitian form with domain $\mathcal D\subset H$, let $V\subset\mathcal D$ be
finite-dimensional, and suppose the operator of $q|_{V^\perp\cap\mathcal D}$ has $0$ in its resolvent set. Then
$q(v+p)=q(p-p(v))+S(v)$ ($v\in V$, $p\in V^\perp\cap\mathcal D$) for a unique bounded linear map $v\mapsto p(v)$ and a
Hermitian form $S$ on $V$ (the Schur complement), and $n_\pm(q)=n_\pm(q|_{V^\perp})+n_\pm(S)$, $n_0(q)=n_0(S)$; if
$q|_{V^\perp}$ is coercive, $q\succ0$ if and only if $S\succ0$.
\end{proposition}

For matrices this is Haynsworth's theorem \cite{Haynsworth1968} (the complement block need only be invertible, not
positive). The form version follows from the completion of squares used in the proof of
Proposition~\ref{w1:prop:polar}, with Lax--Milgram replaced by the invertibility of the operator of
$q|_{V^\perp\cap\mathcal D}$ from the form domain onto its dual, which holds because $0$ is in its resolvent set; the
decomposition $\mathcal D=\{v+p(v)\}\oplus(V^\perp\cap\mathcal D)$ is then $q$-orthogonal and the inertia is additive.
Applied with $V$ the polar plane, the proposition is Proposition~\ref{w1:prop:polar}.

Four ideas have been pursued: to reduce to a finite Schur complement $S_{T,k}$ on a space $V_k$ (piecewise
exponentials, Dirichlet sines, archimedean eigenvectors) off which $\aT$ is coercive, and to prove $S_{T,k}\succeq0$ by
a structural argument; to count the negative directions of the archimedean form, all of which the prime form must
cancel; to encode positivity in a determinant (Fredholm, Birman--Schwinger, Feshbach) or a recursion (Schur and
Verblunsky parameters, Jacobi matrices, canonical systems); and to use the Szeg\H{o}, Widom and Landau--Widom
asymptotics of truncated Toeplitz and Wiener--Hopf operators.

\emph{The finite reduction.} For every $T$ there is an explicit $k_0(T)$ such that for $k\ge k_0(T)$ the form $\aT$ is
coercive on the functions of $\PT$ whose two Tate moments vanish on each of $k$ consecutive blocks of $\IT$; the
complement $V_k$ is spanned by piecewise exponentials, has dimension $2k-2$, and
\[
  \lam>0\iff S_{T,k}\succ0,\qquad \lam\ge0\iff S_{T,k}\succeq0
\]
\cite[Theorem~10.14]{P46} (Section~\ref{w2:sec:reductions}). When $S_{T,k}$ is nonsingular,
Proposition~\ref{w5:prop:haynsworth} splits the inertia of the complete form as
$n_-(\WT)=n_-(\aT|_{\PT\ominus V_k})+n_-(S_{T,k})+n_-(\mathfrak G_T)$, in which the first term vanishes
\cite[Remark~10.15]{P46}. The rows of $S_{T,k}$ are indexed by the blocks of the discretization, not by the contacts
$n\in\NT$, and refining the discretization does not change the negative index: $n_-(S_{T,k})=n_-(\aT|_{\PT})$ for every
admissible $k$.

\emph{Counting negative directions.} The archimedean form has $n_-(\mT|_{\PT})\ge N(T)-2$ with
$N(T)=(2t_0/\pi)T+O(1)\approx4.004\,T$, the lower bound being split evenly between the parities
(Theorem~\ref{w2:thm:archcount}; $t_0\approx6.2898$ is the positive zero of $m$), and we expect the matching upper bound
$(2t_0/\pi)T+O(\log T)$ (Section~\ref{w2:sec:structure}). A proof of $\lam\ge0$ must therefore show that the prime form
cancels a negative space of dimension growing linearly in $T$; primitivity removes only two directions
(Proposition~\ref{w2:prop:interlace}). On small windows $n_-(\mT|_{\PT})=0$ for $T<\Tarch$, and the computations
reported in Section~\ref{w2:sec:structure} indicate that it equals $1$ up to an odd threshold in $[\log2,0.81]$ and $2$
up to $T=0.90$, with $\Tarch\in[0.58,0.60]$ numerically. For the truncated form itself the count is governed by the
polar term.

\begin{proposition}[even sector and one negative direction]\label{w5:prop:evenone}
Let $\aT^{\mathrm{ev}}$ and $\WT^{\mathrm{ev}}$ be the restrictions of $\aT$ and $\WT$ to the even functions of $\DT$. If
$\WT^{\mathrm{ev}}\ge0$, then $n_-(\aT^{\mathrm{ev}})\le1$.
\end{proposition}

\begin{proof}
For even $F$, $A(F)=B(F)$, so $\WT(F)=\aT(F)+2|A(F)|^2$. On the hyperplane $\{A=0\}$ of the even subspace,
$\aT=\WT\ge0$; a form nonnegative on a subspace of codimension one has at most one negative direction.
\end{proof}

The same argument applies on all of $\DT$: the polar term vanishes on $\{A=0\}$, so $\WT\ge0$ implies $n_-(\aT)\le1$.
Computations with the function $\cosh(t/2)\mathbf 1_{\IT}$ indicate that $n_-(\aT^{\mathrm{ev}})\ge1$ for $T\ge0.12$, in
which case $n_-(\aT)=1$ on every such window with $\WT\ge0$. On very small windows the index is zero: for $T\le0.01$ one
has $\aT=\mT\ge(C_0+\int_{2T}^\infty K)\|F\|_2^2>0$ by Proposition~\ref{w2:prop:levy}, so $n_-(\aT)=0$, while $\WT\succ0$
because the polar term is at most $4\sinh T\,\|F\|_2^2$ in absolute value. We expect, but have not proved, the converse
of Proposition~\ref{w5:prop:evenone}, namely that $n_-(\aT^{\mathrm{ev}})\le1$ for all $T$ implies the even part of the
complete window problem. The criterion would then be equivalent to the even complete window problem; it moves the
difficulty and does not reduce it. On the zero side, a pair of zeros $\rho$, $1-\bar\rho$ off the critical line
contributes to \eqref{w1:eq:explicit} the term $2\Rel\bigl(\PhiF(\rho)\overline{\PhiF(1-\bar\rho)}\bigr)$, a hyperbolic
plane in the zero-side realization of the form; it changes the negative index of that realization but not its signature,
so that only the absolute index can detect it.
Numerically, and granted the verified zeros of $\zeta$ up to a large height, the negative index of the even block of
the first $\approx4T$ Dirichlet sines of the complete form is at most one at the points of a grid in $[0.75,4.245]$; we
do not know how to pass from the grid to the intervals between its points, and the computation stops near $T=4.25$ by
loss of margin.

\emph{Determinants.} Every determinant that was proposed as a criterion (of a finite Galerkin block, of a Schur
complement, of a Fredholm or Feshbach reduction) was found to factor as $\lam\cdot R$ with $R$ of fixed sign in its
regime (Section~\ref{w3:sec:barriers}), so that it contains no information beyond $\lam$. The Birman--Schwinger form of
the problem shows the mechanism. Fix $c>\max(-C_0,\sum_{n\in\NT}v_n)$, so that $\mT+c\|\cdot\|_2^2$ and
$\pT+c\|\cdot\|_2^2$ are positive, and let
\[
  \beta_{\max}=\sup_{0\ne F\in\PT}\frac{\pT(F)+c\|F\|_2^2}{\mT(F)+c\|F\|_2^2},
\]
the largest Birman--Schwinger eigenvalue of the shifted prime part relative to the shifted archimedean part
\cite{ReedSimon4}. Since $\aT(F)=(1-\beta(F))(\mT(F)+c\|F\|_2^2)$, where $\beta(F)$ is the quotient in the display, one
has $\lam>0$ if and only if $\beta_{\max}<1$ (for the direct implication use the coercivity of $\aT$ on $\PT$, as in
the proof of Proposition~\ref{w1:prop:polar}), and for $\beta_{\max}\le1$
\[
  (c+C_0)\,(1-\beta_{\max})\ \le\ \lam\ \le\ \frac{c+\sum_{n\in\NT}v_n}{\beta_{\max}}\,(1-\beta_{\max}) :
\]
the lower bound follows from $\mT\ge C_0\|\cdot\|_2^2$, and the upper bound from
$\aT(F)=(1-\beta(F))\beta(F)^{-1}(\pT(F)+c\|F\|_2^2)$ applied to a maximizing sequence. So the criterion is exact, and
bounding $\beta_{\max}$ away from $1$ is the same problem as bounding $\lam$ away from $0$.

\emph{Nonnegative form plus bounded remainder.} In the decompositions of $\aT$ as a manifestly nonnegative form plus a
bounded convolution remainder that were tried (the full Sonin decomposition of Section~\ref{w6:sec:prolate}), the
remainder appears to have infinitely many negative directions on $\PT$ for every $T>\frac12\log2$, as discussed there;
such a decomposition does not exhibit the sign.

\emph{Toeplitz and Wiener--Hopf asymptotics.} On Dirichlet sines the prime form is a Toeplitz plus a Hankel matrix
(Proposition~\ref{w2:prop:combsine}). The classical asymptotic theory does not apply, for a simple reason.

\begin{proposition}[the symbol is never positive]\label{w5:prop:symbolneg}
For every $T>0$, $\sT(0)=C_0-\sum_{n\in\NT}v_n\le C_0<0$, and $\sT(\tau)<0$ on a neighbourhood of $0$. Thus the symbol $\sT$ is not
positive on any window; where $\aT$ is positive on $\PT$, a positive symbol for it can only be obtained by modifying $\sT$
by terms invisible to the window (for instance the transform of an exterior measure and the polar terms, as in Proposition~\ref{w5:prop:weakdual}).
\end{proposition}

\begin{proof}
$v_n\ge0$ and $C_0=m(0)\approx-5.37$; continuity of $\sT$.
\end{proof}

The strong Szeg\H{o} theorem and the Widom and Landau--Widom asymptotics \cite{LandauWidom1980,BottcherSilbermann2006}
describe truncations of positive, or at least elliptic, symbols at sampling rates where the window is long compared
with the scale of the symbol; compare the time--frequency limiting operators of \cite{SlepianPollak1961}. Here the
natural symbol $\sT$ is negative near $0$, a positive representative exists only after an exterior correction, and the
problem is at critical sampling. For every $\lambda>\sT(0)$ the symbol $\sT-\lambda$ changes sign, so the truncated
Wiener--Hopf operators attached to it are not elliptic and no Wiener--Hopf index or factorization is available
\cite[Theorem~4.7]{P46}. Numerically, the set where the discretized symbol is negative has relative measure about
$0.4$ while the negative index of the truncated form is zero or one, so the form is not described by its symbol in the
sense of generalized locally Toeplitz sequences; and a Landau--Widom main term, used heuristically as a predictor,
gives $\beta_{\max}>1$, the wrong side.

\emph{Quantitative attempts.} Several attempts to estimate the Schur complement directly were made; the values are
numerical unless a reference is given. Uniform $4\times4$ Schur blocks: the lower constant drops to $0.013$ at $T=0.69$.
Operator-norm bounds for the coupling between adjacent blocks: the norm equals $\pi$, which is too large
\cite[Remark~10.16]{P46}. IMS localization in position certifies nothing beyond $T\approx0.822$. A proposed ``prime-free
slack'', equal to the bottom of $\mT$ on $\PT$, is negative for $T>\Tarch$, in particular for $T>\frac12\log5$
\cite[Proposition~7.1 and Lemma~7.5]{P46}. A proposed log-convexity of $T\mapsto\lam$, and a margin of size $O(1)$ in
a cone defined by the prime number theorem, are false: the relevant infimum decays roughly like $\lam^{0.7}$. Lastly,
criteria that demand a positive crossing form at a zero of $\lam$ are impossible, since $\lam$ is nonincreasing
(Proposition~\ref{w1:prop:elementary}) and the crossing form at a crossing of zero is therefore $\le0$.

In summary, inertia, determinants and spectral invariants are data of the window: either they are blind to the
weight, the same identities holding for every weight sequence, including those of Table~\ref{w3:tab:falsifiers} with a
negative window, or they reproduce $\lam$ as a factor. Schur complements are exact, but by the two-sided comparison of
$\lambda_{\min}(S_{T,k})$ with $\lam$ in \cite[Theorem~10.14]{P46} they inherit the exponentially small size of $\lam$
(Section~\ref{w2:sec:margin}), and operator-norm and fixed-size block estimates lose constants larger than it. The
Toeplitz asymptotics need a positive or elliptic symbol and supercritical sampling, and neither holds.

Several questions are left open by this discussion.
\begin{enumerate}[label=(\alph*)]
\item Find a structure of the finite Schur complement $S_{T,k(T)}$ that does not inherit the size of $\lam$, for
instance a relative bound of the prime coupling by $\theta$ times the archimedean block inside the complement,
$\theta<1$, proving $S_{T,k(T)}\succeq0$ for $T>\frac12\log5$. Even just beyond $\frac12\log5$ positivity is known only
from the per-window certificates of Table~\ref{w2:tab:certified}, not from any structure of $S_{T,k}$.
\item Prove a two-variable index bound in the sense of \cite{P28}; such a bound is not a restatement of the window
problem and has not been attempted.
\item Bound the dimension of the low negative block from arithmetic data alone.
\item Find a computable detector of $n_-(\aT)\ge2$; none has been computed.
\item Establish a two-sided asymptotic law for $\lam$; no statement of Section~\ref{w3:sec:barriers} excludes one.
\end{enumerate}

The negative index of the Weil form as a Pontryagin index, with conditional bounds, is studied in \cite{P11}; the
operator framework of window forms and their symbols is in \cite{P26}, the tensor amplification of the negative index
in \cite{P28}, and a list of related obstructions in \cite{P41}.

\subsection{Kernel positivity and pole-restricted test spaces}\label{w5:sec:kernels}\label{w5:fam:R9}

The approaches of this subsection replace the form on $\PT$, for all $T$ at once, by a single real kernel or a single
scale of explicit test spaces, and try to prove positivity by comparison with a model, by total positivity, or by
explicit computation. Most of what is known here beyond \cite{P46} comes from computations and arguments that are not written up; we
state it as numerical evidence or as expectation, and no statement of this paper depends on it.

\emph{Temporal kernels and comparison models.} Since $\lam\ge0$ for all $T$ is a statement about all compactly supported
primitive test functions, it can be written as positive semidefiniteness of one real symmetric kernel $A(t,s)$ on a
half-line, obtained by expressing the complete form in a Laplace (theta) representation of the test functions;
equivalently, as positivity of all finite Gram matrices $\Gamma_{ij}=\WT(E_i,E_j)$ of a fixed sequence of primitive
functions of heat-polynomial type, or of all their tails. One then tries to prove this positivity by layers, by
comparison with a model kernel whose positivity is known, or by total positivity.

\begin{remark}[equivalent kernel forms]\label{w5:obs:kernelequiv}
We expect the window problem for all $T$ to be equivalent to each of the following: positive semidefiniteness of the
temporal kernel $A$ on its half-line; positive semidefiniteness of $A$ on any set of real nodes accumulating suitably
(by holomorphy, local positivity propagates); positivity of the Gram matrices $\Gamma_{[J,N]}$ for every tail $J\le N$,
for any fixed $J$. There is no equivalence of the minima at fixed $T$. One implication of this type is proved: with
$h_a(y)=e^{-(a-\frac12)y}\mathbf 1_{y>0}$ and $K(a,b)$ the value of the polarized complete form on $(h_a,h_b)$, positive
semidefiniteness of the matrices $[K(32+i,32+j)]_{0\le i,j\le m}$ for every $m$ implies $\WT\ge0$ on $\DT$ for every $T$
\cite[Proposition~11.1]{P46}. A kernel of this type also appears as a ``theta correlation'' kernel in unrefereed
sources; we could not locate a published reference, and we do not claim the equivalences as new.
\end{remark}

The partial results are these. A hyperbolic comparison model, the kernel obtained by replacing the theta function by
$\exp(-a\cosh2t)$, is expected to have positive Gram matrices of all orders, through a Green function representation.
The true kernel is not a mixture, a convolution or a product of $\cosh$-factors of the model, and no differential
equation with fixed prefactor or Green function closure realizes it; we do not know how to transfer positivity from the
model. The first primitive vector $E_1$ has $\WT(E_1)>0$, and the Gram matrix of $\{E_1,E_2\}$ is positive definite
(interval arithmetic), although a natural partial form is negative on $E_1+E_2$. A resonant part $R_{\mathrm{res}}$ of
the form admits an identity of graph-Laplacian type, a sum of squares
$\sum_qD_q|f_q|^2+\sum_{q,m}\Lambda(m)I_q|f_q-f_{q/m}|^2$ with $D_q>0$, positive definite at all orders with a uniform
margin. The non-resonant remainder is not controlled by it: there is no $\delta>0$ with $A\ge\delta R_{\mathrm{res}}$, no
lower bound of the off-diagonal part by a multiple of $R_{\mathrm{res}}$, and an isolated frequency already gives an
indefinite $2\times2$ block.

\begin{remark}[obstructions to comparison]\label{w5:obs:kernelbar}
The following is indicated by computations (in ball arithmetic for (i)) and by arguments that we have not written in
full. (i) For no $c>0$ is $K_M-c/(u+v)$ positive semidefinite at all orders on any interval (no uniform Cauchy margin), and there
is no $\varepsilon>0$ with $K_\Phi\ge\varepsilon K_{2\pi}$; here $K_M$, $K_\Phi$ are the
Laplace-side kernels of the problem and $K_{2\pi}$ the kernel of the model at parameter $2\pi$. (ii) There are positive
analytic approximants of the kernel, arbitrarily close in every finite set of moments and rates, whose temporal kernel is
not positive semidefinite; hence any proof must use more than finitely many analytic data of the kernel.
(iii) No finitely supported minorant, no real exponential-moment minorant and no fixed-support positive certificate
exists. (iv) Additive comparators (Robin, $\operatorname{asinh}$, hyperbolic, completely monotone types) and all even gauges
tried fail; a positive comparator with defect $o(1/u)$ forces the defect to vanish.
\end{remark}

A function $\Phi$ is a P\'olya frequency function of order $k$ ($\mathrm{PF}_k$) if all minors of order $\le k$ of
$(\Phi(x_i-y_j))$, $x_1<\dots<x_r$, $y_1<\dots<y_r$, are nonnegative \cite{Schoenberg1951}; $\mathrm{TP}_2$ is the
corresponding notion for kernels at order $2$. Total positivity of the theta kernel would give positivity of many Gram
determinants at once. It fails: the temporal kernel is not $\mathrm{TP}_2$; positive semidefiniteness does not imply
conditional positive definiteness of $\log A$ (there is a quartic counterexample); and the $\mathrm{PF}_k$ properties
needed appear to fail for $k\ge5$. The classical P\'olya-frequency and Laguerre--P\'olya criteria for $\Xi$ are recalled
in \cite{P14}.

In summary, the true kernel belongs to no positive class that contains the model. Finite analytic data never suffice
(Remark~\ref{w5:obs:kernelbar}(ii)), so the arithmetic must enter in infinitely many places, consistently with
Theorem~\ref{w3:bar:B11}, and positivity of a favourable part (the resonant sum of squares, the first two primitive
vectors) does not control the remainder. What remains open is positive semidefiniteness of the true kernel by an
argument using its full theta and arithmetic structure: a joint bound of the resonant and non-resonant parts, degrees
beyond the plane $\{E_1,E_2\}$, or a positive Green identity. In our view what is missing is a uniform structural
estimate of the full Schur block, not positivity of favourable summands.

\phantomsection\label{w5:subsec:power}%
\emph{Pole-restricted test spaces.} Here one restricts to test functions whose Laplace transform is rational with
prescribed real poles $u_j>\frac12$: atoms $s^k/(s+u)^{N+k}$, combinations of finitely many powers or poles, Laguerre,
Hermite or thermal bases. On such a subspace the complete form is a finite matrix whose entries are explicit sums over
zeros, which can be split into a verified range (where the zeros are known to lie on the line) and a tail; the aim is to
prove positivity on growing subspaces and to reduce the general case to a table of scalar inequalities for the atoms.
The results, all computer-assisted and not reproduced here except where a reference is given, are as follows. There is a
reduction to atoms: positivity of all atoms $f_{N,k,u}$ (all $N,k$) would suffice for positivity on the whole family, and
so would uniform lower bounds $A_{N,k}\ge-\varepsilon_N$ with $\varepsilon_N\to0$ on $0\le k<10^4N^2$. Positivity holds on
planes of two powers, on blocks with up to $13$ arbitrary poles (with explicit lower constants: the complete form, that
is, the zero-side sum of \eqref{w1:eq:explicit}, exceeds $\|F\|^2/4000$; this is reported without proof, for poles
at least $63/2$, in \cite[Section~11]{P46}), on
Laguerre blocks of dimension $12$, and on large rows of the table of atoms. Using the verification of the zeros up to
height $3\cdot10^{12}$ \cite{PlattTrudgian2021} and explicit bounds for the zero-counting tail \cite{Trudgian2014}, the
complete form is positive on every combination of up to $10\,000$ powers; with growing orders the floor of positivity
tends to the bottom of the complete form on the corresponding window, and the certified floors become extremely small
(of the order of $e^{-3.3\cdot10^5}$ for $4096$ poles). In the other direction there is a heuristic lower bound for
the dimension: a zero $\frac12+\delta+i\Gamma$ forces a negative eigenvalue only on blocks with at least of the order of
$\Gamma\log\Gamma\log(2/\delta)/\delta$ powers. And a model form preserving all the structural properties used
(positivity on all blocks of bounded size, the reduction to atoms, the tail estimates) is negative at
$(N,k)=(40402,40401)$, so these properties alone do not force the sign.

Each of these results covers a bounded order, a bounded number of crossings, or a small reciprocal budget of poles.
By the dimension bound an off-line zero at height $\Gamma$ is invisible to blocks of smaller size, so positivity on a
family of fixed size is implied by the verification of zeros up to a corresponding height and decides nothing beyond
it, while the dense confluent regime (arbitrary numerators, all orders) is the window problem itself, as
\cite[Proposition~11.1]{P46} shows for the exponents $32+i$. Induction from pairs fails: a $3\times3$ principal minor
of a reciprocal Pfaffian family equals $-13/270000$. The rows $N\ge18$ of the table of atoms, in the range
$k<10^4N^2$, are undecided, and for the reason just given we do not expect a finite computation to decide them.

\phantomsection\label{w5:subsec:hankel}%
\emph{Small Hankel and Prony blocks.} A last idea is to detect off-line zeros by the sign of small Hankel determinants
of smoothed zero-side or prime-side readings $S_j$, for instance $S_0S_2-S_1^2$ at a moving scale; positivity of such
blocks for all scales would be a criterion. In sublinear regimes (block index $j\to\infty$ with $j/d\to0$) we expect
the following: the eventual sign of the $2\times2$ determinant is positive for every such sequence if and only if there
are finitely many orbits, under $\rho\mapsto1-\rho,\bar\rho$, of zeros off the critical line. In particular a finite
set of exceptions is invisible. Heuristically, $3\times3$ Hankel determinants of five readings give a criterion of the
form ``$\det\ge-A_j$ for large $j$'' that is equivalent to the Riemann hypothesis, and in planes spanned by a pulse and
its translate positivity is observed up to large distances. So for blocks of fixed size in sublinear regimes detection
is either insensitive to finitely many exceptions or equivalent to the conclusion.

For the Hodge-index and hyperbolicity reformulations and the P\'olya-frequency and Laguerre criteria see \cite{P14};
for eventual versus uniform positivity of hierarchies \cite{P15}; for window sums and finite-impulse-response blocks of
the Li sequence \cite{P38}.

\section{Operator-theoretic and function-theoretic criteria}
\label{w6:sec:operators}

The approaches of this section rewrite the sign of the Weil form as a property of an operator, of a
reproducing-kernel space, of an analytic function or of a sequence of numbers attached to $\xi$,
rather than estimating the form on the prime side as in the preceding sections. Each of them ends
either in an exact restatement of the window problem (Problem~\ref{w1:prob:window}) or in a
statement that holds for every weight sequence and so cannot see the primes. Throughout, $\WT$,
$\aT$, $\PT$ and $\lam$ are as in Section~\ref{w1:sec:setting}, and the zero side of $\WT$ is the
right-hand side of the explicit formula~\eqref{w1:eq:explicit}.

\subsection{Indefinite metrics, positivity criteria for \texorpdfstring{$\xi$}{xi} and exact identities}
\label{w6:fam:R10}

The explicit formula writes $\WT(F)=\sum_\rho \PhiF(\rho)\overline{\PhiF(1-\bar\rho)}$. If every
zero is on the critical line, $1-\bar\rho=\rho$ and each term is $|\PhiF(\rho)|^2\ge0$. Without
that assumption the zero side is a Hermitian form $\langle J\,EF,EF\rangle$, where $E$ is
evaluation at the zeros and $J$ is the involution induced by $\rho\mapsto1-\bar\rho$. The zero side
therefore lives naturally in a Kre\u{\i}n space, and the sign of the form becomes a statement about
the position of the range of $E$ relative to the fundamental decomposition. One may hope to bound
the angular operator of that range, or to show that the negative index of the form is finite and
then zero. The realization itself is available.

\begin{theorem}[Kre\u{\i}n identity and angular operator]\label{w6:thm:krein}
For band-limited test functions $F$ in a dense class $\mathcal D_d$ of the Paley--Wiener space
$PW_d$, the zero side satisfies $\mathfrak t(F)=\langle JEF,EF\rangle=\|P_+EF\|^2-\|P_-EF\|^2$ with
$P_\pm=\frac12(I\pm J)$. The projection $P_+$ is injective on the range $R(E)$, which is the graph
of a (possibly unbounded) operator $K_d$ with domain $P_+R(E)$ (with $\|K_d\|=\infty$ if it is unbounded), and
\[
\inf_{F\neq0}\frac{\mathfrak t(F)}{\|P_+EF\|^2}=1-\|K_d\|^2 .
\]
In particular $\|K_d\|\le1$ if and only if $\mathfrak t\ge0$ on $\mathcal D_d$, and the Riemann
hypothesis holds if and only if $\|K_d\|\le1$ for every $d>0$.
\end{theorem}

This is \cite[Theorem~1.2 and Proposition~1.3]{P08}; the injectivity of $P_+$ uses Selberg's
theorem on a positive proportion of zeros on the line. A bounded self-adjoint realization of
$\mathfrak t$ is obtained in \cite[Theorem~1.1]{P08} only under a coercivity hypothesis
(Coer$_d$), a lower bound for the on-line sampling norm, which is not proved. A local density
condition on on-line zeros in intervals of length $\pi/d$, which holds under the Riemann
hypothesis, is proposed there as a sufficient condition, but its implication to (Coer$_d$) is not
proved, and we do not know whether (Coer$_d$) follows from the Riemann hypothesis. In any case
(Coer$_d$) would make the realization bounded, not signed.

On the side of negative indices, the statements available are conditional. Assume that the
off-line zeros form finitely many quartets $\{\rho,\bar\rho,1-\rho,1-\bar\rho\}$, say $m$ of them.
Then on even real test functions the zero side splits into nonnegative squares over the on-line
zeros and one Lorentzian block per off-line quartet, so that the negative index is at most $m$
\cite[Theorem~3.1 and Corollary~6.5]{P18}; equality is left open there. On all test functions the
negative index $\kappa(Q)$ is at most $2m$ \cite[Lemma~4.6]{P25}, and the explicit Pontryagin
completion in zero coordinates, with index exactly $2m$, is \cite[Proposition~4.7]{P25}. Under a
different hypothesis, namely that the prime tail beyond a cutoff satisfies a one-sided relative
form bound with constant below $1$, the negative index is finite with an explicit bound
\cite[Proposition~2.1]{P11}. A further reformulation in the same spirit is
\cite[Theorem~4.1 and Corollary~4.2]{P35}: for $\omega>0$ the number of negative squares of the
kernel of $\Theta_\omega(z)=\xi(\frac12+\omega+iz)/\xi(\frac12+\omega-iz)$ equals the number
$P_\omega$ of poles of $\Theta_\omega$ in the upper half-plane, by the Kre\u{\i}n--Langer count of
negative squares, so that the Riemann hypothesis is equivalent to this index being $0$ for every
$\omega>0$. Here $P_\omega\le N_\omega$, the number of zeros of $\xi$ with $\Re\rho>\frac12+\omega$
counted with multiplicity, with equality for all $\omega$ outside a countable set
\cite[Theorem~4.1]{P35}. The count is exact, but it converts the index into a count of off-line
zeros and supplies no mechanism that forces it to vanish.

The extension problem of Kre\u{\i}n gives a restatement of the same kind at a single window. The
explicit formula prescribes, on $(-2T,2T)$, the Fourier transform of a positive measure on the
line of ordinates; let $\mathrm{KEP}(T)$ be the problem of extending this datum to a positive
definite tempered distribution on $\R$ \cite[Definition~6.6]{P46}. If $\mathrm{KEP}(T)$ is
solvable then $\lamW\ge0$ \cite[Theorem~6.9]{P46}, and for every $T\ge\frac12\log2$
\[
\lamW>0\ \Longrightarrow\ \mathrm{KEP}(T)\ \text{has more than one solution}\ \Longrightarrow\ \lamW\ge0
\]
\cite[Theorem~6.10]{P46}. Thus indeterminacy of the extension problem and strict positivity of the
complete form coincide except at the windows where $\lamW=0$. A Nevanlinna parametrization, in
$T$, of the set of solutions has not been written.

These approaches cannot give the sign, for the following reason. A faithful indefinite realization
of the zero side contains exactly the information of the explicit formula; the sign is the
statement $J=I$ on $R(E)$, that is, the position of the zeros. By Theorem~\ref{w6:thm:krein} a
bound $\|K_d\|\le1$ is equivalent to the positivity it is meant to prove, so it requires either
information on the zeros, which is circular, or an arithmetic inequality of the same strength as
the window problem. The index statements are finite only under hypotheses (finitely many off-line
quartets, a relative form bound with constant below $1$) that are themselves unproved, and the
index theory counts the negative squares but never forces their number to be zero. The general
theory offers no help either: self-adjointness or unitarity with respect to an indefinite inner
product does not force real spectrum, so a Stone-type argument for a one-parameter group of
$J$-unitary operators does not place the zeros on the line. Computations indicate, moreover, that
the Carleson-type sampling constant of the band realization is identically $1$, so that
incoherence of the primes buys no margin; and numerical estimates (uncertified) of the finite
defect and of the cutoff for a relative form bound slide with the band parameter. The Kre\u{\i}n--Langer
theory of generalized resolvents and of negative squares of kernels is used through
\cite{P35}; the other statements are in \cite{P08,P11,P18,P25}.

\phantomsection\label{w6:fam:R11}
Several classical criteria express the Riemann hypothesis as positivity of an analytic object
attached to $\xi$: a de Branges structure function \cite{deBranges1968}, a Herglotz function, a
Pick matrix, or a self-adjoint operator whose spectrum is the set of ordinates
\cite{BerryKeating1999,Connes1999}. One may try to establish one of these positivity properties
directly, or to build the operator. The simplest of the criteria is the following.

\begin{proposition}[Herglotz criterion]\label{w6:prop:herglotz}
The Riemann hypothesis holds if and only if $\Rel\,\xi'(s)/\xi(s)>0$ for $\Rel s>\frac12$.
\end{proposition}

This is due to Lagarias \cite{Lagarias1999} (see also \cite{Hinkkanen1997}); an exposition with the
boundary identity $\Rel\,\xi'/\xi(\frac12+i\tau)=0$ off the ordinates is \cite[Proposition~2.1 and
Theorem~3.1]{P20}. It is an equivalence and carries no information on how to prove positivity: in
particular it does not see the window $T$ at which a hypothetical first failure of the window
problem would occur.

The de Branges positivity condition proposed as an approach to the Riemann hypothesis is known to
fail: Conrey and Li \cite{ConreyLi2000} showed that the relevant positivity property does not hold.
A realization of the localized Weil form satisfying the three de Branges axioms is a space
$\mathcal H(E)$ in the Pontryagin sense whose negative index is at least the index coming from
off-line quartets \cite[Theorem~8.4]{P18}; so the de Branges structure, when it exists, records the
off-line zeros but does not exclude them.

The arithmetic Pick chain of \cite{P36} is a fifteen-step Nevanlinna--Pick reformulation: an
arithmetic Pick positivity statement, called ARP-P there, implies the Riemann hypothesis by
Nevanlinna--Pick interpolation and a normal-family argument, and conversely
\cite[Theorems~9.241 and 9.245, Corollary~9.246]{P36}. The open input is ARP-P itself, which is
equivalent to the Riemann hypothesis and, at its arithmetic end, to the Li--Keiper criterion
\cite[Corollary~9.247]{P36}. Of the supporting criteria, scalar Pick positivity for fixed finite
cluster patterns at large height is proved, with an effective threshold for a fixed pattern
\cite[Theorem~9.43]{P36} and a non-effective one uniformly in a band \cite[Theorem~9.64]{P36}; it
bears on no window. It is plausible that positivity known at one point does not suffice for such a
criterion, whereas positivity on an interval does; we have not proved this.

For Hilbert--P\'olya operators the situation is as follows. A unitary representation of the
multiplicative group $\R_+^\times$ whose trace on test functions is $\sum_\rho m_\rho\PhiF(\rho)$,
with multiplicities $m_\rho$, would force $\Rel\rho=\frac12$; heuristically, constructing it is
equivalent to the problem, in the spirit of \cite{Connes1999}. An involution anticommuting with
the scaling generator and compatible with the functional equation should not exist, because the
involution induced by the functional equation commutes with the scaling generator; this argument
is heuristic. A third idea uses ladder operators: for the scaling generator $H$ and an operator
$L$ satisfying $[H,L]=cL$ with $c\ne0$, one hopes that $L$ produces arithmetic progressions in the
spectrum. For the specific scaling generator of the construction we expect that a bounded $L$ with
$[H,L]=cL$ must vanish, but we have no proof, and the precise hypotheses matter, as the next
proposition shows.

\begin{proposition}[ladder operators for a bounded generator]\label{w6:prop:ladder}
Let $H$ and $L$ be bounded operators on a Banach space with $HL-LH=cL$ for some $c\neq0$. Then $L$ is
nilpotent. If $H$ is unbounded, $L$ need not be nilpotent.
\end{proposition}
\begin{proof}
By induction $HL^k-L^kH=kcL^k$ for every $k\ge1$, hence $k|c|\,\|L^k\|\le2\|H\|\,\|L^k\|$; for
$k>2\|H\|/|c|$ this forces $L^k=0$. For the second statement take $H=-i\,d/dx$ on $L^2(\R)$ and $L$
multiplication by $e^{icx}$: then $[H,L]=cL$ on the domain of $H$ and $L$ is unitary.
\end{proof}

Nilpotency is not vanishing (on $\C^2$, $H=\mathrm{diag}(0,c)$ and $L=E_{21}$ satisfy $[H,L]=cL$),
so any statement ``$L=0$'' depends on spectral properties of the particular $H$; and for the
unbounded scaling generator the ladder operators exist and carry no arithmetic. We do not know
what happens for an unbounded $L$ and the actual scaling generator of the construction.

A family of operator programmes realizes the form as a Schr\"odinger-type operator, pole plus
archimedean potential minus prime potential, on a finite interval of length about $\log\lambda$,
with the sign equivalent to the absence of negative eigenvalues. Their output is a marginality
statement: the half-width of the potential well, $\delta_\lambda=1/(\sqrt2\log\lambda)$, and the
resolution of the interval, $r_\lambda=\pi/\log\lambda$, are in fixed ratio
$2\delta_\lambda/r_\lambda=\sqrt2/\pi$, and this suggests, heuristically, that no functional of the
magnitudes $\Lambda(n)$ alone decides the sign. The isolated remainders turn out to be of the size
of the Riemann hypothesis itself, and the last step of each programme is an
essential-self-adjointness or limit statement which, as far as we can see, is equivalent to it. A
useful elementary remark from these programmes concerns the pole: the matrix
$G_R=\begin{pmatrix}R&\sqrt R\\ \sqrt R&1\end{pmatrix}$ is positive semidefinite of rank one for
every $R>0$, while $G_R-R\,e_1e_1^{\!\top}$ is indefinite; a divergent pole direction must be
removed by passing to a Schur complement or a quotient, not by subtraction.

None of these criteria gives the sign. The Herglotz and Pick criteria are equivalences whose
positivity is a property of the zeros
(Proposition~\ref{w6:prop:herglotz}; \cite[Corollary~9.246]{P36}); proving them requires a
prime-side bound of the same strength as the problem. De Branges structures exist only under
positivity or record the defect, and the specific positivity hypothesis of de Branges is false
\cite{ConreyLi2000}. Operator constructions either use only the archimedean data and the scaling
action, and then do not distinguish $\zeta$ from models with off-line zeros, or install the zeros
by hand. The prolate operator of Connes, Consani and Moscovici belongs to the first kind; we
return to it in Section~\ref{w6:sec:prolate}.

\phantomsection\label{w6:fam:R12}
A third approach builds a long chain of exact identities---relative Fredholm determinants,
canonical systems, Feshbach--Schur reductions to high and low frequencies, bordered determinants
and adjugates, telescoping identities for differences of consecutive Galerkin quantities---in the
hope that the last link isolates a single arithmetic quantity whose sign is evident. In our
experience each finite identity of such a chain can be established, and so can several
implications between their end-points; we do not reproduce these computations and state none of
them as a theorem. One example illustrates the pattern. Computations indicate that the difference
of consecutive Galerkin quantities $g_{N+2}-g_N$ decomposes exactly into a zero term, a mesh term
of size $O(\sigma/N^2)$ and a boundary term of size $O(\sigma/N^3)$. The zero term carries the
whole difficulty; numerically (uncertified) $N^2$ times the zero term grows, with exponent about
$1.2$. A second pattern is the Douglas-type factorization.

\begin{proposition}[contraction form of a positive difference]\label{w6:prop:douglas}
Let $R$ be a bounded positive operator on a Hilbert space with bounded inverse and $W$ a bounded
operator into another Hilbert space. Then $R-W^*W\ge0$ if and only if $\|WR^{-1/2}\|\le1$.
\end{proposition}
\begin{proof}
Put $x=R^{-1/2}y$. Then $\langle(R-W^*W)x,x\rangle=\|y\|^2-\|WR^{-1/2}y\|^2$, and $x$ ranges over
the whole space when $y$ does.
\end{proof}

This is the bounded-inverse case of Douglas' majorization lemma \cite{Douglas1966}. Suppose the
primitive form is written as $R_T-B_T^*B_T$ with an explicit positive reference $R_T$ and a bounded
operator $B_T$. Proposition~\ref{w6:prop:douglas}, with $W=B_T$, turns the sign into the
contraction condition $\|B_TR_T^{-1/2}\|\le1$. In the factorizations of this kind that we know,
the construction works for every nonnegative weight sequence: the positivity of the reference
comes from the positivity of the weights, and the contraction inequality is then equivalent to the
sign it replaces. Related results are the determinant serialization \cite[Theorem~9.39]{P36}, an
equivalent reformulation, and the one-point jet criteria \cite[Theorems~9.53 and 9.100]{P36},
which are reformulations of Li or Pick type; see \cite[Section~9]{P36}.

An exact identity transports the form without changing its sign content: at the end of the chain
the arithmetic quantity is again the window problem, or a statement equivalent to it. We have no
theorem that excludes this approach, and no chain has been completed to a signed identity; several
end-points appear to be different faces of one missing signed identity. Two reasons for the
failure can be given heuristically. Estimating shells of the chain separately before pairing them
destroys the cancellation between the pole, the archimedean term and the primes; and a positive
lift of the form satisfies at most two of three requirements, namely local positivity, the global
phase of the explicit formula, and descent to primitive functions. Two objects met along the way
were not pursued to the end: a signed identity for the zero term of the telescoping decomposition,
derived from the explicit formula alone, and a Hermitian covariance identity for balanced
quotients.


\subsection{Li--Keiper coefficients and polynomial hierarchies}
\label{w6:sec:li}\label{w6:fam:R13}

Two approaches replace the continuous family of windows by a discrete sequence of conditions, each
of which can be checked: the Li--Keiper coefficients, and the hierarchies of finite-order
positivity conditions (Jensen polynomials, P\'olya frequency orders, heat flow). They share one
structure: every finite level is unconditional or verifiable, and the full statement is the limit.

Let
\[
\lambda_n=\sum_\rho\Bigl[1-\Bigl(1-\frac1\rho\Bigr)^{n}\Bigr]\qquad(n\ge1),
\]
the sum over the nontrivial zeros taken symmetrically. Equivalently
\[
\lambda_n=\frac1{(n-1)!}\frac{d^n}{ds^n}\bigl[s^{n-1}\log\xi(s)\bigr]_{s=1}.
\]

\begin{theorem}[Li's criterion]\label{w6:thm:li}
The Riemann hypothesis holds if and only if $\lambda_n\ge0$ for every $n\ge1$.
\end{theorem}

This is \cite{Li1997}; the coefficients appear earlier in \cite{Keiper1992}, and
\cite{BombieriLagarias1999} extend the criterion to multisets of points and prove that under the
Riemann hypothesis $\lambda_n=\frac n2\log n+c\,n+O(\sqrt n\log n)$ for an explicit constant $c$.

One splits $\log\xi(s)=-\log2+\bigl[\log s-\frac s2\log\pi+\log\Gamma(\frac s2)\bigr]+\log\bigl((s-1)\zeta(s)\bigr)$
(the constant does not contribute to $\lambda_n$) and correspondingly
$\lambda_n=\lambda_n^{\mathrm{arch}}+\lambda_n^{\mathrm{prime}}$, an exact and unconditional
decomposition. The archimedean part is explicit:
\begin{equation}\label{w6:eq:liarch}
\lambda_n^{\mathrm{arch}}=1-\frac n2\bigl(\gamma+\log4\pi\bigr)
+\sum_{r\ge1,\ r\ \mathrm{odd}}\Bigl[\Bigl(1-\frac1r\Bigr)^{n}-1+\frac nr\Bigr],
\end{equation}
and the prime part is expressed through the Stieltjes constants, or, after an Abel regularization
keeping the pole and the prime sum together, through an integral against $d\psi(y)-dy$. The
target is the one-sided inequality
\begin{equation}\label{w6:eq:lionesided}
\lambda_n^{\mathrm{prime}}\ \ge\ -\lambda_n^{\mathrm{arch}}\qquad\text{for every }n,
\end{equation}
which is equivalent to Theorem~\ref{w6:thm:li}.

\begin{proposition}[sign of the archimedean part]\label{w6:prop:liarch}
$\lambda_n^{\mathrm{arch}}<0$ if and only if $1\le n\le7$. Numerically
$\lambda_1^{\mathrm{arch}}=1-\frac\gamma2-\frac12\log4\pi=-0.55412\ldots$,
$\lambda_7^{\mathrm{arch}}=-0.35573\ldots$, $\lambda_8^{\mathrm{arch}}=0.02089\ldots$.
\end{proposition}

This is \cite[Lemma~9.121]{P36}: the second differences of \eqref{w6:eq:liarch} equal
$\sum_{r\ \mathrm{odd}}(1-1/r)^n r^{-2}>0$, and the first differences are evaluated through
$\zeta(2),\dots,\zeta(8)$. For $n=1$ the series in \eqref{w6:eq:liarch} vanishes, and since
$\lambda_1=1+\frac\gamma2-\frac12\log4\pi=0.02309\ldots$ one has $\lambda_1^{\mathrm{prime}}=\gamma$.

\begin{corollary}[absolute domination fails]\label{w6:cor:absdom}
The inequality $|\lambda_n^{\mathrm{prime}}|<\lambda_n^{\mathrm{arch}}$ fails for $1\le n\le7$, in
particular at $n=1$.
\end{corollary}
\begin{proof}
For these $n$ the right-hand side is negative by Proposition~\ref{w6:prop:liarch}.
\end{proof}

So the natural strategy of showing that the archimedean part dominates the prime part in absolute
value can start at $n=8$ at the earliest, and for $n\le7$ the prime part must carry the sign.
Numerically (uncertified) domination fails also at $n=8,9,10$
($\lambda_8^{\mathrm{prime}}=1.4449$ against $\lambda_8^{\mathrm{arch}}=0.0209$;
$\lambda_{10}^{\mathrm{prime}}=1.3238$ against $0.9555$) and first holds at $n=11$
($1.2494<1.5009$). Absolute domination from some $n_0\ge11$ on, combined with a certified finite
range, is not excluded by Corollary~\ref{w6:cor:absdom}; it is not proved, and under the Riemann
hypothesis it holds for large $n$ by \cite{BombieriLagarias1999}.

The finite range is settled numerically with a certificate: $\lambda_n>0$ for $1\le n\le149$ is
certified by ball arithmetic in a script included with \cite{P37}, and
\cite[Proposition~9.137]{P36} gives an independent interval check of $2\lambda_n-A_n>0$ for
$9\le n\le20$ (for $21\le n\le149$ only under an unproved coefficient-tail bound). Here $A_n$ is the
tail quantity of the compact--tail decomposition of \cite[Theorem~2.3]{P37}. The reductions of the
same paper give a cofinal criterion: if $4\lambda_n-A_n>0$ for all large $n$ in the admissible
sense of \cite[Theorem~2.10]{P37}, the compact--tail inequality follows; by Li's criterion and
the theorem of Bombieri and Lagarias the hypothesis is at least as strong as the Riemann hypothesis. A tail
lemma based on the prime number theorem is proved in \cite[Lemma~9.132]{P36}; the remaining
arithmetic inequality implies the target \cite[Theorem~9.135 and Corollary~9.136]{P36}, is open,
is at least as strong as the Riemann hypothesis, and is not known to follow from it except for
large $n$. FIR reformulations of Li's criterion are in \cite[Theorems~3.1 and 4.1]{P38}.

Sixteen attempts to prove the compact--tail inequality by transferring the positivity of an
auxiliary object are shown in \cite{P37} to fail relative to their declared inputs, most by an
exact witness: for instance an explicit $\mathrm{PF}_2$ minor equal to $-1/2916$
\cite[Theorem~6.1 and Witness~6.1]{P37}, a polar--Gamma density equal to
$e^{-\varepsilon x_0}\bigl(-r+(r-1)/\sqrt2\bigr)<0$ at $x_0=\frac12\log2$ for every exponent $r>1$
\cite[Theorem~12.1]{P37}, and the fact that the holomorphy needed for Fej\'er summation is
equivalent to the Riemann hypothesis \cite[Theorems~8.3 and 8.4]{P37}.

Three further observations delimit what the coefficients can do. First, an off-axis zero forces
exponential growth of the Li coefficients, since a pair with $|1-1/\rho|>1$ contributes a geometric
term, while the converse requires a subexponential prime bound of the strength of the problem;
compare the detection remark following \cite[Lemma~9.121]{P36}. Second, the bound
$|\lambda_n^{\mathrm{prime}}|=O(\sqrt n\log n)$ is equivalent to the Riemann hypothesis: it
follows from it by the asymptotic of \cite{BombieriLagarias1999} together with
$\lambda_n^{\mathrm{arch}}=\frac n2\log n+cn+O(1)$, and conversely a lower bound
$\lambda_n\ge-c(\varepsilon)e^{\varepsilon n}$ for every $\varepsilon>0$ implies the Riemann
hypothesis \cite{BombieriLagarias1999}; it is therefore a restatement. Third, the Li functional
cannot serve as a dual certificate for the window problem: Li's criterion is a term-by-term
statement over the zeros, whereas the value function of a flat-top certificate always contains
terms of both signs (Proposition~\ref{w5:prop:noPD}).

The inequality \eqref{w6:eq:lionesided} has not been proved by any of these means. Magnitude bounds
on the prime part, obtained from explicit forms of the prime number theorem, fall short of the
archimedean budget $\frac n2\log n$ by a factor of order $n$: the bound we obtain is of order
$n^2\log n$ against a budget of order $n\log n$, with constants that are not settled. No local
unit---a single prime, a finite shell, a Laguerre lobe, a Fej\'er--Riesz square---appears to have a
sign; numerically the contributions of individual primes or prime towers to
$\lambda_n^{\mathrm{prime}}$ have both signs already at $n=2$, so the cancellation is global.
Positivity transfers from auxiliary objects fail by the explicit witnesses of \cite{P37}. What
remains is \eqref{w6:eq:lionesided} for $n\ge8$, in the regularized form
$\lim_\varepsilon\int(\psi(y)-y+1)f'_{n,\varepsilon}(y)\,dy\ge-\lambda_n^{\mathrm{arch}}$ for the
regularizing kernels of \cite{P36}; with the certified finite range it is equivalent to the
Riemann hypothesis.

\phantomsection\label{w6:fam:R14}
We turn to the polynomial hierarchies. Write $\Xi(t)=\xi(\frac12+it)=\int_\R\Phi(u)e^{itu}\,du$
with the de Bruijn--Newman kernel $\Phi$. The Riemann hypothesis is the statement that $\Xi$ has
only real zeros, and several hierarchies of finite conditions are equivalent to it: hyperbolicity
of all Jensen polynomials of $\Xi$ (P\'olya \cite{Polya1927}); positivity of the Hermite forms of
these polynomials (a polynomial is real-rooted if and only if its Hermite form is positive
semidefinite); the condition that the de Bruijn--Newman constant $\Lambda$ satisfies $\Lambda\le0$
\cite{deBruijn1950,Newman1976}. A related family that is not equivalent to the Riemann hypothesis
is the P\'olya frequency order of $\Phi$: by Schoenberg's theorem \cite{Schoenberg1951} a P\'olya
frequency function of infinite order has a bilateral Laplace transform that is the reciprocal of a
Laguerre--P\'olya function, whereas that of $\Phi$ is entire, so $\Phi$ is not
$\mathrm{PF}_\infty$; the finite orders can be examined as possible sources of sign. The approach
is to prove each finite level and pass to the limit, or to run the heat flow backwards, in the
manner of the theory of stability-preserving operators and interlacing families
\cite{BorceaBranden2009,MSS2015}.

The finite levels are largely known. Griffin, Ono, Rolen and Zagier \cite{GORZ2019} prove that for
each degree $d$ the Jensen polynomials $J^{d,n}$ of $\Xi$ are hyperbolic for all sufficiently large
shifts $n$, and for $d\le8$ for all $n$. Rodgers and Tao \cite{RodgersTao2020} prove $\Lambda\ge0$,
so that the Riemann hypothesis is the statement $\Lambda=0$, the boundary case of the flow. The
Taylor coefficients of $\Xi$ are, up to the sign $(-1)^k$ and the factor $2^{2k+3}/(2k)!$, the
even moments of $\Phi$ \cite[\S3]{P14}, and the classical Hermite, P\'olya frequency and Laguerre
criteria are collected in \cite{P14}. For real polynomials with simple zeros there is a
monotonicity identity for the off-line mass under the backward heat flow \cite[Theorem~2.1]{P13};
it holds for every conjugation-symmetric configuration and so does not use arithmetic. For
scale-free hierarchies eventual and uniform positivity coincide, and a bounded threshold gives
uniform positivity \cite[Theorem~3.1]{P15}; Hodge--Riemann gaps are continuous under convergence
of forms \cite[Lemma~8.2]{P15}, and positivity at every finite level holds regardless of the limit
\cite[Corollary~8.3]{P15}. As for the P\'olya frequency orders, the preprint
\cite{Michalowski2026} states as its main theorem that $\Phi$ is not a P\'olya frequency function
of order $5$, by a $5\times5$ minor certified by interval arithmetic; the preprint is unrefereed
and we have not checked it. If it is correct, no argument through $\mathrm{PF}_k$ with $k\ge5$ is
possible. We do not know whether $\Phi$ is $\mathrm{PF}_3$ or $\mathrm{PF}_4$; such a property
would not by itself imply the Riemann hypothesis. Computations also indicate a second-order
balance law along the de Bruijn--Newman flow, bounding a weighted sum of squared imaginary parts
by $O(T/\log T)$; it carries no sign information.

The hierarchies do not give the sign because each finite level is unconditional---it holds for
$\Xi$ and also for entire functions with off-axis zeros---while the full hierarchy is the
statement itself; uniform convergence of the Hermite forms is strictly stronger. Heuristically,
every finite-depth heat or moment hierarchy is satisfied by a model with a pair of non-real zeros,
and numerically the finite truncations of the prime sum in the explicit formula are negative in
some window at every finite cutoff. The heat flow and the monotone functionals built on it act
identically on every conjugation-symmetric configuration, so they are blind to the primes; the
Riemann hypothesis sits at the boundary $\Lambda=0$ \cite{RodgersTao2020}, where no margin is
available. In the language of \cite{P15}, uniformity requires either finite type or a bounded
threshold, and the Jensen hierarchy of $\xi$ has neither in any known form.


\subsection{Prolate operators, Sonin spaces and the semilocal trace}
\label{w6:sec:prolate}\label{w6:fam:R15}

Connes' trace formula \cite{Connes1999} writes the Weil distribution as a sum of local
contributions over places, and Connes and Consani \cite{ConnesConsani2021} treat the archimedean
place: on test functions supported in a bounded interval of the multiplicative group, the
archimedean contribution is related to the compression of the scaling action to Sonin's space---the
functions which, together with their Fourier transforms, vanish on a fixed interval around the
origin---and to the prolate operator that commutes with time and band limiting
\cite{SlepianPollak1961,LandauWidom1980}. Their result gives Weil positivity at the archimedean
place for small support, through the fact that the relevant eigenvalues are below one; the
remainder between the archimedean distribution and the trace of the compression is controlled only
up to a term of the form $-c|\widehat g(0)|^2$, a rank-one correction in the direction of the pole.
The natural continuation extends the construction to a finite set $S$ of places (semilocal Sonin
spaces) and tries to write
\[
\WT=\mathrm{Pos}_{S}+\Delta_{T,S},
\]
with $\mathrm{Pos}_S$ a manifestly positive semilocal part and a remainder $\Delta_{T,S}$ whose
negative index on the primitive space $\PT$ is controlled, uniformly in $T$ as $S$ grows. A related
group of approaches uses the prolate operators of Connes, Consani and Moscovici \cite{CCM2024},
their Jacobi matrices and finite Weil quotients, whose spectra are expected to approximate the
ordinates of the zeros.

For the operators of the second group the following is proved. The incidence condition proposed as
a spectral reformulation of the Riemann hypothesis through a limit Jacobi operator fails at every
zero on the critical line, unconditionally, because the potential has a boundary singularity
$2m/\eta$ at an ordinate of multiplicity $m$ \cite[Proposition~3.3]{P23}; and a Jacobi operator
whose coefficients tend to zero is compact, so its spectrum cannot be the set of ordinates
\cite[Proposition~3.2]{P23}. Between consecutive finite Weil quotients of the
Connes--Consani--Moscovici construction there are exact transport identities, and cofinal
summability holds at fixed cutoff under the standing condition (C$_L$) of \cite{P39}, while a model
counterexample shows that abstract Galerkin compactness alone cannot give consecutive summability
\cite[Theorems~3.2, 4.1 and 7.4, Proposition~9.1]{P39}; the global limit is not established.
Positive-form convergence or divisor convergence of the finite quotients would imply the Riemann
hypothesis \cite[Theorem~5.1]{P40}. On each finite row the sign reduces to one strict finite
inequality (the source surplus), which is not proved \cite[Section~9]{P40}; the obstructions met
by that inequality are listed in \cite{P41}, a few of them with proofs and the rest as
observations. Finally, the moment problem of the archimedean measure is determinate, a result of
Connes, Consani and Moscovici recalled in \cite[Lemma~2.5]{P22}, and proportionality of all moments
of the discrepancy measure forces the Riemann hypothesis \cite[Theorems~4.1 and 4.4]{P22}; this is
a reformulation, since evaluating the discrepancy measure needs the zeros. The spectra of the
finite prolate operators are real at every finite level, and we do not know whether they converge
to the ordinates of the zeros.

For the Sonin decomposition we have computations and arguments which are not written out as
proofs, and we state their conclusions as expectations. A semilocal Sonin chain should exist for
every finite set $S$ of places and every window, with exact compatibility between consecutive $S$;
it carries no sign. Semilocal Sonin positivity appears to be invariant under the replacement of
the weights $\Lambda(n)$ by $c\Lambda(n)$, $c>0$, and so cannot distinguish $\zeta$ from the
rescaled weight sequences, for which the window problem fails in some window
(Theorem~\ref{w3:bar:B1}). For the full integer Sonin decomposition we expect the remainder to
have infinitely many negative directions on primitive functions as soon as $T>\frac12\log2$: pairs
of bumps at distance $\log2$ should produce an infinite-dimensional negative subspace. A Sonin
parameter uniform in $T$ should not exist, under a hypothesis on the zeros, and norm bounds on the
Sonin remainder are either blind to weight sequences that agree with $\Lambda$ on the contact set
or, when the remainder is defined through the form, circular. Heuristically the $p$-adic Sonin
space is one-dimensional, so no finite place carries a spectrum related to the zeros. A positivity
statement for an Abel-regularized kernel of the prolate construction, which would be convenient in
several of these arguments, appears to be false: by our computation the integral in question is
zero. For the archimedean part alone on primitive functions, a two-value prolate relaxation
indicates at most $5$, $7$ and $8$ negative directions at the windows $\frac12\log3$, $\log2$ and
$0.80$, and no bound uniform in $T$ is possible, by Theorem~\ref{w2:thm:archcount}. In terms of
boundary triples for the prolate operator, the deficiency indices, the Fourier-invariant
self-adjoint extensions and their invariants seem to depend only on the archimedean parameters,
and to be identical for $\zeta$ and for weight sequences that violate the window problem; one
deficiency-index computation underlying this is not settled. Numerically (uncertified), the number
of eigenvalues larger than $1$ of the archimedean prolate compression grows roughly like $4T-2$,
and a purely archimedean obstruction appears near $T\approx0.80$.

Three obstructions recur. The first is blindness to the weights: the prolate and Sonin structures
are built from the archimedean place and the scaling action, and their invariants depend on the
archimedean parameters and on $2\pi$, not on $\Lambda$. A positivity statement that holds equally
for weight sequences violating the window problem cannot imply it
(Proposition~\ref{w3:prop:principle}); that every statement of this type is of that kind is a
heuristic. The second is the inertia of the remainder: if the Sonin remainder has infinite
negative index on $\PT$ beyond $T=\frac12\log2$, the decomposition $\mathrm{Pos}_S+\Delta_{T,S}$
does not reduce the problem to a finite one. The third is convergence to the explicit formula:
heuristically, a positive part whose remainder has bounded negative index for all $T$ must
converge, in the sense of distributions on the zeros, to the explicit formula itself, and the
approach then ends in a reformulation. Operator-algebraic
positivity of the same kind (states, traces, crossed products) is Bochner positivity, which is the
Weil form again.

Some questions in this direction have not been examined. The inertia, on $\PT$, of Connes'
semilocal residue (identity minus the semilocal compression) has never been computed. A Sonin
decomposition with parameter $\eta<1$ adapted to each window has not been tried. Isometric
completions seem to kill the eigenclasses attached to the zeros, and one may ask for a resonant or
nuclear completion that does not; compare \cite[Theorem~10.7]{P40}, where an explicit charged
completion has absolutely continuous scale spectrum.


\subsection{Detectors of off-line zeros and the zero side}
\label{w6:sec:detectors}\label{w6:fam:R16}

The approaches of this subsection do not try to prove an inequality directly. They build an
observable---a truncated Dirichlet sum, a decomposition by the number of prime factors, a localized
Weil matrix, a second moment over zeros, a statistic---that would behave differently if an off-line
zero existed, and hope to show that the anomalous behaviour does not occur. We describe the
arithmetic detectors, the localized Weil form and the second moments, and then the direct readings
of the zero side of the explicit formula, which are the closest to the window problem.

For a Dirichlet polynomial $D(t)=\sum_{n\le N}a_nn^{-1/2-it}$ split $D=\sum_kS_k$ according to
$\omega(n)=k$. One expects an off-line zero $\beta+i\gamma$ to produce growth of $|D(\gamma)|$ of
order $N^{\beta-1/2}$, or a detectable interference between classes. What is proved is the identity
$|D|^2=\sum_k|S_k|^2+\sum_{j\ne k}S_j\overline{S_k}$ and, by Cauchy--Schwarz,
$|D|^2\le K_N^2\max_k|S_k|^2$ with $K_N\sim\log N/\log\log N$ the number of classes
\cite[Theorems~3.1, 4.1 and 5.1]{P02}; a bound $\max_k|S_k|^2\ll N^{\varepsilon}$ would give a
Lindel\"of-type bound \cite[Corollary~4.2]{P02}. Selberg--Delange asymptotics for logarithmic sums
of $q^{\omega(n)}$ are in \cite[Proposition~3.1]{P21}, with no link to the zeros claimed there;
the interference statistics of \cite{P03} and the per-prime weights of \cite{P05} are numerical.
The decisive test is the Davenport--Heilbronn function, which satisfies a functional equation of
the right type and has zeros off the line \cite{DavenportHeilbronn1936}: its partial sums converge
to $L_{\mathrm{DH}}(\frac12+it)$ with an explicit rate,
$D_{\mathrm{DH}}(t;N)=L_{\mathrm{DH}}(\frac12+it)+O((1+|t|)N^{-1/2})$ \cite[Proposition~2.3]{P04}.
Thus off-line zeros produce no growth of partial sums on the line, and numerically no statistic
separates $L_{\mathrm{DH}}$ from $L(s,\chi_5)$ \cite{P04}. The distribution of $\omega(n)$
\cite{ErdosKac1940} and mean-value estimates of the type of \cite{MontgomeryVaughan1974} govern
these sums, and multiplicativity alone does not suppress resonances: random multiplicative models
behave in the same way.

The localized Weil form restricts the explicit formula to a Hermite--Gauss basis concentrated near
a height $\gamma_0$ and studies the matrix $Q=M_{\mathrm{zeros}}-M_{\mathrm{arith}}$. Splitting a
double point $\frac12+i\gamma$ into $\frac12\pm\delta+i\gamma$ changes the zero-side matrix by
$\delta^2A(\gamma)+O(\delta^4)$ with $A$ of rank at most $3$ and a strictly negative eigenvalue
\cite[Lemma~3.1 and Theorem~3.2]{P07}. The operator norm of the prime tail beyond $X$ is at most
$A\exp(-\alpha(\log X)^2)$ for every $\alpha<\sigma^2/4$, using only a Chebyshev bound
\cite[Theorem~4.2]{P07}, and for a model configuration the truncated matrix has a negative
eigenvalue beyond an explicit cutoff \cite[Theorem~5.1]{P07}. Under the Riemann hypothesis the
zero-side matrix is positive semidefinite for every basis \cite[Proposition~6.1(b)]{P07}, and then
$\lambda_{\min}(M^{(X)}_{\mathrm{arith}})\ge-A\exp(-\alpha(\log X)^2)$ \cite[Corollary~6.2]{P07}.
For model forms of the same shape, the matrix and its spectrum depend on the coefficients beyond a
prime $P$ only up to an explicit Gaussian tail \cite[Theorem~4.1]{P09}. On band-limited tests the
prime side reduces exactly to $n\le e^{2d}$ \cite[Proposition~2.1]{P12}, and on even real tests the
zero side is a sum of squares over on-line zeros plus one Lorentzian block per off-line quartet
\cite[Theorem~3.1]{P18}. Two limitations are expected, without proof: every admissible localized
window should force an oversampling factor tending to infinity; and a statistical relative form
bound for the prime tail holds, heuristically, with constant $C\ge e^{d}>1$, whereas $C\le1$ would
be needed, and no density estimate is expected to bring $C$ below $1$ without assuming the
conclusion. Coercivity of the positive part of the localized form through a density condition on
zeros in short intervals is the hypothesis (Coer$_d$) discussed after
Theorem~\ref{w6:thm:krein}.

A third detector controls the window sum of the zero side through the second moment
$V(d,T)=\int_{-2d}^{2d}\bigl|\sum_{0<\gamma\le T}e^{i\gamma\xi}\bigr|^2d\xi$. One has the
identity and the bounds $0\le V\le4dN^2$ \cite[Theorem~2.1]{P16}, and an upper bound for the
windowed zero sum in terms of $V$ \cite[Proposition~2.2]{P16}, which gives no lower bound for the
full form. Under the Riemann hypothesis, by Montgomery's theorem on the pair correlation function
\cite{Montgomery1973}, there is an asymptotic for the low band and a bound $V_w(d,T)\ll dN(T)$ for
a weighted version \cite[Propositions~3.1 and 4.1]{P16}. A bound $V\ll dN(\log T)^c$ is a
hypothesis in \cite{P16} (Hypothesis $H_c$); we expect that it can be proved unconditionally with
some exponent $c$, but have no written proof. Log-correlated models of $S(t)$, in the spirit of
\cite{Selberg1946,KeatingSnaith2000}, and the correlation statistics of the zeros
\cite{RudnickSarnak1996} describe the zeros and $\zeta$ on the line but force nothing.

We come to the zero side itself. The explicit formula~\eqref{w1:eq:explicit} expresses $\WT(F)$ as
a sum over zeros. Every on-line pair contributes a square; an off-line quartet contributes an
indefinite block. Several statements make precise how far the zero side can be read.

\begin{proposition}[an arbitrary off-line zero is detected]\label{w6:prop:detect}
If $\xi$ has a zero off the critical line, then there are real primitive test functions of compact
support, even and odd, on which the complete Weil form is negative.
\end{proposition}

This is the classical form of Weil's criterion (Yoshida \cite{Yoshida1992}; Connes--Consani); a proof
in the present normalization, with real primitive even and odd tests of compact support, is
\cite[Theorem~6.2]{P45}, and the bilateral parity version is \cite[Theorem~13.2]{P39}. The
support of such tests cannot be small. We expect a quantitative version, giving a window
$\tau(\rho)$ beyond which an off-line zero $\rho$ makes $\lamW$ negative, and we expect that
detection uniform in the height of the zero is false; neither statement is proved.

Two facts restrict any argument on the zero side. Termwise positivity is impossible: the value
function of a flat-top certificate always contains terms of both signs, even if every ordinate is
real (Proposition~\ref{w5:prop:noPD}). And no finite set of zeros carries a certificate: for a
nonempty finite symmetric multiset $Z$ of real ordinates and an even, continuous, real, integrable
$w$ equal to $1$ on $[-2T,2T]$, the sum $\sum_{\gamma\in Z}\widehat w(t-\gamma)$ takes negative
values \cite[Theorem~6.5]{P46}. So any zero-side argument must use infinitely many zeros. Three
further statements are expected but not proved. Fixed-support families of tests should be blind to
an off-line displacement $\delta$ of a zero at second order, the effect being $O((\delta L)^2)$,
where $L$ is the length of the support in the logarithmic variable. A lower bound
$\lam\ge-Ce^{\kappa T}$ for all $T$ should force the supremum $\Theta$ of real parts of zeros to
satisfy $\Theta-\frac12=O(\kappa)$, the case $\kappa=0$ being equivalent to the Riemann hypothesis
in the manner of von Koch's theorem. And a quartet off the line should weigh, in the even block,
like two zeros on the line, its displacement entering only at second order.

A simple logical filter applies to all such readings. Any argument whose every step remains valid
for a configuration of zeros containing an off-line quartet cannot prove the window problem for
every $T$, or for all sufficiently large $T$, since by Proposition~\ref{w6:prop:detect} that
configuration violates it. The filter does not apply to an argument for a single fixed window,
since detection requires large support. The Davenport--Heilbronn function and the linear
combinations of Euler products of Bombieri and Hejhal \cite{BombieriHejhal1995} are the standard
test models: they share the functional equation and, in the second case, conditionally, have
almost all zeros on the line, while having zeros off it.

Thus detectors carry information on the location of an off-line zero but do not force its
absence. On the zero side, the sign of the window problem is the statement that the indefinite
quartet blocks are absent; reading it from properties of the zeros is the problem restated, and the
uniform passage from a localized form to the complete form is equivalent to the window problem.
Some of the negative conclusions above rest on heuristics that were never tested numerically, and
we do not regard them as obstructions. Two questions were not pursued: uniform local sparsity of
deep off-line zeros, and the cone in the space of test functions cut out by the zero datum.

\section{Arithmetic approaches: prime sums, Euler products and Beurling systems}\label{w7:sec:primeside}

The form $\aT=\mT-\pT$ is the difference of an archimedean form, which is explicit, and a prime form
$\pT(F)=\sum_{n\in\NT}v_n\Rel\cF(\log n)$, which involves the finitely many values $\Lambda(n)$, $n<e^{2T}$. The
most direct strategy is therefore to bound $\pT$ from above by the tools of analytic number theory. We describe
three groups of attempts of this kind: the use of classical estimates for prime sums, the reformulation of the
sign as a one-sided inequality for a weighted prime sum, and the attempts to use the Euler product, the
functional equation and the integrality of $\N$ jointly. None of them proves a sign in a window beyond those
certified directly (Section~\ref{w2:sec:certified}). In the first two groups the reason is the same: the
available estimates control the size of a prime sum, while the window problem asks for an inequality of one
sign, exact to the size of the margin.

\subsection{Classical estimates for prime sums}\label{w7:fam:R17}

The idea is to bound the prime form by estimates for $\psi(y)-y$: the prime number theorem with remainder,
zero-free regions of Vinogradov--Korobov type \cite{Ford2002}, zero-density estimates, upper-bound sieves and
the large sieve \cite{MontgomeryVaughan1973}, the Montgomery--Vaughan mean-value theorem
\cite{MontgomeryVaughan1974}, Vaughan's identity and the van der Corput $B$-process \cite{Vaughan1977}, moments
of $\zeta$ and mollifiers, Tauberian theorems, and the passage from averaged to individual statements
(Bohr--Voronin universality, linear independence of the $\log p$). Each tool controls the size of a prime sum,
and the hope is that the control is fine enough to dominate $\pT$ by $\mT$ on $\PT$. The following proposition
excludes every estimate that is insensitive to a change of one coefficient.

\begin{proposition}[asymptotic estimates decide no window beyond $\frac12\log2$]\label{w7:prop:finmod}
Let $\mathcal S$ be a set of conditions on nonnegative weights $a=(a(n))_{n\ge2}$ supported on prime powers, and
suppose that $\mathcal S$ is preserved when $a(2)$ is replaced by $a(2)+M$ for any $M>0$. If $\Lambda$ satisfies
$\mathcal S$, then for every $T>\frac12\log2$ there is a weight satisfying $\mathcal S$ whose form has
$\lambda(T;a)<0$. Hence no argument whose only arithmetic input is $\mathcal S$ proves $\lam\ge0$ in any window
$T>\frac12\log2$.
\end{proposition}

\begin{proof}
For $T>\frac12\log2$ we have $2\in\NT$. Lemma~\ref{w3:lem:isolate} with $n_0=2$, $s=1$ gives a real
$F\in\PT$, $F\neq0$, with $\pT[a](F)=\frac12 v_2[a]\|F\|_2^2$ for every weight $a$, where
$v_2[a]=2a(2)/\sqrt2$. Put $a=\Lambda+M\delta_2$. Then
$\aT[a](F)=\mT(F)-\frac12\sqrt2\,(\log2+M)\|F\|_2^2<0$ once $M$ is large, and $a$ satisfies $\mathcal S$.
\end{proof}

The hypothesis holds for every statement that fixes $\sum_{n\le x}a(n)$, or a weighted variant of it, only up
to an error term that absorbs a bounded change of the sum: a prime number theorem with any remainder $o(x)$
and unspecified implied constant, or a pole-free region or a density estimate for the poles of
$\sum a(n)n^{-s}$, since the change alters the partial sums by $O(M)$ and the Dirichlet series by the entire
function $M2^{-s}$. It does not hold for asymptotics with an exact constant term, such as
$\sum_{n\le x}a(n)/n=\log x-\gamma+o(1)$ with $\gamma$ Euler's constant, whose constant term the change shifts
by $M/2$, nor for estimates with explicit constants such as $\psi(x)\le c\,x$ for all $x$ with a numerical $c$.
The proposition is the case $q=2$ of Theorem~\ref{w3:bar:B11}; classes of conditions that contain all
rescalings $c\Lambda$ or all large truncations are excluded in the same way by Theorem~\ref{w3:bar:B1}. For
estimates with explicit constants or exact secondary terms we have no theorem. Rough computations suggest that
the classical estimates of this kind decide no window $T>0.75$, and none beyond $0.749$ when upper sieves are
included, and it is plausible that the threshold $\frac12\log2$ of the proposition persists for this wider
class; neither statement is proved. The value $0.75$ lies just above $T_0(0)=0.74921$ (a numerical value), the
window beyond which the archimedean count alone forces a negative direction on $\PT$
(Section~\ref{w2:sec:structure}).

The individual tools fail for reasons that can be stated separately, although only some of them are theorems.
Positivity of the weights, support on prime powers, upper-bound sieves and the large sieve together do not
suffice, because the twisted weights $b_\vartheta$ of Section~\ref{w3:sec:falsifiers} satisfy
$0\le b_\vartheta\le\Lambda$ and every premise of an upper sieve, while they reproduce oscillations of size
$x^{1/2}$ and are falsifiers in the sense of that section; the Montgomery--Vaughan inequality gives the same
exponent $\frac12$ for every weight, and so cannot separate $\Lambda$ from $b_\vartheta$. Bounds for
exponential sums of Weyl or Vinogradov--Korobov type, used to control the prime form on the window, lead in
our estimates to error terms of size at least $2T\exp\big(e^{T}(1+o(1))\big)$, far above the scale of the
margin (Section~\ref{w2:sec:margin}). Vaughan's identity and the $B$-process, applied to $\pT$, return the
explicit formula itself after one step. For mollified moments there is a barrier at length exponent
$\theta\ge\frac12$, and in the moment expansions of $|\zeta|^{2k}$ that we examined (including the
fourth-moment formula of Motohashi type) every constant of order at least two enters the window form with net
coefficient zero; moreover the results on the proportion of zeros on the line hold, in the forms we examined,
also for the Davenport--Heilbronn and Epstein functions, and so cannot carry a sign. These last statements are
heuristic. So are the following two. The mean square of the window error obeys an identity of
Montgomery--Vaughan type that holds in every model with an explicit formula, with ceiling $T/\log T$, and the
nonnegativity of Montgomery's pair-correlation function \cite{Montgomery1973} gives an inequality of the
wrong orientation. Universality (Bagchi--Voronin), the linear independence of $\{\log p\}$ over $\Q$ and local
Tauberian theorems give density statements on the phase torus, not statements on the inertia of $\aT$ on
$\PT$; a Kronecker-type alignment of the phases $\gamma\log p$, $p\le Y$, requires heights
$\log\gamma\gtrsim Y/\log Y$, so an alignment argument cannot be run without information on the zeros.
Finally, the heuristic that multiplicativity suppresses resonances in truncated Euler products is contradicted
by numerical experiments (uncertified, and not reproduced here) with the Liouville function.

Two features are common to all these cases. The unconditional tools bound magnitudes: sieves and mean-value
theorems give upper bounds for nonnegative sums, or two-sided bounds with an error, whereas positivity of $\aT$
needs an upper bound for the signed quantity $\pT(F)$ that is exact to the size of the margin
(Observation~\ref{w3:bar:B12}), and the error of the available bounds exceeds the margin by many orders of
magnitude. Further, an estimate in absolute value destroys the signed cancellation of exponent $\frac12$ among
the roughly $e^{2T}$ terms of $\pT$ on which the sign depends. What would be needed is an evaluation of the
prime form in each window finer than the Montgomery--Vaughan inequality, namely a signed cancellation of
exponent $\frac12$ correlated with the samples of the test function. None of the estimates above provides it,
and we do not know whether one exists. Method barriers for mollified moments and classical inputs are discussed
heuristically in \cite{P33,P34}; \cite{P21} contains the Selberg--Delange asymptotics for sums of
$q^{\omega(n)}$, without a link to the zeros. For the classical background we refer to
\cite{Ingham1932,Titchmarsh1986}.

\subsection{One-sided arithmetic criteria}\label{w7:fam:R18}

By the explicit formula and Landau's theorem on Dirichlet integrals with nonnegative coefficients, the Riemann
hypothesis is equivalent to many eventual one-sided bounds for weighted prime sums, since a zero with real
part $\frac12+\Delta$, $\Delta>0$, forces oscillations of size $x^{1/2+\Delta-\varepsilon}$ of both signs.
Through Weil's criterion (Section~\ref{w1:sec:setting}) such statements bear on the window problem. One takes
a criterion of this type and tries to prove the one-sided bound by localizing, smoothing or regrouping the
prime sum and estimating each piece by known results. The prototype is the following classical statement; see
\cite[Ch.~V]{Ingham1932} and \cite[Ch.~XIV]{Titchmarsh1986} for criteria of this kind.

\begin{proposition}[a criterion of Landau type]\label{w7:prop:landau}
Let $S(x)=\sum_{p\le x}\log p\,\log(x/p)$. The Riemann hypothesis holds if and only if $S(x)\le x$ for all
sufficiently large $x$.
\end{proposition}

\begin{proof}
The Riesz mean $\sum_{n\le x}\Lambda(n)\log(x/n)=\int_1^x\psi(t)\,dt/t$ equals
$x-\sum_\rho x^\rho/\rho^2+O(\log x)$, the sum over zeros being absolutely convergent
\cite{Ingham1932}. The prime squares contribute
$\sum_{p\le\sqrt x}\log p\,\log(x/p^2)=2\int_1^{\sqrt x}\theta(t)\,dt/t=2\sqrt x\,(1+o(1))$ by the prime number
theorem, and higher powers $O(x^{1/3}\log^2x)$. Hence
\[
  S(x)=x-2\sqrt x\,(1+o(1))-\sum_\rho\frac{x^\rho}{\rho^2}.
\]
Under the Riemann hypothesis $|\sum_\rho x^\rho/\rho^2|\le\sqrt x\sum_\rho|\rho|^{-2}$, and
$\sum_\rho|\rho|^{-2}=\sum_\rho(\frac14+\gamma^2)^{-1}=2+\gamma_E-\log4\pi=0.0461\ldots<2$, so $S(x)<x$ for large $x$.
Conversely, suppose $S(x)\le x$ for $x\ge x_0$ and $\Theta=\sup\Rel\rho>\frac12$. Then
$g(x)=x+x^{1/2+\delta}-S(x)$, with $0<\delta<\Theta-\frac12$, is eventually nonnegative, and its Mellin transform
$\int_1^\infty g(x)x^{-s-1}dx$ (which equals $\frac1{s-1}+\frac1{s-\frac12-\delta}-\frac1{s^2}\sum_p\log p\,p^{-s}$
for $\Rel s>1$) continues meromorphically to $\Rel s>\frac12+\delta$, with poles at the zeros $\rho$ with
$\Rel\rho>\frac12+\delta$ and none on the real segment $(\frac12+\delta,\infty)$; this contradicts Landau's theorem
\cite[Ch.~V]{Ingham1932}. Hence $S(x)>x$ for arbitrarily large $x$.
\end{proof}

Many criteria of this type can be written down by the argument of Proposition~\ref{w7:prop:landau} or a
variant of it; each is equivalent to the Riemann hypothesis or implies it. We describe the main shapes that
were considered and the point at which the attempt to prove the one-sided bound stopped; in no case was the
inequality proved, and the obstructions stated are in part heuristic. For Riesz means with sparse sampling
one asks that the smoothed sums $\sum_{n\le x}(1-(n/x)^k)\Lambda(n)n^{-s_0}$ be bounded on a mesh $x=j^m$; for
nonnegative weights the mesh exponent is at most $4$, and what can be established covers finite ranges only.
Divisorial minorants are positive combinations $\sum_mf(mt)$ that minorize the prime weight; the optimal one
loses a proportion $\frac23+o(1)$ of the main term, and minorants supported on composite numbers need length
$\gtrsim x^{9/4}$. Bilinear reductions of type~II (pairs of factors near $\sqrt x$, rows indexed by the
cofactor, symmetric M\"obius forms) leave a residual signed sum that must be bounded at root precision, and
the passage to moduli loses; criteria on local correlations of primes, with tents and ramps in $\log(p/q)$,
reduce to the same bilinear problem. Exponent detectors measure the growth rate of the negative part of a
sparse family (Gaussian averages, or one row of a Hankel matrix together with its diagonal): the detection of
an exponent $>\frac12$ is exact, but the upper bound on the family that would exclude it is equivalent to the
Riemann hypothesis. For weighted barycentres of consecutive primes and for binomial carries, the versions with
positive carries lose a term of order $x$, and the signed versions are again equivalent to the Riemann
hypothesis. For criteria on the M\"obius, Liouville and Mertens functions (increments in bands of ratio tending
to $1$, pair energies) the available savings, which come from the theorem of Matom\"aki and Radziwi\l\l\
\cite{MatomakiRadziwill2016}, are sub-exponential, and there are non-multiplicative models with the same local
statistics. In the moving sieve one considers a M\"obius pulse on the integers free of prime factors
$\le z=(\log x)^2$, with a terminal prime; the zero-free region of \cite{Ford2002} appears to give the
one-sided bound for $\log z\ge C(\log\log x)^{5/3}(\log\log\log x)^{1/3}$, whereas the criterion needs
$\log z=2\log\log x$, a cut beyond the reach of any region of Vinogradov--Korobov type. Lastly, for a fixed
compactly supported Gaussian detector with square bias $-\frac18$ and ceiling $-\frac1{16}$, a residual balance
remains at constant scale.

Several of these inequalities were checked numerically on long finite ranges with interval arithmetic, in
computations that we do not reproduce here: a smoothed remainder of Landau type is negative for all real
$2\le x\le999983$, a barycentre inequality holds on $3\le x\le99991$, and a prime-indexed ramp inequality
holds for every prime $p\le10^{12}$. A verification on a finite range never implies an eventual bound.
Moreover, these criteria are statements about the whole set of zeros; they are not statements about a window,
and a certified range of $x$ does not certify any window (Section~\ref{w2:sec:certified}).

The common reason for the failure is the following. However the sum is localized or regrouped, the reduction
leaves one signed piece at root precision whose bound is itself equivalent to the Riemann hypothesis; the
pieces estimated in absolute value lose a power of $x$, and the available savings (zero-free regions,
cancellation of the M\"obius function in short intervals \cite{MatomakiRadziwill2016}) are sub-exponential
where an exponential saving is needed. Heuristically, models with positive weights reproduce any finite
statistics of a criterion, and non-multiplicative models with the correct local statistics defeat any argument
that does not use multiplicativity. A proof along these lines needs an unconditional one-sided arithmetic
estimate at root scale, which is of a different order from an improvement of constants in the known
estimates. The sign was eventually obtained along this line, by a one-sided criterion for Dirichlet
$L$-functions weighted by $1-\chi_K$, \cite[Theorem~12.9]{P46}, with the required bound coming from the
arithmetic of automorphic forms rather than from estimates of prime sums; see Section~\ref{w8:sec:open}.

\subsection{Euler product, functional equation and integrality}\label{w7:sec:euler}
\phantomsection\label{w7:fam:R19}

The falsifiers of Section~\ref{w3:sec:falsifiers} (see Table~\ref{w3:tab:falsifiers}) indicate that three
properties of $\zeta$ are each indispensable:
\begin{itemize}
\item[(E)] the Euler product over all primes (equivalently, the support and exact values of $\Lambda$);
\item[(F)] the functional equation with gamma factor $\pi^{-s/2}\Gamma(s/2)$ and conductor $1$;
\item[(I)] integrality: the integers form the additive lattice $\Z_{\ge1}$, $N(x)=\lfloor x\rfloor$.
\end{itemize}
The Davenport--Heilbronn function and the Epstein function of $x^2+5y^2$ satisfy a functional equation of
Riemann type (with conductor $5$, respectively degree two and conductor $20$) but not (E); twisted, truncated
and edited weights satisfy parts of (E) and not (F); discrete Beurling systems have an Euler product over
their generalized primes, but in general no functional equation, and they do not satisfy (I). One therefore
looks for an identity or an inequality in which (E), (F) and (I) enter jointly and which implies a lower bound
for $\aT$ on $\PT$. The candidates were the convolution identity $\Lambda\ast1=\log$ and its consequences,
Selberg's identity, discriminants of multiplicativity, cones on $\log\zeta$ or $-\zeta'/\zeta$, the
classification of the Selberg class \cite{Selberg1992}, product formulas, and two-variable explicit formulas.
Two elementary facts frame the discussion.

\begin{proposition}[locality of the window form]\label{w7:prop:local}
The form $\aT$ depends on the weight only through $\Lambda(n)$, $n<e^{2T}$. In particular, a prime
$p\ge e^{2T}$, or any modification of the Euler factors at such primes, does not change $\aT$ or $\lam$.
\end{proposition}

\begin{proof}
$\pT(F)=\sum_{n\in\NT}v_n\Rel\cF(\log n)$ and $\supp\cF\subset[-2T,2T]$.
\end{proof}

So the Euler product over all primes can enter a proof only through an argument that couples infinitely many
windows; within one window it is the finite list $\Lambda(n)$, $n<e^{2T}$, and Theorem~\ref{w3:bar:B11} shows
that no finite list of values suffices by itself.

\begin{proposition}[Selberg's identity and its rigidity]\label{w7:prop:selberg}
For an arithmetic function $a$, write $L(n)=\log n$. Then
\begin{equation}\label{w7:eq:selberg}
  a\cdot L+a\ast a=\mu\ast L^2
\end{equation}
holds if and only if $a=\Lambda$.
\end{proposition}

\begin{proof}
For $a=\Lambda$: in terms of Dirichlet series, $\Lambda\leftrightarrow-\zeta'/\zeta$, $\Lambda L\leftrightarrow
(\zeta'/\zeta)'=\zeta''/\zeta-(\zeta'/\zeta)^2$, $\Lambda\ast\Lambda\leftrightarrow(\zeta'/\zeta)^2$ and
$\mu\ast L^2\leftrightarrow\zeta''/\zeta$, so \eqref{w7:eq:selberg} holds. Conversely, evaluate
\eqref{w7:eq:selberg} at $n=1$: $a(1)^2=(\mu\ast L^2)(1)=0$, so $a(1)=0$. For $n\ge2$ the identity reads
$a(n)\log n+\sum_{d\mid n,\,1<d<n}a(d)a(n/d)=(\mu\ast L^2)(n)$, which determines $a(n)$ from the values at
proper divisors since $\log n\ne0$. By induction there is exactly one solution, and it is $\Lambda$.
\end{proof}

The relation of degree one $\Lambda\ast1=\log$ also characterizes $\Lambda$, by M\"obius inversion. Inserted
linearly in the window form, however, it appears to be equivalent to enlarging the window,
$T\mapsto T+\frac12\log X$, and the squares of ``sum$\times$product'' type built from it contain the truncated
Weil square and are invariant under a change of sign of the archimedean part, hence blind to the sign. More
generally, an identity that is homogeneous linear in the weight holds for $c\Lambda$ as well and carries no
sign (Theorem~\ref{w3:bar:B1}(ii)). Identity \eqref{w7:eq:selberg} is, as far as we can see, the first exact
relation quadratic in $\Lambda$ that characterizes $\Lambda$. In the two-variable explicit formula on thickened
diagonals the primes$\times$primes term is
$\sum_N g(N)\sum_{d\mid N}\Lambda(d)\Lambda(N/d)\varphi(d^2/N)$ plus explicit finite terms
\cite[Theorem~5.4]{P28}, so that the convolution $\Lambda\ast\Lambda$ enters with the sign of a square; in the
two-variable form built from it we were nevertheless unable to extract a sign from the prime block, and we
expect that it cannot be extracted from that block alone.

The functional equation at conductor one is rigid. By Hamburger's theorem \cite{Hamburger1921}, if
$f(s)=\sum a_nn^{-s}$ and $g(s)=\sum b_nn^{-s}$ converge absolutely for $\Rel s>1$, $P(s)f(s)$ is entire of
finite order for some polynomial $P$, and
$\pi^{-s/2}\Gamma(\frac s2)f(s)=\pi^{-(1-s)/2}\Gamma(\frac{1-s}2)g(1-s)$, then $f=g=c\,\zeta$. Thus (F) in the
exact form stated already forces $f=c\zeta$ among Dirichlet series over $\N$, and (E) and (I) enter such a
statement only as identification. The theorem says nothing about weights $a\ne\Lambda$ in the window form or
about series over other frequency sets, and it does not by itself give any bound on $\aT$. In the Selberg
class the elements of degree one are $\zeta$ and the shifts $L(s+i\theta,\chi)$ with $\chi$ a primitive
non-principal Dirichlet character and $\theta\in\R$ \cite{KaczorowskiPerelli1999}, so $\zeta$ is isolated
among elements of degree one and conductor one. Multiplicativity can also be used as a discriminant: an Euler
product forces the anatomy of a localized explicit-formula form to vanish off prime powers
\cite[Theorem~3.1]{P10}, and in degree one, under a non-vanishing hypothesis and the axioms of the Selberg
class, the anatomy identifies $\zeta$ or the $L$-function of a real primitive character
\cite[Theorem~10.1]{P10}. A localized numerical discriminant built on this idea does not separate $\zeta$ from
smooth non-multiplicative controls: in an uncertified computation, not reproduced here, its value $1.9798$ for
$\zeta$ lies below the values $2.4391$, $2.1118$, $2.2085$ of three smooth controls, so that the separating
feature is smoothness and not multiplicativity. We found no cone condition on $\log\zeta$, or on
$-\zeta'/\zeta$, whose verification gives $\lam\ge0$; the scaled weights $c\Lambda$ satisfy the same cone
conditions, for $c$ near $1$ under a non-degeneracy hypothesis that we have not removed. As for positivity
place by place, each local form at a prime is indefinite, the coupling of the places through a product formula
either fixes $\Lambda$ or admits a negative comb, and we know of no factorization of the form by product
kernels on the solenoid; see Section~\ref{w4:sec:local}.

In every chain of this kind that we examined, one of the inputs is used through a prior linear identity, after
which it is lost. Poisson summation and the functional equation of $\theta$ do not involve $\Lambda$ at all:
they use only the additive lattice $\Z$, so they cannot distinguish $\Lambda$ from any other weight on $\N$. A
joint use of the three inputs seems to require a step that is non-linear in $\Lambda$, such as
\eqref{w7:eq:selberg}, or one of the inputs used as an inequality. We do not know of a single step that uses
(E), (F) and (I) jointly and fails for every falsifier, and we do not know of an obstruction to its existence;
a first test for a candidate is a prime-side identity that fails for each of the falsifiers of
Table~\ref{w3:tab:falsifiers}. We return to this requirement in Section~\ref{w8:sec:constraints}. Related
material is in \cite{P19} and \cite{DavenportHeilbronn1936}.

\phantomsection\label{w7:fam:R20}
We turn to integrality. To decide whether it is needed, one compares $\N$ with systems of Beurling generalized
integers whose counting function is $N(x)=cx+O(x^\theta)$. If every Beurling system with small $\theta$
satisfied the analogue of the Riemann hypothesis, a proof could use only the size of $N(x)-cx$; if not, the
exact lattice structure of $\N$ must enter. The literature answers this for $\theta>\frac12$. A Beurling
system can have $N(x)=cx+O(x^\theta)$ with $\frac12<\theta<1$ while its prime-counting function oscillates as
much as the de la Vall\'ee-Poussin zero-free region permits, the zeros of its zeta function approaching the
line $\Rel s=1$ \cite{DMV2006}; systems satisfying the analogue of the Riemann hypothesis and systems with
large oscillation are both constructed in \cite{Zhang2007}; see \cite{DiamondZhang2016} for the general theory
and \cite{BDR2023} for systems with power savings in both counting functions, where the two exponents cannot
both be below $\frac12$. In the window setting, the discrete system with primes $2,\frac{16}5,5,7,11,\dots$ is
a falsifier for $T\ge\frac12\log5+10^{-4}$, with the level of verification stated in
Section~\ref{w3:sec:falsifiers}.

Below $\frac12$ we have constructions that we expect to be correct but that are not written out with complete
proofs, and whose relation to the systems of \cite{DiamondZhang2016} we have not checked: a discrete system
with $N(x)=cx+O(x^{83/208})$ and infinitely many zeros on $\Rel s=\frac{11}{16}$, and, for each
$\theta\in(\frac13,\frac12)$, discrete systems with zeros on $\Rel s=\frac{1+\theta}2$. If they stand, an
argument whose only input on the integers is $N(x)=cx+O(x^\theta)$ with some $\theta>\frac13$ cannot prove the
sign, and must use integrality beyond the size of the remainder. We do not know whether a remainder exponent
$\theta\le\frac13$ suffices, that is, whether there are off-lattice discrete systems with $\theta\le\frac13$
violating the analogue of the Riemann hypothesis, nor whether there is rigidity at $\theta=0$. Three further
remarks point the same way. An analytic theorem of the form ``$\theta<\frac12$ implies \dots'' must use the
integer multiplicities of the primes, and sieve bounds hold in every discrete system with
$N(x)=Ax+O(x^\theta)$, hence cannot see integrality. Reductions of linear or factor type (``$N=\rho x+O(1)$
forces a lattice'', factor arguments for all $\theta<1$) are false or are contradicted by explicit discrete
models. And in every transfer between a Beurling system and $\N$ that we examined, the failure falls into one
of two classes: analytic self-duality about $\frac12$, or the exact additive geometry of $\N$; we do not have
this dichotomy as a general statement.

The reason these comparisons do not lead to the sign is that the integer count does not see the window form:
$N(x)$ is a cumulative quantity, while $\aT$ depends on the exact positions $\log n$ of the first $e^{2T}$
prime powers and on their weights. Beurling systems that agree with $\N$ in the size of every remainder term
reproduce the analytic side of every argument we examined, and the separation occurs only at the passage from
the multiplicative semigroup to the additive lattice $\Z$.

Finally, several objects of automorphic origin were examined as possible sources of the sign: the positivity of
the Petersson norm and of Rankin--Selberg integrals, the Maass--Selberg relations, metaplectic Eisenstein series
of weight $\frac12$, the character of the Weil representation, the class number relations of
Kronecker--Hurwitz and Eichler--Selberg type, and the summation formulas of Voronoi and Kuznetsov. None of
them gave it, and the reasons repeat those met above. The positivity of a Petersson norm or of a
Rankin--Selberg integral is a statement about an $L$-function of higher degree: for a cusp form $f$ the form
attached to $L(s,f\times\bar f)$ is the sum of the form of $\zeta$ and that of the symmetric square, and to
subtract the latter is to return to the problem. The identities coming from the Weil representation reproduce
the local trace formula and hold for every weight. In the class number relations the function $\Lambda$ appears
only through a logarithmic derivative, and the positivity present before that operation does not survive it.
The Voronoi formula holds equally for Epstein zeta functions of class number greater than one, whose forms are
negative on some windows (Table~\ref{w3:tab:falsifiers}). These conclusions are heuristic, in the sense that
each concerns the constructions that were actually tried. What they have in common is that the automorphic
input was used as a source of analytic positivity. The input of \cite{P46} is automorphic too, but it is of
another nature: it is the integrality of algebraic modular forms and of their theta lifts, an arithmetic
property that none of these constructions used.

\section{Geometric approaches and deformations}\label{w7:sec:geometry}

The approaches described so far work with the form $\aT$ in a fixed window. This section describes attempts to
transport to $\zeta$ the mechanism of Weil's proof for curves over finite fields, or to realize the zeros on a
global space where positivity would be structural, and then attempts that move the weight sequence, the window,
or a map between windows, and try to carry positivity along.

For a curve $C$ of genus $g$ over $\mathbb F_q$, Weil proved the Riemann hypothesis by intersection theory on
$C\times C$ \cite{Weil48}: the graphs $\Gamma_n$ of the powers of Frobenius have intersection numbers
$\Gamma_n\cdot\Delta=N_n$ with the diagonal, the Lefschetz identity turns these into the explicit formula, and
the Hodge index theorem (in the form of the Castelnuovo--Severi inequality) gives a sign on the classes
orthogonal to the two rulings, which is the Riemann hypothesis. Grothendieck's proof of the Hodge index
inequality for surfaces from the Riemann--Roch theorem \cite{Grothendieck1958} gives this sign in every
characteristic, and his standard conjectures \cite{Grothendieck1969} isolate the positivity statement needed in
higher dimension.

\subsection{Intersection frames and the Hodge index theorem}\label{w7:fam:R21}

The idea is to equip the bilinear form of the explicit formula with the structure of Weil's proof: a degree, an
effective cone, a radical, an ample class, an affine duality $E\mapsto K-E$ and a Riemann--Roch functional
$h^0$; and then to run Hartshorne's proof of the Hodge index theorem \cite[Ch.~V, Theorem~1.9]{Hartshorne1977}.
Variants look for arithmetic analogues of the four ingredients (surface, Frobenius correspondences, Lefschetz
identity, Hodge index), for hard-Lefschetz structures, tensor amplification, Frobenius lattices, and tropical or
Arakelov versions; for Hodge theory outside geometry the model is \cite{AdiprasitoHuhKatz2018}.

In \cite{P47} the frame is built on a class $\mathcal D^W$ of exponentially decaying test functions, with the
extended Weil form $\mathcal I$, the sign being normalized so that the Weil inequality reads
$\mathcal I(F,F)\le0$ on primitive tests. The P\'olya class $K$, with Mellin transform $\Phi_K=\xi$, is positive
and of degree one, and the space splits orthogonally into a hyperbolic plane and the primitive sector. The
results that bear on the sign are the following two theorems.

\begin{theorem}[Hodge-index transfer {\cite[Theorem~6.1]{P47}}]\label{w7:thm:transfer}
If $h^0:\mathcal D^W\to\R_{\ge0}$ satisfies (a) $h^0\ge0$; (b$'$) $h^0(E)>0$ implies that $E$ is equivalent
modulo the radical to a nonzero effective class; (c) $h^0(E)+h^0(K-E)\ge\frac12\mathcal I(E,E)-\frac12\mathcal
I(E,K)+\chi$ for some constant $\chi$ and all $E$; then $\mathcal I(F,F)\le0$ for every real primitive
$F\in C_c^\infty(\R)$.
\end{theorem}

\begin{theorem}[frame defect {\cite[Theorems~8.6 and 8.8]{P47}}]\label{w7:thm:defect}
No function $h^0$ satisfies (a), (b$'$), (c), for any $\chi\in\R$. The obstruction is a cosh inequality on the
effective cone: for an effective class $E'$ and a zero $\frac12+i\gamma$,
$|\Rel(e^{i\theta}\Phi_{E'}(\frac12+i\gamma))|\le\frac12 s(\gamma,\theta)\deg E'$ with
$s(\gamma,\theta)=\sup_u|\cos(\gamma u+\theta)|/\cosh(u/2)$, and $s_1=s(\gamma_1,\frac\pi2)=0.99846\ldots<1$ at
the first zero produces classes of positive self-intersection (of degree $\ge325$ when $\chi=0$) with no
effective representative.
\end{theorem}

The proof of Theorem~\ref{w7:thm:defect} in \cite{P47} depends on numerical inputs: the numerically known
simplicity and isolation of the first zero (no other zero with $0<\Iml\rho\le21$) and Trudgian's bound for
$S(T)$; the constants are computed in 40-digit floating-point arithmetic, not interval arithmetic. The same
defect is stated for $\Xi_a(s)=\xi(s+a)\xi(s-a)$, $0<a<\frac12$ \cite[Theorem~9.4]{P47}; the proof there relies on
an adaptation of \cite[Theorem~6.2]{P45} that is indicated but not written out in full. It also holds for curves
over finite fields, where the Riemann hypothesis is true: for $y^2+xy=x^3+1$ over $\mathbb F_2$ the sharp
constant is $\sqrt{7/11}$ (exact), with a witness in degree $3$ \cite[Proposition~10.7, Theorem~10.8]{P47}. The
constants $s_1$ and $\sqrt{7/11}$ are the same invariant, computed for $\zeta$ and for the curve. Thus the
defect is a property of the effective cone of the frame and does not depend on the zeros. The transfer theorem
remains a valid implication, with a hypothesis that cannot be satisfied.

Other papers of the series contribute the following. The arithmetic Lefschetz programme of \cite{P42} proposes
objects over $\Z$ for the four ingredients; it shows that infinite-rank contact excludes a finite-rank geometric
approach \cite[Theorem~9.1]{P42}, and that the symmetrized logarithmic derivative
$\zeta'/\zeta(s)+\zeta'/\zeta(1-s)$ is an elementary gamma expression, so that functionals built from it cannot
see the zeros \cite[Proposition~1.1]{P42}. On semilocal id\`ele class groups one has a correspondence algebra
with trace \cite{P43}, a Frobenius family with contact of weight $\Lambda(n)/\sqrt n$ with the diagonal
\cite{P44}, and the explicit formula for the self-intersection of a radial class, with a real primitive compactly
supported test detecting any off-line zero \cite[Theorem~6.2]{P45}; the sign itself, the fourth ingredient, is
not provided. An $\mathfrak{sl}_2$-string of length $\ge3$ graded by the ordinates of the zeros forces a
three-term arithmetic progression among them \cite[Proposition~5.1]{P17}, which is excluded under the hypothesis
that the ordinates are linearly independent over $\Q$. Tensor powers of a form of signature $(p,q)$ have negative
index $((p+q)^k-(p-q)^k)/2$ \cite[Theorem~4.1, Proposition~4.3]{P28}, so tensor amplification multiplies a
negative index and does not remove it.

A further obstruction, to Frobenius structures defined over $\overline{\Q}$, is a direct consequence of a
classical theorem.

\begin{proposition}[transcendence of Frobenius exponents]\label{w7:prop:rational}
Let $\rho$ be a non-real complex number and $p_1,p_2,p_3$ distinct primes. Then at least one of
$p_1^{\rho},p_2^{\rho},p_3^{\rho}$ is transcendental. In particular no non-trivial zero of $\zeta$ is a common
Frobenius exponent, in a structure defined over $\overline{\Q}$, at three primes.
\end{proposition}

\begin{proof}
We apply the six exponentials theorem \cite{Lang1966} with $x_1=1$, $x_2=\rho$, which are linearly independent
over $\Q$ since $\rho\notin\R$, and $y_j=\log p_j$, which are linearly independent over $\Q$ by unique
factorization. The numbers $e^{x_1y_j}=p_j$ are algebraic, so some $e^{x_2y_j}=p_j^\rho$ is transcendental.
\end{proof}

Several further variants were examined and abandoned; we state the reasons without proof, and nothing in this
paper depends on them. A pairing on the prime-power coefficients does not see the spectrum, whatever the kernel.
For the lattice count $h^0_\Z(mK)$ we only know a lower bound of order $(\log m)^{3/2}/\log\log m$ and no upper
bound, not even its finiteness, while clause (c) requires quadratic growth; the comparison is unsettled. No
primitive class has all its Frobenius traces nonnegative, as one expects from Landau's theorem. An arithmetic
transplant of the Castelnuovo--Severi argument applies an inequality valid for classes of positive square to
classes of negative square. A holomorphic lamination with a single pole carries at most one pair of zeros, so a
carrier would have to fail a Torelli property. The Galois group ring yields a positive structure that is inert
for $\zeta$, its effect in function fields coming from the discreteness of $(\log q)\Z$. A frame on the tensor
square, and six weakenings of clause (b$'$), tested along explicit rays $tF+K$, also fail.

In every frame examined, the positivity is postulated in one of the axioms (effectivity, a Riemann--Roch
inequality, an ample class), and verifying the axiom is equivalent to the sign; the Frobenius and Galois
structures available over $\Z$ act far from the zeros, or as translations. Outside the class of functions $h^0$
of Theorem~\ref{w7:thm:defect} we know no general obstruction. One would need an object on the ad\`elic square
that does not pass through an $h^0$ of this type: for example a non-radial regularization of the trace on the
square combining degree and zero data, a non-invertible correspondence of degree $q$, or a $\Z$-polarized
object with an ample class and quadratic $h^0$. None has been constructed.

\subsection{Moment extraction and global spaces}\label{w7:fam:R22}

Deligne's proof of the Weil conjectures \cite{Deligne1974} bounds individual Frobenius eigenvalues by bounding
high moments of traces over a family and letting the order of the moment tend to infinity. The analogue we
considered is an additive pencil over Farey arcs: one takes Gaussian moments of order $2k$ of exponential sums
over the arcs, and hopes to extract $\psi(N)-N$ with an individual bound. The elementary core of the difficulty
is the following.

\begin{observation}[no extraction from $L^2$ identities]\label{w7:obs:L2}
Let $X_f\ge0$ be indexed by a finite family with weights $w_f>0$, and let $M_k=\sum_fw_fX_f^k$. The power means
$(M_k/M_0)^{1/k}$ are nondecreasing in $k$ and tend to $\max_fX_f$. A Parseval identity fixes only $M_1$, and
from $M_1$ alone one obtains only $\max_fX_f\le M_1/\min_fw_f$ and $\max_fX_f\ge M_1/M_0$. An individual bound
therefore requires an independent upper bound on $M_k$ for large $k$ of the right size.
\end{observation}

\begin{proof}
Monotonicity of power means is H\"older's inequality; the limit is standard for finite families; the bounds
from $M_1$ are immediate from $X_f\ge0$.
\end{proof}

In Deligne's setting the bound on $M_k$ comes from cohomology (dimension times weight), and we see no counterpart
of it in the additive pencil. Our analysis of the pencil is not written up, and we report its conclusions without
proof. The extraction identity holds, with three corrections, but it holds uniformly for $\Xi_a$ as well, so it
carries no arithmetic, and a positive (Bochner) version is blind to the zeros. The boundedness of the moments of
the needed order reduces, for one range of a parameter $c\in[\frac12,1)$, to the large sieve, and is obtained
only where it gives nothing; it is open for $c<\frac12$, and in the form needed it is equivalent to a statement
of Hardy--Littlewood $k$-tuple type. Nonnegative families have abscissa of convergence at least
$2k\beta_{\max}$, so no nonnegative family yields the pole a priori, and the conversion of a power saving into
the sign is equivalent to a zero bound. A full transposition of Deligne's weight argument, with a cohomological
source of moment bounds, has not been attempted. A dictionary with Laumon's local Fourier transform
\cite{Laumon1987} was considered only formally; it appears to meet the obstruction of
Proposition~\ref{w7:prop:rational}.

\phantomsection\label{w7:fam:R23}
A different way to import geometry is to realize the zeros as a spectrum on a global space in which positivity
is structural: Connes's ad\`ele class space and trace formula \cite{Connes1999} with its archimedean positivity
\cite{ConnesConsani2021}, Deninger's dynamical systems on foliated spaces \cite{Deninger1998}, the Bost--Connes
system \cite{BostConnes1995}, condensed or liquid analytic structures, $K$-theory and $KK$-theory, compact
quantum groups. Connes reduced the Riemann hypothesis to the validity of a trace formula on the ad\`ele class
space \cite{Connes1999}; in this sense the ad\`ele class space is a reformulation of Weil's criterion and
provides no sign by itself, and we reach the same conclusion, heuristically, for Meyer's nuclear realization of
the cohomology \cite{Meyer2005}. Connes and Consani proved Weil positivity at the archimedean place for test
functions with support in a small interval \cite{ConnesConsani2021}.

For the other frameworks we report what we found, without proof. Deninger's axioms (cohomology, flow, Hodge-type
data) can be satisfied, and an object satisfying them has as its whole content a single positivity statement
equivalent to the problem. In each of eighteen published frameworks that we tested against this requirement, the
positivity is derived only in settings where $\zeta$ cannot live. The partition function of the Bost--Connes
system is $\zeta(\beta)$, and an inertia metric built from the system gives no sign. For condensed and liquid
structures we have negative statements for three specific classes of constructions only. In realizations by
compact quantum groups, Haar orthogonality appears to kill the interference terms, and invariants of $KK$- and
$L$-theoretic type meet a $\Q/\Z$ parity obstruction. Where a construction was carried out, the positivity it
provides is either equivalent to Weil's criterion or holds uniformly for functions without an Euler product; the
missing datum is the coupling of the archimedean and prime parts (see also \cite{P42}). We know no obstruction
to these frameworks in general: the GNS construction over the ad\`ele class algebra, a foliated space in
Deninger's sense, and a quadratic Riemann--Roch theorem over $\operatorname{Spec}\Z$ remain unconstructed.

\subsection{Deformation of the weight and asymptotic laws}\label{w8:sec:families}\label{w8:fam:R24}

We now turn to arguments that move the weight or the window. They all meet the same two facts: the form
$\aT[a]$ is affine in the weight (Proposition~\ref{w1:prop:convex}), and the margin is very small
(Section~\ref{w2:sec:margin}). Throughout, $\lambda(T;a)$ is the bottom of $\aT[a]$ on unit vectors of $\PT$, so
that $\lam=\lambda(T;\Lambda)$.

The first idea is to join $\Lambda$ to a reference weight $\rho$ for which positivity is known, for instance
$\rho=c\Lambda$ with $c$ small in a window where $\mT\ge0$ on $\PT$, or a smooth density; to follow the bottom
along the path $a_t=(1-t)\rho+t\Lambda$, or along $c\mapsto c\Lambda$; and to show that it cannot cross zero
before $t=1$. A variant works with $T$ instead of $t$: one looks for a two-sided asymptotic law
$\lam=L(T)(1+o(1))$ with an explicit positive $L(T)$, whose lower half would contain the sign. A second variant
controls the one-sided derivative of $T\mapsto\lam$ by $\lam$ itself, through a Gr\"onwall inequality
$-\frac{d}{dT}\lam\le C(T)\lam$, so that a positive value cannot reach zero. The first statement is elementary.

\begin{proposition}[the admissible interval of a scaling]\label{w8:prop:scaling}
Fix $T>0$ and a weight $a$. The function $c\mapsto\lambda(T;ca)$ is finite and concave on $\R$. Hence
$J_T(a)=\{c\in\R:\lambda(T;ca)\ge0\}$ is a closed interval (possibly empty or unbounded). If $\lambda(T;a)>0$, then
$J_T(a)$ contains a neighbourhood of $c=1$. If moreover $F\in\PT$, $\|F\|_2=1$, satisfies $\pT[a](F)>0$, then
$\sup J_T(a)\le 1+\aT[a](F)/\pT[a](F)$. More generally, for two weights $\rho,a$ the set
$\{t\in[0,1]:\lambda(T;(1-t)\rho+ta)\ge0\}$ is a closed interval, and if it contains $0$ and $1$ it is all of $[0,1]$.
\end{proposition}

\begin{proof}
For fixed $F\in\PT$ the map $c\mapsto\aT[ca](F)=\mT(F)-c\,\pT[a](F)$ is affine. So $c\mapsto\lambda(T;ca)$ is an
infimum of affine functions, hence concave. It is $<+\infty$ (evaluate at one $F$). It is $>-\infty$ by the lower bound
$\aT[ca](F)\ge C_0\|F\|_2^2-|c|\sum_{n\in\NT}2|a_n|n^{-1/2}\|F\|^2$, which uses the bound $\mT(F)\ge C_0\|F\|_2^2$ of Section~\ref{w1:sec:setting}, $|\cF(u)|\le\|F\|_2^2$ and
Definition~\ref{w1:def:lambda}. A finite concave function on $\R$ is continuous, so its superlevel set $\{\ge0\}$ is a
closed interval. Continuity at $c=1$ gives the neighbourhood when $\lambda(T;a)>0$. The upper bound is
Proposition~\ref{w2:prop:concave}(ii) applied to $a$:
$\lambda(T;ca)\le\aT[a](F)-(c-1)\pT[a](F)<0$ for $c-1>\aT[a](F)/\pT[a](F)$. The same argument applies to
$t\mapsto\lambda(T;(1-t)\rho+ta)$, which is concave in $t$.
\end{proof}

On every certified window (Table~\ref{w2:tab:certified}), $J_T(\Lambda)$ is therefore a closed interval around
$1$. For $T\le\frac12\log2$ one has $\pT=0$ and every rescaling is positive. For larger windows, taking $F$ close
to the minimizer, we expect the upper end of $J_T(\Lambda)$ to be within $O(\lam)$ of $1$, because numerically
$\pT(F)$ stays of order one on near-minimizers (Section~\ref{w2:sec:margin}); we have no proof that the interval
contains a neighbourhood of $1$ of that size. Thus positivity singles out $c=1$ only to the precision of the
margin, and a proof must explain a coincidence between $\mT$ and $\pT$ to that order.

The concavity in Proposition~\ref{w8:prop:scaling} has a consequence for homotopies. Along a segment from a
positive reference $\rho$ to $\Lambda$, positivity holds on an interval $[0,t_\infty(T)]$, and it holds at $t=1$
only if it already holds at $\Lambda$. Thus concavity transfers positivity between two positive endpoints and
never towards an unknown one. A homotopy argument therefore needs an input other than convexity to exclude a
crossing at some $t<1$, and a crossing at $t<1$ is exactly what the rescalings produce: every hypothesis that is
stable under small changes of the weight is satisfied by $(1+\delta)\Lambda$ for $\delta$ above the margin
(Proposition~\ref{w2:prop:concave}(ii), Remark~\ref{w2:obs:precision}). In particular no argument can
prove $\lambda(T;c\Lambda)\ge0$ for all $c$ with $|c-1|\le\delta$ when $\delta>\aT(F)/\pT(F)$ for some unit
vector $F\in\PT$ with $\pT(F)>0$. The attempts made were the following. For the path from a smooth density to
$\Lambda$, a heuristic argument gives a dichotomy $\lim_{T\to\infty}t_\infty(T)\in\{0,1\}$; families built from
wave packets reduce to it under an additional hypothesis, and two-copy constructions do not seem to change it. We
do not know whether the statement that $t_\infty>0$ in the limit is strictly weaker than the window problem. A
Neumann series of finite degree for the resolvent of the archimedean part reaches the margin, heuristically, only
with degree growing like $\log(1/\lam)$. A stability statement in the weight, of the form
$\lambda(T;a)\ge-C_T\,\Delta_T(a,\Lambda)$ for a distance $\Delta_T$ on weights that vanishes at $\Lambda$, is
open; our attempts at it failed because the rescalings $(1+\delta)\Lambda$ satisfy every hypothesis used, with
slack.

As for asymptotic laws, the upper half is accessible, since it is produced by test functions
(Section~\ref{w2:sec:margin}). The bound $\lam\le4\cdot10^{9}\,T\,e^{-15.449\,T}$ for $T\ge1$ of
\cite[Corollary~9.11]{P46} excludes a law $\lam\asymp e^{-cT}$ with $c<15.449$, and numerically (uncertified)
the decay is much faster: the logarithm of the bottom is of order $-c_1e^{2T}$ with $c_1$ close to $4\pi$, the
third digit of the fits of Section~\ref{w2:sec:margin} not being stable. A two-sided law with the main term of
Landau--Widom type of the prolate problem on $\IT$ is conjectured in \cite{Zhu2026}. The lower half carries the
sign. Since $\lam\to0$ in the regime of interest, a lower bound $\lam\ge L(T)(1-o(1))$ with $L(T)>0$ is the
window problem together with a quantitative margin, and we see no reason for it to be easier. By the soft-target
reduction (Proposition~\ref{w2:prop:soft}), already $\lam\ge-\varepsilon(T)$ with $\varepsilon(T)\to0$ suffices,
so no weakening in the size of the error is available. We note that Observation~\ref{w3:bar:B4b} excludes
functionals of $\{\lambda(T'):T'\le T\}$ on a bounded range of windows, and does not exclude two-sided asymptotic
laws. The Gr\"onwall inequality $-\frac{d}{dT}\lam\le C(T)\lam$ with a free constant is, heuristically, a restatement: it
is equivalent to positivity together with a local Lipschitz bound. A version with $C(T)=8\pi e^{2T}(1+o(1))$,
the rate suggested by the fits, has not been attempted.

One-parameter continuation arguments, with the classical theorems attributed, are surveyed heuristically in
\cite{P27}, which also states a coupled two-dimensional monotonicity problem with a falsifier
\cite[Problem~6.1]{P27}; the synthesis \cite{P29} restates the propagation trichotomy. The distinction between
eventual and uniform positivity in scale-dependent hierarchies is in \cite[Theorem~3.1]{P15}.

A related question concerns all admissible weights at once. Let $W_\infty$ be the set of sequences $w\ge0$ on
the integers $n\ge2$ such that $\aT[w]\ge0$ on $\PT$ for every $T$. It is convex
(Proposition~\ref{w1:prop:convex}) and contains $\Lambda$ by \cite[Corollary~16.3]{P46}. For $w\in W_\infty$ the
form is the pairing of the autocorrelation $\cF$ with an even distribution $u_w$ on $\R$, the sum of the
archimedean kernel and of negative atoms at the points $\pm\log n$, proportional to $w_n$. Since
every $F=(-\partial^2+\frac14)G$ with $G$ smooth of compact support is primitive and has
$\cF=(-\partial^2+\frac14)^2(G\star\widetilde G)$, the distribution $(-\partial^2+\frac14)^2u_w$ is of positive
type, and the Bochner--Schwartz theorem gives a positive even tempered measure $\nu_w$ and real numbers
$\alpha,\beta$ with
\[
u_w=\check\nu_w+\alpha\cosh(u/2)+\beta\,u\sinh(u/2).
\]
For $w=\Lambda$ this is the explicit formula, and $\nu_\Lambda$ is carried by the ordinates of the
zeros. We do not know a description of $W_\infty$. In particular we do not know whether $\Lambda$ is determined
in $W_\infty$ by arithmetic conditions on $\nu_w$, such as being purely atomic with integer masses; a rigidity
statement of this kind would explain the sign rather than merely establish it.

\subsection{First contact and transport between windows}\label{w8:fam:R25}

Three related ideas propagate positivity in the window parameter. If positivity fails, then by
\cite[Lemmas~3.10 and~10.17]{P46} the function $T\mapsto\lam$ is continuous and there is a first window $\Tst$
with $\lambda(\Tst)=0$, attained by a null vector of definite parity (see Section~\ref{w2:sec:structure}). One
may look for a contradiction in the structure at this first contact: the moments of the null vector, its
boundary jets, a quadrature at its nodes, a transmission condition at $\pm\Tst$, a conservation law forcing zero
flux. One may try to show that positivity in the window $T$ implies positivity in $T+dT$, or across the entrance
of a new contact $n$ at $T_n=\frac12\log n$. Or one may look for a map from tests of a larger window to tests of
smaller windows, or to a Hilbert space, with an exact remainder of manifest sign.

\phantomsection\label{w8:tab:contact}
We begin with the first contact. For $F$ supported in $\IT$ with $T<\Tst$ one has $\aT(F)=\mathfrak
a_{\Tst}(F)$, by the proof of Proposition~\ref{w2:prop:monotone}: the forms in smaller windows are evaluations
of the form at $\Tst$ on functions extended by zero. So no information descends from the contact to smaller
windows, and proving that the support of a null vector descends to a smaller window is, as far as we can see,
the window problem itself. Several detectors of a first contact were written, for instance the boundary flux of
the null vector between two contacts, measured by the norm of the difference of two evaluation operators. They
are exact and computable, but each of them is known to vanish at a first contact only if one already knows that
there is none. We also tried to read the sign from quantum-information inequalities applied to the null vector
or to the prime part. None of these attempts succeeded, each for a specific reason, which we state without
proof. For recovery maps of Petz type, the relevant fibres have one element, so recovery is trivial. For
covariance and uncertainty inequalities of canonical-anticommutation type, the connected correlations give an
orientation without a sign, and the full covariance gives a sign but loses the orientation. In the
Kadison--Schwarz inequality the equality case is reached and carries no information. Under complex polarization,
every positive square obtained contains the full quotient comb again. For an isometric channel built from the
explicit formula, the coherent channel has infinite rank, and the Tate moments remove only two modes. A
transversality lemma for supports of Euler--Poisson type fails at $N=4$. In a two-prime harmonic model of the
flux, harmonicity gives a constant, nonzero flux.

There is a common reason behind these failures. The structure at a hypothetical first contact is shared by
models that do have a first contact: finite off-line models, the quasicrystal deformations $f_t$ and the other
falsifiers of Table~\ref{w3:tab:falsifiers}. A contradiction drawn only from that structure would also apply to
them; what is needed is a property of the null vector that they do not share. The moments of the null vector
enter only through the eigenvalue equation, and the boundary identities are invariant under change of
orientation. Two objects have never been written down: a global boundary relation, with its Green identity at
the contact, that does not impose the vanishing of the flux, and a conservation law for the flux in an
indefinite metric.

Propagation in $T$ meets two obstacles. Extension by zero gives $\lambda(T')\le\lambda(T)$ for $T<T'$
(Proposition~\ref{w2:prop:monotone}). This is the only inclusion between windows, and it points the wrong way
for an induction: positivity in a large window implies positivity in smaller ones, never conversely. The second
obstacle is the entrance of a new contact, for which the following first-order heuristic applies.

\begin{observation}[the cost of a new contact]\label{w8:obs:jump}
Let $F$ be real, even, continuous up to the boundary of $\IT$, with $T$ slightly larger than $T_n=\frac12\log n$ and
$\delta=2T-\log n$ small. Then $\cF(\log n)=\int_{T-\delta}^{T}F(t)F(t-\log n)\,dt\approx\delta\,F(T)F(-T)=\delta F(T)^2$.
So the new term $-v_n\cF(\log n)$ in $\aT(F)$ has $T$-derivative $\approx-2v_nF(T)^2$ at the entrance of the contact.
Heuristically, a minimizer that does not vanish at the edge therefore loses form value at a rate proportional to
the square of its boundary value, and the loss has no compensating sign.
\end{observation}

The hypothesis is the weak point: a null vector at a first contact is not expected to be more regular than
$H^s$ for $s<\frac12$ (Section~\ref{w2:sec:structure}), so the boundary value need not exist, and the
observation is not a statement about $\lam$. Every propagation argument through the contact set that we examined
reduces to the Gr\"onwall inequality of Subsection~\ref{w8:fam:R24}.

We turn to transport. Convolution maps two windows to their sum, and gives an exact identity. To state it, write
$\mathfrak L$ for the functional on autocorrelations defined by the full comb,
$\mathfrak L(c)=\frac1{2\pi}\int m\,\widehat c-\sum_{n\ge2}v_n\Rel c(\log n)$. Then $\aT(F)=\mathfrak L(\cF)$ for
every $F\in\DT$ and every $T$, since the terms with $\log n\ge2T$ vanish.

\begin{proposition}[transport identity]\label{w8:prop:transport}
Let $F\in C_c^\infty(\IT)$ and let $G\in C_c^\infty(I_S)$ be real. For $s\in\R$ put
$\Delta_sF(t)=F(t+\tfrac s2)-F(t-\tfrac s2)$. Then $F*G\in C_c^\infty(I_{T+S})$, $\Delta_sF\in C_c^\infty(I_{T+|s|/2})$,
\begin{equation}\label{w8:eq:transport}
  \mathfrak a_{T+S}(F*G)=\widehat G(0)^2\,\aT(F)-\frac12\int_{|s|\le2S}c_G(s)\,\mathfrak a_{T+|s|/2}(\Delta_sF)\,ds,
\end{equation}
and $F*G$ and $\Delta_sF$ are primitive whenever $F$ is.
\end{proposition}

\begin{proof}
$c_{F*G}=\cF*c_G$, by Fubini. Expanding the product,
$c_{\Delta_sF}(u)=2\cF(u)-\cF(u-s)-\cF(u+s)$. Since $G$ is real, $c_G$ is even, and $\int c_G=\widehat G(0)^2$. Hence
$\frac12\int c_G(s)c_{\Delta_sF}(u)\,ds=\widehat G(0)^2\cF(u)-(\cF*c_G)(u)$, that is,
$c_{F*G}=\widehat G(0)^2\cF-\frac12\int c_G(s)c_{\Delta_sF}\,ds$, with $c_G(s)=0$ for $|s|\ge2S$. Apply the linear
functional $\mathfrak L$. The interchange with the $s$-integral is legitimate because all functions are smooth with
uniformly bounded supports. Then use $\mathfrak L(c_H)=\mathfrak a_{T'}(H)$ for $H$ supported in $I_{T'}$. For
primitivity, $\Phi_{F*G}=\PhiF\Phi_G$ and
$\Phi_{\Delta_sF}(z)=\PhiF(z)\bigl(e^{-(z-\frac12)s/2}-e^{(z-\frac12)s/2}\bigr)$, and both vanish at $z=0,1$ when
$\PhiF$ does.
\end{proof}

The identity uses only the linearity of $\mathfrak L$, so it holds for every weight sequence, in particular for
the rescalings and for Beurling weights, and it cannot carry a sign by itself (Theorem~\ref{w3:bar:B1}). Its
remainder has no manifest sign. For $G\ge0$ one has $c_G\ge0$, and the remainder is minus an average of the
forms $\mathfrak a_{T+|s|/2}(\Delta_sF)$ in windows up to $T+S$, whose sign is the unknown; for general $G$ the
weight $c_G$ may in addition change sign. Since these are the very windows one wants to reach, the step is
circular. A support count gives an obstruction to the simplest linear transport. For $H\in\mathcal P_{T'}$,
$T'>T$, the autocorrelation $c_H$ in general has support $[-2T',2T']$, while every element of the closed convex
cone generated by $\{\cF:F\in\PT\}$ is supported in $[-2T,2T]$. So $c_H$ is in general not in the closed cone
generated by autocorrelations of the smaller window. We read this as saying that a useful transport must be
nonlinear in the test function, with an exact remainder that is nonnegative and is not evaluated in the larger
window. No such map is known.

Finally, one may look for a map $U$ from tests to a Hilbert space with $\aT(F)=\|UF\|^2$ or
$\aT(F)\ge\|UF\|^2$, together with a uniform comparison to an energy norm. Such a comparison cannot hold.

\begin{proposition}[no uniform coercivity]\label{w8:prop:nocoerc}
There is no constant $c>0$ with $\aT(F)\ge c\|F\|_{\DT}^2$ for all $F\in\PT$ and all $T\ge1$.
\end{proposition}

\begin{proof}
$\|F\|_{\DT}\ge\|F\|_2$ up to a fixed constant, so such a $c$ would give $\lam\ge c'>0$ for all $T\ge1$. We show
$\inf_{T\ge1}\lam\le0$. If $\xi$ has a zero off the line, Theorem~\ref{w1:thm:criterion} gives a real primitive
$F_+\in C_c^\infty(\IT)$, for some $T$, with $\aT(F_+)<0$, so $\lam<0$ for that $T$ and, by
Proposition~\ref{w2:prop:monotone}, for all larger $T$. If all zeros lie on the line, then
\eqref{w1:eq:explicit} gives $\aT(F)=\sum_\gamma|\Fh(\gamma)|^2$ for smooth $F\in\PT$. Fix $\tau_0>0$ that is not an
ordinate, $\phi$ smooth, even, supported in $(-1,1)$, $\phi\ne0$, and put
$F_R=(-\partial_t^2+\frac14)\bigl[R^{-1/2}\phi(t/R)\cos(\tau_0t)\bigr]$, supported in $I_R$. Then
$\widehat{F_R}(\tau)=(\tau^2+\frac14)\,\widehat{G_R}(\tau)$ with $G_R$ the bracket, so $\widehat{F_R}(\pm\frac i2)=0$,
that is $A(F_R)=B(F_R)=0$, and $F_R$ is real and primitive. As $R\to\infty$, $\|F_R\|_2^2\to\frac12(\tau_0^2+\frac14)^2\|\phi\|_2^2>0$,
while $\widehat{G_R}(\gamma)$ is $R^{1/2}$ times a sum of two translates of $\widehat\phi(R\,\cdot)$ centred at $\pm\tau_0$;
since $\widehat\phi$ decays faster than any power, the ordinates have at most logarithmic density and none equals
$\pm\tau_0$, the sum $\sum_\gamma(\gamma^2+\frac14)^2|\widehat{G_R}(\gamma)|^2$ tends to $0$. Hence
$\lambda(T)\le\aT(F_R)/\|F_R\|_2^2\to0$ for $T>R\to\infty$. In both cases $\inf_{T\ge1}\lam\le0$.
\end{proof}

For maps without a uniform constant we know no general obstruction, only failed attempts, which we report
without proof. For an explicit map of Mellin--Gauss type, the domination $\aT(F)\ge\|UF\|^2$ fails at
$T=\frac45$ and frequency $\frac{5\pi}8$: a computation in interval arithmetic, not reproduced here, gives a form value
below $0.00236$ while $\|UF\|^2$ exceeds $0.2613$. The identity $\aT=\|U\cdot\|^2$ fails for the same map on
polar-corrected tests of high frequency. Every realization of the prime operators as compressions
$PU(\log n)J$ of a unitary group, with fixed bounded maps $P,J$, appears to need an infinite spectral moment. On
the other hand, we expect that positivity on a dense polar-free even image implies global positivity, so that a
domination on a dense core of polar-free tests would suffice; the map would have to use the exact values of
$\Lambda$ in its definition. Auxiliary positive structures were found (anticorrelation of Gaussian prime
variables, contractive transfer operators, a symmetry of the polar term), but they survive a finite erasure of
the weights and are therefore not specific to $\Lambda$ (Theorem~\ref{w3:bar:B1}).

For the finite quotients of the Connes--Consani--Moscovici construction, the exact transport between
consecutive quotients is known: it consists of two rank-one channels of opposite parity and a metric correction
\cite[Theorem~4.1 and the theorem `Complete transport defect' of \S5]{P39}. In that setting every von Mangoldt
atom is load-bearing at the sharp constant \cite[Theorem~24.1]{P39}, and the signed selection of the branch
remains open. Propagation arguments in one parameter are discussed further in \cite{P27,P29}. The paper
\cite{P30} concerns a different one-parameter family, the vertical shifts of $\zeta$, and gives a conditional
dichotomy under sequential self-approximation; it does not bear on the window problem and is mentioned only for
the analogy.

\section{What a proof of the sign must use}\label{w8:sec:constraints}

Each barrier theorem of Section~\ref{w3:sec:barriers} and each test weight of Section~\ref{w3:sec:falsifiers}
excludes one class of arguments. Read together they say a good deal about what a proof of the window problem
(Problem~\ref{w1:prob:window}) has to contain, and this section collects what they say. The requirements below are
not theorems about all possible proofs. Each of them has exactly the strength of the statement it rests on, which
is cited in each case; where the supporting statement is only a heuristic, we say so.

\subsection{The density of the primes is not enough}

One requirement can be proved completely here, because it concerns the archimedean symbol only. Split the prime
form into a main term and a fluctuation. For real $F\in\DT$, with $\psi(t)=\sum_{n\le t}\Lambda(n)$,
\[
\pT(F)=\int_{[1,e^{2T}]}2t^{-1/2}\cF(\log t)\,d\psi(t),
\]
and replacing $d\psi(t)$ by $\kappa\,dt$ gives the \emph{main term with density $\kappa$},
\begin{gather*}
\kappa\int_0^{2T}2e^{u/2}\cF(u)\,du=\kappa\int_\R e^{|u|/2}\cF(u)\,du
=\kappa\Bigl(2A(F)B(F)-\mathfrak e_T(F)\Bigr),\\
\mathfrak e_T(F)=\frac1{2\pi}\int_\R\frac{|\Fh(\tau)|^2}{\frac14+\tau^2}\,d\tau .
\end{gather*}
This uses $e^{|u|/2}=2\cosh(u/2)-e^{-|u|/2}$, $\int\cF(u)e^{\pm u/2}du=A(F)B(F)$ for real $F$, and
$\int e^{-|u|/2}e^{i\tau u}du=(\frac14+\tau^2)^{-1}$. On $\PT$ the term $2AB$ vanishes. So the form with the
main term of density $\kappa$ in place of the prime form is
$\frac1{2\pi}\int\bigl(m(\tau)+\kappa(\tfrac14+\tau^2)^{-1}\bigr)|\Fh|^2$.

\begin{proposition}\label{w8:prop:kappa}
Let $\kappa^*=\sup_{\tau\in\R}\bigl(-m(\tau)\bigr)(\tfrac14+\tau^2)$. If $\kappa\ge\kappa^*$, the form
$\mT+\kappa\,\mathfrak e_T$ is nonnegative on $\PT$ for every $T$. If $\kappa<\kappa^*$, it is negative somewhere on
$\PT$ for all large $T$. One has $\kappa^*\ge-C_0/4\approx1.3430$, the value at $\tau=0$.
\end{proposition}

\begin{proof}
The first claim is pointwise positivity of the symbol
$s_\kappa(\tau)=m(\tau)+\kappa(\frac14+\tau^2)^{-1}$. For the second, $s_\kappa$ is continuous and negative at some
$\tau_0$, hence negative on an interval around $\tau_0$; since $s_\kappa$ is even, we may take this interval inside
$(0,\infty)$, and we choose three distinct points $0<\tau_1<\tau_2<\tau_3$ in it. Let
$\phi$ be smooth, even, supported in $(-1,1)$, with $\|\phi\|_2=1$, and put
$F^{(j)}_R(t)=R^{-1/2}\phi(t/R)\cos(\tau_jt)$. As $R\to\infty$, $|\widehat{F^{(j)}_R}|^2$ concentrates at $\pm\tau_j$.
So the matrix of the form $\frac1{2\pi}\int s_\kappa\widehat F\overline{\widehat G}$ on these three functions tends to
$\frac12\mathrm{diag}(s_\kappa(\tau_j))$, which is negative definite, and their $L^2$ Gram matrix tends to
$\frac12 I$. For large $R$ the form is therefore negative definite on the three-dimensional span. The conditions
$A=B=0$ cut out a subspace of codimension at most two, which meets the span nontrivially. For $T>R$ this gives an
element of $\PT$ with negative form.
The lower bound $\kappa^*\ge-C_0/4$ is the value at $\tau=0$, which uses only
$(\Gamma'/\Gamma)(\frac14)=-\gamma-\frac\pi2-3\log2$.
\end{proof}

A numerical maximization of $\tau\mapsto(-m(\tau))(\frac14+\tau^2)$ gives $\kappa^*\approx7.4310$, attained at
$\tau_*\approx3.7769$; the function is negative beyond the zero $\tau\approx6.2898$ of $m$. We do not use these
values. What matters is that the main term of $\Lambda$ has density $\kappa=1<-C_0/4\le\kappa^*$, an inequality
that needs only the value at $\tau=0$. Hence positivity cannot come from the density of the primes.
Write
\begin{equation}\label{w8:eq:Phi}
\Phi_T(F)=-\int_{[1,e^{2T}]}2t^{-1/2}\,\cF(\log t)\,d\bigl(\psi(t)-t\bigr)
\end{equation}
for the fluctuation form. By the computation above with $\kappa=1$, $\pT(F)=2A(F)B(F)-\mathfrak e_T(F)-\Phi_T(F)$,
and therefore
\begin{equation}\label{w8:eq:aTsplit}
\aT(F)=\bigl(\mT+\mathfrak e_T\bigr)(F)+\Phi_T(F)-2A(F)B(F)\qquad(F\in\DT\text{ real}).
\end{equation}
On $\PT$ the fluctuation must satisfy $\Phi_T\ge-(\mT+\mathfrak e_T)$, a form whose symbol
$-s_1(\tau)=-m(\tau)-(\frac14+\tau^2)^{-1}$ is positive near the origin (for instance $-s_1(0)=-C_0-4\approx1.3722$).
The inequality $\Phi_T\ge(\kappa^*-1)\mathfrak e_T$ on $\PT$ would suffice and is not necessary.

\subsection{The requirements}

We list the requirements, grouped by what they demand. Three test weights recur: a multiple $c\Lambda$ of the
von Mangoldt function; a weight obtained from $\Lambda$ by changing finitely many values, for example by erasing
one prime power; and a Beurling system with the same initial data (Table~\ref{w3:tab:falsifiers}).

\begin{enumerate}[label=(\arabic*),leftmargin=*]
\item\label{w8:con:exact} \emph{Exactness at the scale of the margin.} An argument that proves positivity for
every weight in a neighbourhood of $\Lambda$ fails as soon as the neighbourhood contains $(1+\delta)\Lambda$ with
$\delta>\aT(F)/\pT(F)$ for some $F\in\PT$ with $\pT(F)>0$ (Proposition~\ref{w2:prop:concave}; see also
Proposition~\ref{w8:prop:scaling}). Since $\pT(F)$ stays of order one near the minimizer, this $\delta$ is,
heuristically, of the order of the margin, which decays at least exponentially in $T$
(Section~\ref{w2:sec:margin}). Main term and remainder of any estimate must therefore agree to that precision.
This does not mean that the sign has to be located with precision $\lam$ at each window: a lower bound
$\lam\ge-\varepsilon(T)$ with $\varepsilon(T)\to0$ suffices (Proposition~\ref{w2:prop:soft}).
\item\label{w8:con:controls}\label{w8:con:separate} \emph{A step that fails for every test weight.} A proof must
contain a step that is false for each of the three test weights, and more generally for every model that shares
the premises used: a finite model with zeros off the line, a quasicrystal deformation, an Epstein zeta function,
the Davenport--Heilbronn function, a truncated Euler product, a Beurling system
(Proposition~\ref{w3:prop:principle}, Theorem~\ref{w3:bar:B1}; see also Observation~\ref{w3:bar:B9}). An argument valid
for every weight, or affine in the weight with a remainder independent of it, proves nothing.
\item\label{w8:con:moments} \emph{Both primitivity conditions.} Without $A(F)=B(F)=0$ the form is not
nonnegative: every nonzero finite prime part is indefinite on $\PT$ (Theorem~\ref{w3:bar:B6}), and the second
condition, orthogonality of the autocorrelation to $u\sinh(u/2)$, cannot be dropped (Observation~\ref{w3:bar:B6b}).
\item\label{w8:con:couple}\label{w8:con:fluct} \emph{Archimedean and prime parts in one step, through the
fluctuation.} The archimedean form has at least $(2t_0/\pi)T-3\approx4T-3$ negative directions on $\PT$
(Theorem~\ref{w2:thm:archcount}), and beyond the threshold $\Tarch$ the primes must compensate them
(Remark~\ref{w2:obs:arch}). By Proposition~\ref{w8:prop:kappa} they cannot do so through their density: the
negative part of $\mT+\mathfrak e_T$ has to be paid by the fluctuation form $\Phi_T$. The two parts cannot be
bounded separately.
\item\label{w8:con:inputs} \emph{Euler product, functional equation and integrality together.} Each of the Euler
product over all primes, the functional equation with conductor one and gamma factor $\Gamma_\R$, and the
integrality of the multiplicities is necessary, as is shown by a test object that keeps the other two: truncated
Euler products, the Davenport--Heilbronn function, Beurling systems (Observation~\ref{w3:bar:B8} and
Table~\ref{w3:tab:falsifiers}). We expect, without proof, that the three must enter in one step: in a chain of
steps, each of these objects passes every link that does not use the input it lacks.
\item\label{w8:con:nonlin} \emph{Nonlinearity in the weight.} All the identities we know are affine in $\Lambda$,
and an identity valid for every weight carries no sign (Theorem~\ref{w3:bar:B1}). The exact relations coming from
$\Lambda*1=\log$ are linear in $\Lambda$. The only nonlinearity we have met comes from the minimizer of the form,
and using it is circular. A proof must contain a step that is nonlinear in $\Lambda$, or an input that enters as
an inequality.
\item\label{w8:con:nonlocal}\label{w8:con:finite} \emph{No finite or local data.} Positivity cannot be read zero
by zero, nor from finitely many zeros (Observation~\ref{w3:bar:B5}), nor prime by prime, since the form of each prime is
indefinite (Section~\ref{w4:sec:local}). No finite list of exact linear data of $\Lambda$ decides all windows
(Theorem~\ref{w3:bar:B11}), and finitely many moments, Euler values or verified windows do not determine the sign
(Theorem~\ref{w3:bar:B4}). Each fixed window, on the other hand, is a condition on finitely many values of
$\Lambda$ (Section~\ref{w2:sec:structure}); the difficulty is the uniformity in $T$.
\item\label{w8:con:criteria} \emph{No determinant, moment or bounded-range criterion.} Determinants carry
$\lam$ as a factor (Observation~\ref{w3:bar:B3}). Criteria on the low part of the symbol, and functionals of
$\{\lambda(T')\}$ on a bounded range of $T'$, are excluded (Theorem~\ref{w3:bar:B4}, Observation~\ref{w3:bar:B4b}).
Two-sided asymptotic laws for $\lam$ are not excluded (\S\ref{w8:fam:R24}).
\item\label{w8:con:import}\label{w8:con:sos} \emph{No positivity that holds for every weight.} A quantity that is
nonnegative for every positive weight cannot equal $\aT$ (Observation~\ref{w3:bar:B10}); an auxiliary positive object is
useful only if its exact identity with the form is proved. In the same direction, representations of $\aT$ on
$\PT$ as sums of squares of Dirichlet polynomials with arithmetic coefficients are not available for
$T\ge\frac12\log3$ (\S\ref{w4:fam:R1}); we do not know whether they exist for $\frac12\log2<T<\frac12\log3$.
\item\label{w8:con:cancel}\label{w8:con:sign} \emph{Cancellation must be kept, and the bound needed is an upper
bound.} Expanding the prime form term by term in a basis loses too much: absolute majorants between frequency
bands lose a square root of the number of bands. The compensation between the even and odd sectors, and between
integers with different numbers of prime factors, has to be preserved. Moreover the unconditional tools on the
prime side give lower bounds for the prime form where an upper bound is needed (Observation~\ref{w3:bar:B12}). These
statements are heuristic.
\item\label{w8:con:nearnull} \emph{A nearly null subspace of growing dimension.} Numerically there are eight
directions with Rayleigh quotient below $4\cdot10^{-6}$ at $T=1$, and more as $T$ grows; heuristically their
number is of order $e^{2T}$. An identity of the form $\aT(F)=\sum_jc_j|\ell_j(F)|^2$ must therefore have nearly
vanishing functionals on a subspace whose dimension grows with $T$.
\item\label{w8:con:uniform} \emph{One construction for all $T$.} The bottom $\lam$ is nonincreasing in $T$
(Proposition~\ref{w2:prop:monotone}), so there is no induction from small windows to large ones, and no rate can
be obtained from the total mass of the weights (Observation~\ref{w3:bar:B4b}). The object that carries the sign must
be defined for all windows at once.
\end{enumerate}

\subsection{The shape of an identity that would give the sign}\label{w8:fam:R26}

Instead of looking for a proof, one may ask what any proof by an identity would have to produce, and test
candidates against that description before trying to prove anything.

\begin{definition}\label{w8:def:signobject}
A \emph{sign object} is a family $(\mathcal H_T,V_T,\rho_T)_{T>0}$, where $\mathcal H_T$ is a Hilbert space,
$V_T:\PT\to\mathcal H_T$ a linear map, and $\rho_T$ a functional on $\PT$, such that
\begin{equation}\label{w8:eq:signobject}
\aT(F)=\|V_TF\|^2+\rho_T(F)\qquad(F\in\PT,\ T>0),
\end{equation}
and such that $\rho_T\ge0$ follows from a structure given before the identity is proved (a positive measure, a
positive operator, a positive kernel) and not from \eqref{w8:eq:signobject} itself.
\end{definition}

Without the last condition the definition is empty: if $\lam\ge0$, then the square root $V_T$ of the operator of
the closed form and $\rho_T=0$ satisfy \eqref{w8:eq:signobject}. Conversely, a sign object gives $\lam\ge0$
(cf.\ Proposition~\ref{w4:prop:taut}). The content lies in what the requirements above impose on the identity:
\begin{enumerate}[label=(\roman*)]
\item it is an exact equality for every $F$ and every $T$, with no $o(\cdot)$ terms, absolute values or losses by
the Cauchy--Schwarz inequality (requirement~\ref{w8:con:exact});
\item it fails for each of the three test weights and for each object of Table~\ref{w3:tab:falsifiers}, and one
can point to the place in the construction where it fails (requirement~\ref{w8:con:controls});
\item $V_T$ or $\rho_T$ depends nonlinearly on the weights, and $V_T$ is written without inverting or factorizing
the form, since otherwise the construction is a polarization of the form itself (requirement~\ref{w8:con:nonlin} and \S\ref{w4:fam:R2});
\item $\|V_TF\|^2$ does not vanish on test functions that see near-coincidences $\ell\approx kn$ of integers,
and $V_T$ nearly annihilates a subspace of growing dimension (requirement~\ref{w8:con:nearnull});
\item the construction is one for all $T$ (requirement~\ref{w8:con:uniform}).
\end{enumerate}

\begin{problem}\label{w8:prob:signobject}
Construct a sign object satisfying (i)--(v).
\end{problem}

Many candidates were tested against this description, among them algebraic, combinatorial, analytic and
physical constructions, and factorizations through a projection whose kernel contains the radical of the form.
None satisfied it: each fell into a class with a common counterexample weight, or turned out to be a change of
coordinates inside the radical of the form. We know of no construction that meets the description, and none is
excluded beyond the classes listed above.

\section{Conclusion}\label{w8:sec:open}

We end with the present state of the problem: how the sign was eventually obtained, what that proof does and does
not say about the approaches described above, and which of the questions met along the way remain open.

\subsection{How the sign was obtained}

None of the approaches of Sections~\ref{w4:sec:squares}--\ref{w8:sec:families} gave the sign of the Weil form on
all windows by working with the form itself. The sign is now a theorem, proved by a different road. In
\cite{P46} it is shown that no Dirichlet $L$-function vanishes in $\Rel s>\frac12$
\cite[Theorem~15.5]{P46}; in particular the Riemann hypothesis holds \cite[Theorem~16.1]{P46}, and by Weil's
criterion (Theorem~\ref{w1:thm:criterion}) the form $\aT$ is nonnegative on $\PT$ for every $T$
\cite[Corollary~16.3]{P46}. The proof starts from a one-sided criterion of Landau type
\cite[Theorem~12.9]{P46}, of the kind discussed in \S\ref{w7:fam:R18}: the absence of zeros in a half-plane
follows from an upper bound for a prime sum weighted by $1-\chi_K$, where $\chi_K$ is the character of an
imaginary quadratic field $K$. The bound itself does not come from analytic number theory. It is obtained from
three theorems on integral algebraic modular forms on a chain of theta lifts between unitary groups over a CM
field, the Short Forms Theorem, the Spanning Theorem and the Covolume Theorem, which give, for suitable arithmetic
data, the bound required by the criterion. The mathematics behind these three theorems, and the way from them to
the criterion, are due to Alpoge \cite{CharacterFamily}; the account in \cite{P46} and its companion papers follows that work.

It is instructive to compare this proof with the requirements of Section~\ref{w8:sec:constraints}. The criterion
is one-sided, so the inequality that is proved is an upper bound for a prime sum, as it must be
(requirement~\ref{w8:con:cancel}). The weight $1-\chi_K$ is nonnegative and vanishes on the primes that split in
$K$, so the sum is sensitive to the values of the weight on individual primes and not only to their density
(requirement~\ref{w8:con:couple}); and the argument uses the integrality of arithmetic objects in an essential
way, which is what separates the integers from the Beurling systems of \S\ref{w7:fam:R20}
(requirement~\ref{w8:con:inputs}). It also leaves the setting of all the attempts described here in one
respect: it works with Dirichlet $L$-functions of conductor greater than one and with the characters $\chi_K$ of
imaginary quadratic fields, whereas those attempts concern $\zeta$ alone, and the average of the Weil form over
a family of characters, which was tried as well, carries no sign. On the other hand the proof does not produce an identity of the form
\eqref{w8:eq:signobject}, nor a dual certificate, nor any of the positive objects looked for in
Sections~\ref{w4:sec:squares}--\ref{w6:sec:operators}. It gives the sign of $\aT$ through the zeros, by Weil's
criterion, and says nothing about the size of $\lam$. The questions about the form itself are therefore untouched
by it, and the barrier theorems of Section~\ref{w3:sec:barriers} keep their interest: they explain why the many
attempts to find the sign inside the explicit formula did not succeed, and they apply verbatim to other
$L$-functions for which no proof of this kind is available.

\subsection{Questions that remain}

Now that the sign is known, the exact reformulations of the window problem collected in the previous sections
are true statements. For each of them one may still ask for a proof in its own terms, and such a proof would be
of interest because it would be likely to extend to situations where the route of \cite{P46} does not apply. We
single out those we find most natural.

\emph{The even sector.} The form $\aT$ is invariant under $F(t)\mapsto F(-t)$, which exchanges $A$ and $B$ and
preserves $\PT$. Hence $\PT=\PT^{\mathrm{ev}}\oplus\PT^{\mathrm{odd}}$ is an orthogonal decomposition for both
$L^2$ and $\aT$, and $\lam=\min(\lambda^{\mathrm{ev}}(T),\lambda^{\mathrm{odd}}(T))$. Numerically the minimizer
is even (Section~\ref{w2:sec:margin}). On even real functions $A(F)=B(F)=\int F(t)\cosh(t/2)\,dt$, so an even
function is primitive if and only if $\ell_T(F):=\langle F,\cosh(\cdot/2)\mathbf 1_{\IT}\rangle=0$.

\begin{problem}[the polar--prime inequality]\label{w8:prob:polarprime}
Prove directly that for every $T>0$ and every real even $F\in\DT$ with $\ell_T(F)=0$,
\[
Q_T(F):=\bigl(\mT+\mathfrak e_T\bigr)(F)+\Phi_T(F)\ \ge\ 0 ,
\]
with $\mathfrak e_T$ and $\Phi_T$ as in \eqref{w8:eq:Phi}--\eqref{w8:eq:aTsplit}.
\end{problem}

\begin{proposition}\label{w8:prop:polarprime}
For each $T$, the inequality of Problem~\ref{w8:prob:polarprime} at $T$ is equivalent to
$\lambda^{\mathrm{ev}}(T)\ge0$. The window problem is equivalent to it together with
$\lambda^{\mathrm{odd}}(T)\ge0$ for all $T$. For fixed $T$, the set of weight vectors $(\Lambda'(n))_{n\in\NT}$
for which the inequality holds with $\psi$ built from $\Lambda'$ is closed and convex.
\end{proposition}

\begin{proof}
On real even $F$ with $\ell_T(F)=0$ we have $A(F)=B(F)=0$, so \eqref{w8:eq:aTsplit} gives $Q_T(F)=\aT(F)$. The form
is real on real functions, and a complex even primitive function has real and imaginary parts that are even and
primitive, so the bottom on $\PT^{\mathrm{ev}}$ is attained on real functions. The second assertion follows from the
sector decomposition. The third is Proposition~\ref{w1:prop:convex}: $Q_T$ is affine in the weights through
$d\psi$.
\end{proof}

In these coordinates the archimedean part and the main term of the primes are combined in the single form
$\mT+\mathfrak e_T$, and all the arithmetic sits in the fluctuation $\Phi_T$. On a basis of about $4T$ even
Dirichlet sines the compression of $\Phi_T$ is explicit, a Toeplitz part plus a Hankel part
(Proposition~\ref{w2:prop:combsine}), and by the min--max principle the inequality forces $\Phi_T$ to be positive
on at least as many directions orthogonal to $\cosh(t/2)$ as $\mT+\mathfrak e_T$ has negative ones there. A model
with the same archimedean factor and the same weights $(\Lambda(n))_{n\in\NT}$ has the same $Q_T$, so the test
objects that defeat most approaches are invisible in these coordinates unless they change the weights in the
window. This is the formulation in which a direct proof seems to us least hopeless.

\emph{Other reformulations.} The window problem is equivalent to the positivity of explicit finite Schur
complements (Section~\ref{w2:sec:reductions}); what is missing is a proof of that positivity by a structure that
does not inherit the loss of the margin through the reduction. It is also equivalent to the existence, for every
$T$ and every $\varepsilon>0$, of dual certificates of value at least $-\varepsilon$
(Section~\ref{w5:sec:dual}); here one would like an arithmetic description of the optimal certificate, and a
lower bound for the deficit of certificates sampled on a lattice in terms of the window and the sampling step.
On the zero side, the one-sided inequality for the tail of the Li coefficients after splitting them into
archimedean and prime parts (Section~\ref{w6:sec:li}) and the coercivity hypothesis of \cite{P08} for the
Kre\u{\i}n realization (Section~\ref{w6:sec:operators}) are in the same position: for neither do we have
a proof in its own terms.

\emph{Quantitative questions.} The proof in \cite{P46} gives $\lam\ge0$ and nothing more. Explicit positive
lower bounds are known only on bounded windows (Table~\ref{w2:tab:certified}). A two-sided law
$\lam=L(T)(1+o(1))$ with an explicit $L$ (\S\ref{w8:fam:R24}), and even the exact rate of decay of the margin
(Section~\ref{w2:sec:margin}), are open. So is the question whether $\aT$ is a sum of squares of Dirichlet
polynomials with arithmetic coefficients on $\PT$ for $\frac12\log2<T<\frac12\log3$
(requirement~\ref{w8:con:import}), and Problem~\ref{w8:prob:signobject} on the existence of an identity that
exhibits the sign.

\emph{Approaches that were never carried out.} A few directions were identified and not pursued, and we mention
them without recommending them, since some may fall into one of the excluded classes once written down:
the Gelfand--Naimark--Segal construction over the adele
class algebra for the state defined by the explicit formula; capacity-theoretic energies on the Berkovich
projective line; the use of pair-correlation and $n$-level statistics of the zeros
\cite{Montgomery1973,RudnickSarnak1996} as inputs to the sign; the Carath\'eodory--Fej\'er extension problem for
the Toeplitz part of the prime form on Dirichlet sines; extensions, in the sense of Kre\u{\i}n, of the symmetric
operator underlying the zero-side realization whose spectral measure carries atoms at $\log p$;
and a reflection-positive model of the complete form, in the sense of Osterwalder and Schrader
\cite{OsterwalderSchrader1973}, that is not a product over the primes, product models being excluded factor by
factor (\S\ref{w4:fam:R6}).

\medskip
The picture that comes out of this survey is consistent. The Weil form on a finite interval is positive by a
margin that is exponentially small, on a space where the archimedean part alone has a growing number of negative
directions; every approach that treats the primes through their density, through finitely many of them, or
through a property shared by nearby weights, is defeated by an explicit counterexample; and the approaches that
avoid these defects end in exact reformulations of the problem. The proof that exists goes around the form
rather than through it. Whether the sign can be seen in the explicit formula itself, by an identity or by an
inequality that uses the arithmetic of $\Lambda$ at the scale of the margin, is the question this survey leaves
to the reader.

\appendix
\section{The approaches at a glance}\label{w8:sec:table}

Table~\ref{w8:tab:master} lists the approaches discussed in Sections~\ref{w4:sec:squares}--\ref{w8:sec:families}
with, in a few words, the point at which each one stops and the place where it is discussed. It is an index to the
text and does not replace it.

{\small
\setlength{\tabcolsep}{4pt}
\begin{longtable}{>{\raggedright}p{4.1cm}>{\raggedright}p{7.6cm}>{\raggedright\arraybackslash}p{2.2cm}}
\caption{The approaches and where they stop.}\label{w8:tab:master}\\
\toprule
Approach & Where it stops & Discussed in\\\midrule\endfirsthead
\toprule
Approach & Where it stops & Discussed in\\\midrule\endhead
\bottomrule\endfoot
Sums of squares of Dirichlet polynomials & not available for $T\ge\frac12\log3$; open in between & \S\ref{w4:fam:R1}\\
A signed object plus one that sees the weight & the identity holds for every weight, and the remainder does not see $\Lambda$ & \S\ref{w4:fam:R2}\\
Decomposition by prime or by place & each local form is indefinite & \S\ref{w4:fam:R3}\\
Cells, lattices, near-coincidences & exact reformulation of the problem & \S\ref{w4:fam:R4}\\
Truncations and regularized energies & the self-energy is unbounded below & \S\ref{w4:fam:R5}\\
Positivity imported from probability, geometry or physics & holds for every positive weight, or reproduces $\zeta$ & \S\ref{w4:fam:R6}\\
Dual certificates and linear programming & works window by window, with margins that decay; in general an exact reformulation & \S\ref{w5:fam:R7}\\
Determinants, Wiener--Hopf, Schur complements & determinants carry $\lam$ as a factor; the Schur reduction inherits the margin & \S\ref{w5:fam:R8}\\
Positive kernels and total positivity & the kernel lies in no positive class that contains the models & \S\ref{w5:fam:R9}\\
Kre\u{\i}n spaces & the zero side does not control the angular operator & \S\ref{w6:fam:R10}\\
de~Branges spaces and Herglotz--Pick criteria & equivalent statements & \S\ref{w6:fam:R11}\\
Chains of exact identities & identities carry no sign & \S\ref{w6:fam:R12}\\
Li--Keiper coefficients & the archimedean part does not dominate term by term; the tail is of the strength of the problem & \S\ref{w6:fam:R13}\\
Jensen polynomials and the heat flow & each level holds unconditionally; the limit is the statement & \S\ref{w6:fam:R14}\\
Prolate operators, Sonin spaces, semilocal trace & no sign is produced & \S\ref{w6:fam:R15}\\
Statistical and localized detectors & data on the zero side do not decide & \S\ref{w6:fam:R16}\\
Classical estimates for prime sums & they bound sizes and are stable under a change of one coefficient & \S\ref{w7:fam:R17}\\
One-sided arithmetic criteria & equivalent statements; the bound needs an input of a different kind, which is how the sign was obtained in \cite{P46} & \S\ref{w7:fam:R18}\\
Euler product and functional equation & each is necessary, and chaining them fails & \S\ref{w7:fam:R19}\\
Beurling generalized integers & integrality is needed, and is not used by the methods available & \S\ref{w7:fam:R20}\\
Transfer of the Hodge index theorem & conditional on a count of sections that the frame does not have & \S\ref{w7:fam:R21}\\
Moments in an additive pencil & the moments carry no arithmetic & \S\ref{w7:fam:R22}\\
Global spaces and changes of category & not carried out & \S\ref{w7:fam:R23}\\
Deformation of the weight; two-sided laws & concavity transfers only between positive endpoints; the lower half of a law is the problem & \S\ref{w8:fam:R24}\\
First contact and transport between windows & monotonicity goes the wrong way; no uniform coercivity & \S\ref{w8:fam:R25}\\
An identity exhibiting the sign & no candidate satisfies the requirements & \S\ref{w8:fam:R26}\\
\end{longtable}
}

\begin{thebibliography}{99}
\bibitem{AdiprasitoHuhKatz2018} K.~Adiprasito, J.~Huh and E.~Katz, \emph{Hodge theory for combinatorial geometries}, Ann.\ of Math.\ (2) \textbf{188} (2018), 381--452.
\bibitem{BakryEmery1985} D.~Bakry and M.~\'Emery, \emph{Diffusions hypercontractives}, in: S\'eminaire de Probabilit\'es XIX, Lecture Notes in Math.\ \textbf{1123}, Springer, 1985, 177--206.
\bibitem{BerryKeating1999} M.~V.~Berry and J.~P.~Keating, \emph{The Riemann zeros and eigenvalue asymptotics}, SIAM Rev.\ \textbf{41} (1999), 236--266.
\bibitem{BeurlingDeny1959} A.~Beurling and J.~Deny, \emph{Dirichlet spaces}, Proc.\ Nat.\ Acad.\ Sci.\ U.S.A.\ \textbf{45} (1959), 208--215.
\bibitem{Bochner1932} S.~Bochner, \emph{Vorlesungen \"uber Fouriersche Integrale}, Akademische Verlagsgesellschaft, Leipzig, 1932.
\bibitem{Bombieri2000} E.~Bombieri, \emph{Remarks on Weil's quadratic functional in the theory of prime numbers, I}, Rend.\ Mat.\ Acc.\ Lincei (9) \textbf{11} (2000), 183--233.
\bibitem{BombieriHejhal1995} E.~Bombieri and D.~A.~Hejhal, \emph{On the distribution of zeros of linear combinations of Euler products}, Duke Math.\ J.\ \textbf{80} (1995), 821--862.
\bibitem{BombieriLagarias1999} E.~Bombieri and J.~C.~Lagarias, \emph{Complements to Li's criterion for the Riemann hypothesis}, J.\ Number Theory \textbf{77} (1999), 274--287.
\bibitem{BorceaBranden2009} J.~Borcea and P.~Br\"and\'en, \emph{The Lee--Yang and P\'olya--Schur programs.\ I.\ Linear operators preserving stability}, Invent.\ Math.\ \textbf{177} (2009), 541--569.
\bibitem{BostConnes1995} J.-B.~Bost and A.~Connes, \emph{Hecke algebras, type III factors and phase transitions with spontaneous symmetry breaking in number theory}, Selecta Math.\ (N.S.) \textbf{1} (1995), 411--457.
\bibitem{BottcherSilbermann2006} A.~B\"ottcher and B.~Silbermann, \emph{Analysis of Toeplitz Operators}, 2nd ed., Springer, 2006.
\bibitem{BrascampLieb1976} H.~J.~Brascamp and E.~H.~Lieb, \emph{On extensions of the Brunn--Minkowski and Pr\'ekopa--Leindler theorems, including inequalities for log concave functions, and with an application to the diffusion equation}, J.\ Funct.\ Anal.\ \textbf{22} (1976), 366--389.
\bibitem{BDR2023} F.~Broucke, G.~Debruyne and Sz.~Gy.~R\'ev\'esz, \emph{Some examples of well-behaved Beurling number systems}, Trans.\ Amer.\ Math.\ Soc.\ \textbf{378} (2025), 477--501.
\bibitem{CohnElkies2003} H.~Cohn and N.~Elkies, \emph{New upper bounds on sphere packings I}, Ann.\ of Math.\ (2) \textbf{157} (2003), 689--714.
\bibitem{Connes1999} A.~Connes, \emph{Trace formula in noncommutative geometry and the zeros of the Riemann zeta function}, Selecta Math.\ (N.S.) \textbf{5} (1999), 29--106.
\bibitem{ConnesConsani2021} A.~Connes and C.~Consani, \emph{Weil positivity and trace formula, the archimedean place}, Selecta Math.\ (N.S.) \textbf{27} (2021), Paper No.~77.
\bibitem{CCM2024} A.~Connes, C.~Consani and H.~Moscovici, \emph{Zeta zeros and prolate wave operators}, arXiv:2310.18423.
\bibitem{ConreyLi2000} J.~B.~Conrey and X.-J.~Li, \emph{A note on some positivity conditions related to zeta and $L$-functions}, Internat.\ Math.\ Res.\ Notices 2000, no.~18, 929--940.
\bibitem{DavenportHeilbronn1936} H.~Davenport and H.~Heilbronn, \emph{On the zeros of certain Dirichlet series}, J.\ London Math.\ Soc.\ \textbf{11} (1936), 181--185 and 307--312.
\bibitem{deBranges1968} L.~de~Branges, \emph{Hilbert Spaces of Entire Functions}, Prentice-Hall, 1968.
\bibitem{deBruijn1950} N.~G.~de~Bruijn, \emph{The roots of trigonometric integrals}, Duke Math.\ J.\ \textbf{17} (1950), 197--226.
\bibitem{Deligne1974} P.~Deligne, \emph{La conjecture de Weil. I}, Publ.\ Math.\ Inst.\ Hautes \'Etudes Sci.\ \textbf{43} (1974), 273--307.
\bibitem{Deninger1998} C.~Deninger, \emph{Some analogies between number theory and dynamical systems on foliated spaces}, Doc.\ Math.\ Extra Vol.\ ICM Berlin 1998, Vol.~I, 163--186.
\bibitem{DMV2006} H.~G.~Diamond, H.~L.~Montgomery and U.~M.~A.~Vorhauer, \emph{Beurling primes with large oscillation}, Math.\ Ann.\ \textbf{334} (2006), 1--36.
\bibitem{DiamondZhang2016} H.~G.~Diamond and W.-B.~Zhang, \emph{Beurling Generalized Numbers}, Math.\ Surveys Monogr.\ \textbf{213}, Amer.\ Math.\ Soc., 2016.
\bibitem{Douglas1966} R.~G.~Douglas, \emph{On majorization, factorization, and range inclusion of operators on Hilbert space}, Proc.\ Amer.\ Math.\ Soc.\ \textbf{17} (1966), 413--415.
\bibitem{ErdosKac1940} P.~Erd\H{o}s and M.~Kac, \emph{The Gaussian law of errors in the theory of additive number theoretic functions}, Amer.\ J.\ Math.\ \textbf{62} (1940), 738--742.
\bibitem{Ford2002} K.~Ford, \emph{Vinogradov's integral and bounds for the Riemann zeta function}, Proc.\ London Math.\ Soc.\ (3) \textbf{85} (2002), 565--633.
\bibitem{Garding1959} L.~G{\aa}rding, \emph{An inequality for hyperbolic polynomials}, J.\ Math.\ Mech.\ \textbf{8} (1959), 957--965.
\bibitem{GORZ2019} M.~Griffin, K.~Ono, L.~Rolen and D.~Zagier, \emph{Jensen polynomials for the Riemann zeta function and other sequences}, Proc.\ Natl.\ Acad.\ Sci.\ USA \textbf{116} (2019), 11103--11110.
\bibitem{Grothendieck1958} A.~Grothendieck, \emph{Sur une note de Mattuck--Tate}, J.\ Reine Angew.\ Math.\ \textbf{200} (1958), 208--215.
\bibitem{Grothendieck1969} A.~Grothendieck, \emph{Standard conjectures on algebraic cycles}, in: Algebraic Geometry (Internat.\ Colloq., Tata Inst.\ Fund.\ Res., Bombay, 1968), Oxford Univ.\ Press, 1969, 193--199.
\bibitem{Hamburger1921} H.~Hamburger, \emph{\"Uber die Riemannsche Funktionalgleichung der $\zeta$-Funktion}, Math.\ Z.\ \textbf{10} (1921), 240--254.
\bibitem{Hartshorne1977} R.~Hartshorne, \emph{Algebraic Geometry}, Grad.\ Texts in Math.\ \textbf{52}, Springer, 1977.
\bibitem{Haynsworth1968} E.~V.~Haynsworth, \emph{Determination of the inertia of a partitioned Hermitian matrix}, Linear Algebra Appl.\ \textbf{1} (1968), 73--81.
\bibitem{Hinkkanen1997} A.~Hinkkanen, \emph{On functions of bounded type}, Complex Variables Theory Appl.\ \textbf{34} (1997), 119--139.
\bibitem{Ingham1932} A.~E.~Ingham, \emph{The Distribution of Prime Numbers}, Cambridge Tracts \textbf{30}, Cambridge Univ.\ Press, 1932.
\bibitem{KaczorowskiPerelli1999} J.~Kaczorowski and A.~Perelli, \emph{On the structure of the Selberg class, I: $0\le d\le1$}, Acta Math.\ \textbf{182} (1999), 207--241.
\bibitem{KeatingSnaith2000} J.~P.~Keating and N.~C.~Snaith, \emph{Random matrix theory and $\zeta(1/2+it)$}, Comm.\ Math.\ Phys.\ \textbf{214} (2000), 57--89.
\bibitem{Keiper1992} J.~B.~Keiper, \emph{Power series expansions of Riemann's $\xi$ function}, Math.\ Comp.\ \textbf{58} (1992), 765--773.
\bibitem{Lagarias1999} J.~C.~Lagarias, \emph{On a positivity property of the Riemann $\xi$-function}, Acta Arith.\ \textbf{89} (1999), 217--234.
\bibitem{LandauWidom1980} H.~J.~Landau and H.~Widom, \emph{Eigenvalue distribution of time and frequency limiting}, J.\ Math.\ Anal.\ Appl.\ \textbf{77} (1980), 469--481.
\bibitem{Lang1966} S.~Lang, \emph{Introduction to Transcendental Numbers}, Addison-Wesley, 1966.
\bibitem{Laumon1987} G.~Laumon, \emph{Transformation de Fourier, constantes d'\'equations fonctionnelles et conjecture de Weil}, Publ.\ Math.\ Inst.\ Hautes \'Etudes Sci.\ \textbf{65} (1987), 131--210.
\bibitem{Li1997} X.-J.~Li, \emph{The positivity of a sequence of numbers and the Riemann hypothesis}, J.\ Number Theory \textbf{65} (1997), 325--333.
\bibitem{MSS2015} A.~W.~Marcus, D.~A.~Spielman and N.~Srivastava, \emph{Interlacing families I: Bipartite Ramanujan graphs of all degrees}, Ann.\ of Math.\ (2) \textbf{182} (2015), 307--325.
\bibitem{MatomakiRadziwill2016} K.~Matom\"aki and M.~Radziwi\l\l, \emph{Multiplicative functions in short intervals}, Ann.\ of Math.\ (2) \textbf{183} (2016), 1015--1056.
\bibitem{Meyer2005} R.~Meyer, \emph{On a representation of the idele class group related to primes and zeros of $L$-functions}, Duke Math.\ J.\ \textbf{127} (2005), 519--595.
\bibitem{Michalowski2026} W.~Micha{\l}owski, \emph{On the P\'olya frequency order of the de Bruijn--Newman kernel: certified failure at order five}, preprint, arXiv:2602.20313 (2026).
\bibitem{Montgomery1973} H.~L.~Montgomery, \emph{The pair correlation of zeros of the zeta function}, in: Analytic Number Theory, Proc.\ Sympos.\ Pure Math.\ \textbf{24}, Amer.\ Math.\ Soc., 1973, 181--193.
\bibitem{MontgomeryVaughan1973} H.~L.~Montgomery and R.~C.~Vaughan, \emph{The large sieve}, Mathematika \textbf{20} (1973), 119--134.
\bibitem{MontgomeryVaughan1974} H.~L.~Montgomery and R.~C.~Vaughan, \emph{Hilbert's inequality}, J.\ London Math.\ Soc.\ (2) \textbf{8} (1974), 73--82.
\bibitem{Mori15} R.~Mori, \emph{Weil positivity on the window of width $\log 15$}, Zenodo, doi:10.5281/zenodo.23281236 (2026).
\bibitem{Newman1976} C.~M.~Newman, \emph{Fourier transforms with only real zeros}, Proc.\ Amer.\ Math.\ Soc.\ \textbf{61} (1976), 245--251.
\bibitem{OsterwalderSchrader1973} K.~Osterwalder and R.~Schrader, \emph{Axioms for Euclidean Green's functions}, Comm.\ Math.\ Phys.\ \textbf{31} (1973), 83--112.
\bibitem{PlattTrudgian2021} D.~Platt and T.~Trudgian, \emph{The Riemann hypothesis is true up to $3\cdot10^{12}$}, Bull.\ London Math.\ Soc.\ \textbf{53} (2021), 792--797.
\bibitem{Polya1927} G.~P\'olya, \emph{\"Uber die algebraisch-funktionentheoretischen Untersuchungen von J.~L.~W.~V.~Jensen}, Kgl.\ Danske Vid.\ Sel.\ Math.-Fys.\ Medd.\ \textbf{7} (1927), 3--33.
\bibitem{Prekopa1973} A.~Pr\'ekopa, \emph{On logarithmic concave measures and functions}, Acta Sci.\ Math.\ (Szeged) \textbf{34} (1973), 335--343.
\bibitem{ReedSimon4} M.~Reed and B.~Simon, \emph{Methods of Modern Mathematical Physics IV: Analysis of Operators}, Academic Press, 1978.
\bibitem{RodgersTao2020} B.~Rodgers and T.~Tao, \emph{The de Bruijn--Newman constant is non-negative}, Forum Math.\ Pi \textbf{8} (2020), e6.
\bibitem{RudnickSarnak1996} Z.~Rudnick and P.~Sarnak, \emph{Zeros of principal $L$-functions and random matrix theory}, Duke Math.\ J.\ \textbf{81} (1996), 269--322.
\bibitem{Schoenberg1951} I.~J.~Schoenberg, \emph{On P\'olya frequency functions I.\ The totally positive functions and their Laplace transforms}, J.\ Analyse Math.\ \textbf{1} (1951), 331--374.
\bibitem{Selberg1946} A.~Selberg, \emph{Contributions to the theory of the Riemann zeta-function}, Arch.\ Math.\ Naturvid.\ \textbf{48} (1946), 89--155.
\bibitem{Selberg1992} A.~Selberg, \emph{Old and new conjectures and results about a class of Dirichlet series}, in: Proc.\ Amalfi Conf.\ Analytic Number Theory (Maiori, 1989), Univ.\ Salerno, 1992, 367--385.
\bibitem{SlepianPollak1961} D.~Slepian and H.~O.~Pollak, \emph{Prolate spheroidal wave functions, Fourier analysis and uncertainty, I}, Bell System Tech.\ J.\ \textbf{40} (1961), 43--63.
\bibitem{Smith1876} H.~J.~S.~Smith, \emph{On the value of a certain arithmetical determinant}, Proc.\ London Math.\ Soc.\ \textbf{7} (1875/76), 208--212.
\bibitem{Titchmarsh1986} E.~C.~Titchmarsh, \emph{The Theory of the Riemann Zeta-Function}, 2nd ed., revised by D.~R.~Heath-Brown, Oxford Univ.\ Press, 1986.
\bibitem{Trudgian2014} T.~S.~Trudgian, \emph{An improved upper bound for the argument of the Riemann zeta-function on the critical line II}, J.\ Number Theory \textbf{134} (2014), 280--292.
\bibitem{Vaaler1985} J.~D.~Vaaler, \emph{Some extremal functions in Fourier analysis}, Bull.\ Amer.\ Math.\ Soc.\ (N.S.) \textbf{12} (1985), 183--216.
\bibitem{Vaughan1977} R.~C.~Vaughan, \emph{Sommes trigonom\'etriques sur les nombres premiers}, C.~R.\ Acad.\ Sci.\ Paris S\'er.~A \textbf{285} (1977), 981--983.
\bibitem{Viazovska2017} M.~S.~Viazovska, \emph{The sphere packing problem in dimension 8}, Ann.\ of Math.\ (2) \textbf{185} (2017), 991--1015.
\bibitem{Weil48} A.~Weil, \emph{Sur les courbes alg\'ebriques et les vari\'et\'es qui s'en d\'eduisent}, Hermann, Paris, 1948.
\bibitem{Weil52} A.~Weil, \emph{Sur les ``formules explicites'' de la th\'eorie des nombres premiers}, Comm.\ S\'em.\ Math.\ Univ.\ Lund (1952), tome suppl\'ementaire, 252--265.
\bibitem{Yoshida1992} H.~Yoshida, \emph{On Hermitian forms attached to zeta functions}, in: Zeta Functions in Geometry (Tokyo, 1990), Adv.\ Stud.\ Pure Math.\ \textbf{21}, Kinokuniya, Tokyo, 1992, 281--325.
\bibitem{Zhang2007} W.-B.~Zhang, \emph{Beurling primes with RH and Beurling primes with large oscillation}, Math.\ Ann.\ \textbf{337} (2007), 671--704.
\bibitem{RStation45} \emph{Weil positivity on the window of width $\log 11$}, computational certificate with supplement,
\url{https://rstation.ai/problems/riemann_hypothesis/papers/45} (2026).

\bibitem{Zhu2026} X.~Zhu, \emph{Weil positivity in compact windows: a finite reduction, certified two-sided bounds, and a Landau--Widom decay law}, preprint, arXiv:2608.24827 (2026).
\bibitem{P01} D.~A.~Trejo Pizzo, \emph{The Critical Line: A Survey of Recent Progress on the Riemann Hypothesis}, Paper 01 of the series, preprint.
\bibitem{P02} D.~A.~Trejo Pizzo, \emph{The $\omega$-class decomposition of Dirichlet polynomials: elementary bounds and computations}, Paper 02 of the series, preprint.
\bibitem{P03} D.~A.~Trejo Pizzo, \emph{On the Sign of Inter-Class Interference in Dirichlet Partial Sums: A Peak--Trough Reconciliation}, Paper 03 of the series, preprint.
\bibitem{P04} D.~A.~Trejo Pizzo, \emph{Partial sums of the Davenport--Heilbronn function at the ordinates of its off-line zeros}, Paper 04 of the series, preprint.
\bibitem{P05} D.~A.~Trejo Pizzo, \emph{The $\omega$-hierarchy of the Riemann zeta function: per-prime weights, moment exponents and the branching picture of the maximum}, Paper 05 of the series, preprint.
\bibitem{P06} D.~A.~Trejo Pizzo, \emph{Positivity formulations of the Riemann hypothesis and the function-field analogy: a heuristic checklist and some remarks}, Paper 06 of the series, preprint.
\bibitem{P07} D.~A.~Trejo Pizzo, \emph{The localized Weil form on a Hermite--Gauss basis: the effect of an off-line zero pair and the truncation of the prime sum}, Paper 07 of the series, preprint.
\bibitem{P08} D.~A.~Trejo Pizzo, \emph{A Kre\u{\i}n-space formulation of the zero side of the Weil quadratic form on Paley--Wiener spaces}, Paper 08 of the series, preprint.
\bibitem{P09} D.~A.~Trejo Pizzo, \emph{Stable Local Arithmetic Fingerprints of Localized Explicit-Formula Forms}, Paper 09 of the series, preprint.
\bibitem{P10} D.~A.~Trejo Pizzo, \emph{The localized-Weil anatomy of $L$-functions: Euler products, Satake power sums, and Rankin--Selberg products}, Paper 10 of the series, preprint.
\bibitem{P11} D.~A.~Trejo Pizzo, \emph{The off-line zeros as a Pontryagin index: a conditional reduction and numerical observations on the prime tail}, Paper 11 of the series, preprint.
\bibitem{P12} D.~A.~Trejo Pizzo, \emph{The Band-Limited Weil--Carleson Form: a numerically saturated generalized eigenvalue, sine-kernel structure of its spectrum, and a comparison of five reformulations}, Paper 12 of the series, preprint.
\bibitem{P13} D.~A.~Trejo Pizzo, \emph{Monotonicity of the Off-Line Mass under the de Bruijn--Newman Zero Flow}, Paper 13 of the series, preprint.
\bibitem{P14} D.~A.~Trejo Pizzo, \emph{Hodge-index and hyperbolicity reformulations of the Riemann hypothesis as positivity statements}, Paper 14 of the series, preprint.
\bibitem{P15} D.~A.~Trejo Pizzo, \emph{Eventual versus uniform positivity in scale-dependent hierarchies}, Paper 15 of the series, preprint.
\bibitem{P16} D.~A.~Trejo Pizzo, \emph{A second-moment identity for zero exponential sums and the localized Weil form}, Paper 16 of the series, preprint.
\bibitem{P17} D.~A.~Trejo Pizzo, \emph{An $\mathfrak{sl}_2$ triple on the archimedean continuum and arithmetic progressions among the zeros of $\zeta$: remarks on a hard-Lefschetz analogy}, Paper 17 of the series, preprint.
\bibitem{P18} D.~A.~Trejo Pizzo, \emph{The localized Weil form on even test functions: off-line Lorentzian blocks and obstructions to positive-kernel realizations}, Paper 18 of the series, preprint.
\bibitem{P19} D.~A.~Trejo Pizzo, \emph{The arithmetic defect of off-line zeros: the modified explicit formula and two-scale structure}, Paper 19 of the series, preprint.
\bibitem{P20} D.~A.~Trejo Pizzo, \emph{A Herglotz positivity criterion for the Riemann hypothesis}, Paper 20 of the series, preprint.
\bibitem{P21} D.~A.~Trejo Pizzo, \emph{Logarithmic sums of $q^{\omega(n)}$ and their shifted correlations: remarks}, Paper 21 of the series, preprint.
\bibitem{P22} D.~A.~Trejo Pizzo, \emph{The Riemann Hypothesis as moment rigidity over Meixner--Pollaczek polynomials}, Paper 22 of the series, preprint.
\bibitem{P23} D.~A.~Trejo Pizzo, \emph{Jacobi operators, a prime-sum potential and the logarithmic derivative of zeta on the critical line}, Paper 23 of the series, preprint.
\bibitem{P25} D.~A.~Trejo Pizzo, \emph{The Weil form and the scaling operator: a Pontryagin-index dictionary in zero coordinates}, Paper 25 of the series, preprint.
\bibitem{P26} D.~A.~Trejo Pizzo, \emph{The Weil--Toeplitz algebra of window forms and a logistic tensor-tower identity}, Paper 26 of the series, preprint.
\bibitem{P27} D.~A.~Trejo Pizzo, \emph{One-dimensional propagation and the Riemann Hypothesis: a heuristic survey of obstructions and two open directions}, Paper 27 of the series, preprint.
\bibitem{P28} D.~A.~Trejo Pizzo, \emph{Tensor amplification of the Weil negative index and a two-variable explicit formula on thickened diagonals}, Paper 28 of the series, preprint.
\bibitem{P29} D.~A.~Trejo Pizzo, \emph{One wall in many coordinates: a heuristic synthesis of the propagation trichotomy and a master quantifier for the Riemann Hypothesis}, Paper 29 of the series, preprint.
\bibitem{P30} D.~A.~Trejo Pizzo, \emph{Finiteness of off-line zeros and sequential self-approximation of zeta: a conditional reduction of the Riemann Hypothesis}, Paper 30 of the series, preprint.
\bibitem{P33} D.~A.~Trejo Pizzo, \emph{An Energy Decomposition of the Riemann Hypothesis: the Off-Line Energy $I$, its Integral Representation, and Structural Barriers}, Paper 33 of the series, preprint.
\bibitem{P34} D.~A.~Trejo Pizzo, \emph{Obstructions in a series of approaches to the Riemann Hypothesis: a heuristic survey}, Paper 34 of the series, preprint.
\bibitem{P35} D.~A.~Trejo Pizzo, \emph{The colligation route to the Riemann Hypothesis: a Kre\u{\i}n--Langer index, a Laplace--pole argument, and heuristic reformulations of Weil positivity}, Paper 35 of the series, preprint.
\bibitem{P36} D.~A.~Trejo Pizzo, \emph{The Riemann Hypothesis: An Obstruction Ledger and an Arithmetic Pick Reformulation}, Paper 36 of the series, preprint.
\bibitem{P37} D.~A.~Trejo Pizzo, \emph{The Limits of the Positivity Family in the Li--Keiper Route to the Riemann Hypothesis}, Paper 37 of the series, preprint.
\bibitem{P38} D.~A.~Trejo Pizzo, \emph{The Li Sequence as a Sampled Control System: FIR Window Stability, Off-Line Modes, and the Arithmetic Closure Problem}, Paper 38 of the series, preprint.
\bibitem{P39} D.~A.~Trejo Pizzo, \emph{Cofinal Weil Quotients: Exact Transport, Fixed-Cutoff Summability, and the Global Branch Obstruction}, Paper 39 of the series, preprint.
\bibitem{P40} D.~A.~Trejo Pizzo, \emph{Branch Selection and Global Polarization in the Connes--CCM Route to RH: Finite Prime--Gamma Inertia, Tate Boundary Geometry, and the Divisor-Sensitivity Obstruction}, Paper 40 of the series, preprint.
\bibitem{P41} D.~A.~Trejo Pizzo, \emph{Obstructions to Weil Positivity: A Catalogue of No-Go Theorems and Observations in the Prime--Gamma Doob Coordinate}, Paper 41 of the series, preprint.
\bibitem{P42} D.~A.~Trejo Pizzo, \emph{An Arithmetic Lefschetz Programme over $\Z$: A Metrized Bivariant Correspondence Square, its Determinant Pairing and its Corner Trace}, Paper 42 of the series, preprint.
\bibitem{P43} D.~A.~Trejo Pizzo, \emph{A Correspondence Algebra and the Weil Distribution on Semilocal Idele Class Groups}, Paper 43 of the series, preprint.
\bibitem{P44} D.~A.~Trejo Pizzo, \emph{An Arithmetic Frobenius Family on Semilocal Idele Class Groups}, Paper 44 of the series, preprint.
\bibitem{P45} D.~A.~Trejo Pizzo, \emph{An Explicit Formula for the Trace on Semilocal Idele Class Groups}, Paper 45 of the series, preprint.
\bibitem{P46} D.~A.~Trejo Pizzo, \emph{The Sign of the Radial Intersection Form on the Idele Class Group}, Paper 46 of the series, preprint.
\bibitem{CharacterFamily} L.~Alpoge, Lean~4 formalisation of the character-family argument
for Dirichlet $L$-functions, source archive in fourteen parts, November 2026,
\url{https://alpo.ge/lean.tar.gz.part01} (first part).
\bibitem{Liu2026} V.~Liu, \emph{Certified Weil positivity beyond the unit window: source-exact
block-Schur and tail-compensation bounds for the Riemann zeta function}, manuscript dated 14 Sep 2026,
alphaXiv preprint.
\bibitem{MoriUpper} R.~Mori, \emph{Improved upper bounds for windowed Weil forms with exponential factor
$e^{-4\pi\mu}$}, Zenodo (2026), v3, doi:\nolinkurl{10.5281/zenodo.23363767}.
\bibitem{P47} D.~A.~Trejo Pizzo, \emph{An Intersection Frame for the Explicit Formula: Hodge-Index Transfer, a Frame Defect Theorem and a Function-Field Model}, Paper 47 of the series, preprint.
\end{thebibliography}
\end{document}

